% Modernised reading — fonds Alexandre Grothendieck, shelfmark 29, pages 1-216.
%
% This is an INTERPRETATION, not a transcription. It re-reads the pages in
% today's notation and vocabulary, states the hypotheses the manuscript leaves
% implicit, and drops the critical apparatus so that the mathematics reads
% straight through. Everything the brackets and underlines carried — a doubtful
% reading, a notation he used differently, a missing passage — moves into a
% footnote.
%
% It is held to being mathematically correct as it stands, not merely faithful
% to the page. Where the manuscript is loose or mistaken, the true statement is
% what appears here, and the footnote says what the page has. In any dispute of
% substance the transcription governs: whoever cites Grothendieck must cite the
% batch-NN.fr.tex files.
%
% Scope: ONE document for the whole shelfmark — the convention every other
% folder of this project follows, and the one the site's manifest is built for.
% The eleven batches of transcription are read here as a single argument,
% because that is what they are: the question the 1967 typescript poses at its
% page 24 is answered at page 190, and the pair (C, s) it defines at page 12
% comes back at page 201 as a presentation of the object it was an
% approximation to. Split across eleven files, those two facts had to be
% asserted eleven times and could be read nowhere.
%
% NOTATION fixed once, for two hundred and sixteen pages:
%   Et(S) the étale site, Fl_S the fibred category of relative schemes;
%   CALLIGRAPHIC R = (C, s) for the 1967 *donnée de ramification*, the
%     auxiliary catalogue;
%   ROMAN (R, r) for the late *domaine de ramification*, the champ itself.
%     These are DIFFERENT objects that the carton gives nearly the same name,
%     and confusing them makes nonsense of both ends of the folder;
%   Rev^R(S) always with the exponent — the typescript writes both;
%   pi^1(S,xi;G) for the pointed-torsor functor, which the 1959 letter writes
%     Z(S,a;G) and the feuilles anciennes write with a gothic Z.
%
% Substantive departures from the pages, each footnoted where it occurs:
%   - page 161's nodal counter-example gives the FULL profinite completion;
%   - page 70 writes Pic(X) where its own argument requires Pic(Y);
%   - page 20's « n' >= n » is divisibility, as pages 2 and 37 settle;
%   - page 137's minimality has the two groups of the contracted product
%     exchanged;
%   - Lemme 6.5 of page 103 needs F and G invertible;
%   - page 30's Galois covering needs the commutation condition of page 36;
%   - page 68 credits purity with what normality gives;
%   - page 154's residue-field tower is drawn against its own inclusions;
%   - page 213's dimension bound is not sharp as written.
%
% Two statements are announced in the folder and never made — the Proposition 5
% of page 196 and the essential surjectivity of alpha at page 206 — and this
% reading records both as absences rather than stepping over them.
%
% Pass: Opus 5 (claude-opus-5), 2026-08-29 — first pass, unchecked against the
% pages by a human.
\documentclass[11pt,a4paper]{article}
% Le préambule commun à toutes les transcriptions du fonds Grothendieck.
%
% Il tient en une idée : **l'incertitude se compose, elle ne se résout pas.**
% Une transcription qui lisse un mot douteux a détruit la seule chose pour
% laquelle on la faisait — un lecteur doit pouvoir savoir, à chaque endroit,
% ce qui était sur le feuillet et ce qui a été ajouté. Les six macros qui
% suivent sont donc l'appareil critique, et non de la décoration.
%
% Les mêmes commandes sont reconnues par `scripts/render.mjs`, qui produit la
% vue de lecture du site. Une macro ajoutée ici doit l'être aussi là-bas,
% faute de quoi le rendu échouera bruyamment — ce qui est voulu : un rendu
% silencieusement faux est pire qu'un rendu absent.

\ProvidesPackage{grothendieck}[2026/08/08 Transcriptions du fonds Grothendieck]

\usepackage{amsmath,amssymb,amsthm}
\usepackage{mathrsfs}
\usepackage{stmaryrd}
% Les diagrammes commutatifs sont partout dans ce fonds : c'est la langue
% même de la géométrie algébrique de Grothendieck, et une transcription qui
% les rendrait en prose ne transcrirait rien.
\usepackage{tikz-cd}
% Français par défaut — c'est la langue des feuillets — mais l'anglais est
% chargé aussi : la traduction et le résumé sont des documents anglais, et la
% typographie française (espaces avant les deux-points) y serait une faute.
% Ces documents basculent par \selectlanguage{english} en tête de corps.
\usepackage[english,french]{babel}
\usepackage{fontspec}
\usepackage{geometry}
\geometry{a4paper,margin=25mm,marginparwidth=28mm}
\usepackage{marginnote}
\usepackage{xcolor}
% ulem, not soul. `soul`'s \st and \ul are built on a character-by-character
% scan that XeLaTeX's font machinery breaks: the document fails with an
% undefined control sequence rather than degrading. ulem does the same two jobs
% and is Unicode-engine safe. `normalem` stops it hijacking \emph, which the
% transcriptions use for Grothendieck's own emphasis.
\usepackage[normalem]{ulem}
\usepackage{hyperref}
\hypersetup{colorlinks=true,linkcolor=black,urlcolor=black,pdfborder={0 0 0}}

% --- Métadonnées du lot -----------------------------------------------------
% Reprises telles quelles par le script de rendu, qui les affiche en tête de
% la vue de lecture. Elles fixent ce que le fichier prétend couvrir : sans
% elles, une transcription flotte sans point d'attache dans un fonds de seize
% mille feuillets.

\newcommand{\folder}[1]{\gdef\grothFolder{#1}}
\newcommand{\batch}[1]{\gdef\grothBatch{#1}}
\newcommand{\leaves}[2]{\gdef\grothFirst{#1}\gdef\grothLast{#2}}
\newcommand{\foldertitle}[1]{\gdef\grothFoldertitle{#1}}
\gdef\grothFolder{?}\gdef\grothBatch{?}\gdef\grothFirst{?}\gdef\grothLast{?}\gdef\grothFoldertitle{}

% --- Le résumé --------------------------------------------------------------
%
% Il ouvre la lecture modernisée, sous le titre « Résumé », et il est composé
% en petit. Ce n'est pas un ornement : c'est le seul passage du document qui
% ne soit pas des mathématiques, et un lecteur doit pouvoir le reconnaître
% avant de l'avoir lu — puis le sauter s'il connaît déjà le sujet, ou s'y
% arrêter s'il ne le connaît pas. Le filet qui le suit marque la reprise du
% corps.

\newenvironment{resume}
  {\small\setlength{\parskip}{0.45em}}
  {\par\medskip\hrule\medskip}

% --- Mots-clés ---------------------------------------------------------------
%
% En anglais, en clôture du résumé de la lecture modernisée : le vocabulaire
% moderne sous lequel chercher ce que les feuillets construisent. Le manifeste
% du site (`npm run manifest`) les extrait pour étiqueter la cote entière —
% cette ligne est la source unique des tags, il n'y a pas de fichier à côté.
\newcommand{\keywords}[1]{\par\smallskip\noindent
  {\footnotesize\textsc{Keywords} --- \textit{#1}}\par}

% --- L'appareil critique ----------------------------------------------------

% Le numéro de feuillet, en marge. C'est l'ancre qui permet de régler un
% désaccord sur une formule en regardant *un* feuillet plutôt qu'une chemise
% de six cents. Le site s'en sert aussi pour tourner les pages du fac-similé
% quand on fait défiler la transcription.
\newcommand{\leaf}[1]{%
  \marginnote{\footnotesize\color{gray!70}\textsf{#1}}%
  \label{leaf:#1}\ignorespaces}

% Illisible : on ne devine pas. Un mot inventé qui se lit comme les autres est
% le pire résultat possible.
\newcommand{\ill}{\textcolor{red!60!black}{[\,\ldots\,]}}

% Lecture douteuse : le mot est donné, et signalé comme incertain.
\newcommand{\uncertain}[1]{\uline{#1}}

% Ajout de l'éditeur — une lettre manquante, un mot sous-entendu.
\newcommand{\add}[1]{\textcolor{blue!50!black}{[#1]}}

% Biffé par Grothendieck lui-même. Ce qu'il a rayé fait partie du document :
% une rature dit souvent où la pensée a changé de direction.
\newcommand{\struck}[1]{\sout{#1}}

% Note du transcripteur, sur l'état matériel du feuillet.
\newcommand{\note}[1]{\par\smallskip{\small\color{gray!60!black}\textit{[#1]}}\par\smallskip}

% Note marginale de Grothendieck, distincte du corps du texte.
\newcommand{\marginal}[1]{\par\smallskip\noindent\hspace{1em}%
  {\small\color{gray!60!black}#1}\par\smallskip}

% --- Titre ------------------------------------------------------------------

\AtBeginDocument{%
  \begin{center}
    {\large\bfseries Fonds Alexandre Grothendieck --- cote n\textsuperscript{o}~\grothFolder}\\[2pt]
    {\itshape\grothFoldertitle}\\[4pt]
    {\small feuillets \grothFirst--\grothLast\ (lot \grothBatch)}\\[2pt]
    {\footnotesize\color{gray!60!black}Transcription établie d'après le fac-similé numérisé
     par l'Université de Montpellier.\\
     Les lectures incertaines sont soulignées, les illisibles notées [\,\ldots\,],
     les ajouts entre crochets.}
  \end{center}
  \vspace{1em}}

\endinput


\folder{29}
\pages{1}{216}
\dating{1959-1969}
\watermark{Édition de démonstration\\interprétation personnelle de l'œuvre}
\foldertitle{Groupe fondamental [Autour de SGA 4 et SGA 7] : lettres (1959, 1967, 1969), tapuscrits et copies de tapuscrit annotés (s.d.), notes manuscrites (s.d.) — lecture modernisée du dossier entier}

\begin{document}

\section*{Résumé}

\begin{resume}
Deux cent seize feuillets, de 1959 à 1969, tournent autour d'une seule
question : \emph{qu'est-ce qu'un revêtement, en géométrie algébrique, quand on
autorise la ramification} ?

Le mot vient de la topologie. Au-dessus d'un espace on en dispose un autre,
qui le recouvre en plusieurs exemplaires sans jamais se déchirer ni se coller :
un escalier en colimaçon au-dessus d'une cour circulaire, qui monte d'un étage
à chaque tour. Toute l'information est dans la façon dont les étages se
permutent quand on fait le tour, et le groupe fondamental n'est que le registre
de ces permutations.

La difficulté commence quand on rebouche le puits. L'escalier, prolongé, retombe
sur un simple plancher, et le nombre de tours n'est plus lisible sur le
résultat. C'est exactement la situation d'un revêtement ramifié : deux données
très différentes ont la même \emph{réalisation géométrique}, et cette
réalisation ne les détermine pas. Un revêtement modérément ramifié n'est donc
pas un revêtement qui satisfait une \emph{propriété} : c'est un revêtement muni
d'une \emph{structure}. Les onze premières pages du carton sont le
procès-verbal de cette découverte, faite à deux. Jean-Pierre Murre, qui
rédigeait pour SGA 1 l'exposé sur le groupe fondamental modéré, écrit qu'il ne
voit pas d'où la structure pourrait venir ; Grothendieck lui répond que ses
propres notes étaient fautives, énumère cinq énoncés faux, et propose de tout
reprendre.

Ce que la reprise apporte est le geste que le carton entier répète :
\emph{cesser d'étudier l'objet, étudier la catégorie des objets}. Un revêtement
isolé ne dit pas de quel type de ramification il provient ; mais la collection
de tous les revêtements d'un type donné, avec ses sommes, ses produits et ses
quotients, en dit assez pour reconstituer un groupe. C'est le principe de la
théorie de Galois --- un corps ne se connaît pas tout seul, il se connaît par
ses extensions --- transposé aux schémas, et poussé jusqu'au bout. Car
l'échafaudage que 1967 doit se donner pour définir cette collection, la
« donnée de ramification », est arbitraire : deux catalogues très différents
définissent la même catégorie. Grothendieck le reconnaît sur-le-champ et donne
l'analogie juste, qui est la sienne : entre le catalogue et la catégorie, il y
a le rapport d'un \emph{site} à son \emph{topos}. Puis il écrit la phrase que
le carton mettra dix ans à honorer --- si l'on veut éliminer l'arbitraire,
« il semble qu'on soit \emph{infailliblement conduit} à prendre comme catégorie
la catégorie des revêtements elle-même qu'on se propose de construire ».

Les dernières pages font exactement cela. Un \emph{domaine de ramification}
n'est plus un catalogue : c'est un \emph{champ} au-dessus du site étale, une
famille de catégories qui se recollent, dont chaque fibre est une catégorie
multigaloisienne, munie d'une réalisation géométrique. Quatre axiomes suffisent.
Et les trois exemples de 1967 reviennent alors, dans l'ordre, comme trois
conséquences ; le couple qui servait de donnée revient comme une
\emph{présentation} de l'objet ; et les exemples kummériens, avec leur
obstruction dans $H^2$, sont ce qu'on appelle aujourd'hui un \emph{champ des
racines}. La boucle se referme dans un seul carton.

Entre les deux, le dossier fait autre chose encore, et il faut le signaler
parce qu'on ne s'y attend pas : une lettre à Serre du 18 octobre 1959, au
milieu du carton, définit le groupe fondamental des revêtements
\emph{infinitésimaux}. En caractéristique $p$ il existe des objets qui
ressemblent à des revêtements de degré $p$ et n'ont qu'\emph{un seul point} ;
le groupe fondamental ordinaire, qui ne sait compter que des points, ne les
voit pas. La lettre construit celui qui les voit, le calcule sur une variété
abélienne, et note sa dualité de Cartier. C'est ce qu'on appelle aujourd'hui le
schéma en groupes fondamental.

Le reste est la théorie à l'épreuve : des théorèmes de finitude ; le groupe
fondamental et le groupe de Picard \emph{birationnels}, obtenus l'un et l'autre
en coupant par un espace linéaire générique jusqu'à tomber sur une courbe ; la
monodromie d'une famille qui dégénère, dont la classe s'avère être le nombre
d'auto-intersection de la fibre centrale --- une identité que le carton énonce
quatre fois, à cent dix pages d'intervalle ; le théorème de monodromie locale,
en dix lignes ; et l'axiomatique des catégories galoisiennes, écrite ici trois
fois, dans trois états dont le dernier est celui de SGA 1.

Ce n'est pas un texte : c'est un dossier de travail, et l'ordre du carton n'est
pas l'ordre de la pensée. 1959 est à la page 124, 1967 au début, la synthèse à
la fin et sans date. On y voit un homme qui recule devant une vérification
(« jetant provisoirement un voile pudique sur ces abîmes »), qui juge sa propre
rédaction (« pénible et finalement peu utile »), qui écrit deux fois noir sur
blanc qu'il ne sait pas, qui annonce une proposition et ne l'énonce pas --- et
qui, en 1959, donne le calendrier du \emph{multiplodoque} : « Sauf difficultés
imprévues ou enlisement, il devrait être fini d'ici 3 ans, ou 4 au maximum. On
pourra commencer à faire de la géométrie algébrique ! »

\keywords{tame fundamental group, ramification datum, fibred category, Kummer covering, cartesian section, fibre functor, tamely ramified covering, normal crossings divisor, Abhyankar's lemma, faithfully flat descent, finiteness of the fundamental group, nonabelian H1, constructible sheaf, Serre quotient category, Grothendieck existence theorem, birational fundamental group, Bertini theorem, Picard scheme, Frobenius, purity theorem, specialization exact sequence, log smooth degeneration, self-intersection, Chern class of a divisor, Raynaud extension, local monodromy theorem, quasi-unipotence, inertia group, Lefschetz pencil, vanishing cycles, fundamental group scheme, Nori fundamental group, infinitesimal group scheme, Cartier duality, pro-representable functor, Galois category, axiomatic Galois theory, finite etale morphism, quotient by a finite group, Kunneth formula for the fundamental group, pro-etale fundamental group, enlarged fundamental group, strict pro-representability, effective descent datum, etale homotopy, ramification domain, gerbe, root stack, 2-category, stack of topoi, multigaloisian category, Galois topos, fundamental groupoid, generating family, degree of a morphism}
\end{resume}

\section*{Le fil du dossier, et l'ordre des feuillets}

Deux cent seize pages ne se lisent pas d'affilée, et celles-ci moins que
d'autres : le carton mêle quatre strates séparées par une décennie. Voici donc
le fil, et les conventions qui valent pour tout ce qui suit.

\subsection*{Les stations}

\begin{itemize}
  \item \textbf{Pages 1 à 40} --- la correspondance de 1967 avec Murre, et le
    tapuscrit « Données de ramification sur un schéma » qu'elle accompagne :
    la donnée $\mathcal{R} = (\mathcal{C}, s)$, la catégorie
    $\mathrm{Rev}^{\mathcal{R}}(S)$, le foncteur de réalisation géométrique
    $r$, le groupe $\pi_1^{\mathcal{R}}(S,s)$, et le manifeste de la page 31.
  \item \textbf{Pages 39 à 76} --- ce à quoi la machine sert : les conjectures
    de finitude, le théorème de finitude relatif, puis le groupe fondamental et
    le groupe de Picard birationnels.
  \item \textbf{Pages 77 à 100} --- la théorie appliquée : l'échange de 1969
    avec Murre sur les revêtements formels, et le calcul du $\pi_1$ modéré d'un
    schéma formel le long de sa fibre spéciale.
  \item \textbf{Pages 101 à 122} --- la boîte à outils : quasi-unipotence,
    Bertini et ses variantes, fonctorialité des groupes d'inertie.
  \item \textbf{Pages 124 à 160} --- la strate ancienne, 1959 et 1960 : la
    lettre à Serre sur le groupe fondamental infinitésimal, les « feuilles
    anciennes », puis les axiomes des catégories galoisiennes.
  \item \textbf{Pages 161 à 187} --- les revêtements infinis : le progroupe
    fondamental, les cinq classes de recouvrements, la descente et son
    formalisme simplicial.
  \item \textbf{Pages 188 à 216} --- la synthèse, sans date : les domaines de
    ramification, et l'axiomatique des catégories multigaloisiennes.
\end{itemize}

\subsection*{Deux objets, presque un seul nom}

C'est le seul piège de vocabulaire du carton, et il court d'un bout à l'autre.

La \emph{donnée} de ramification de 1967, notée ici
$\mathcal{R} = (\mathcal{C}, s)$, est un objet \emph{auxiliaire} : une
catégorie fibrée quelconque au-dessus du site étale, munie d'un foncteur vers
les schémas relatifs, dont on \emph{déduit} la catégorie des revêtements. Le
\emph{domaine} de ramification des pages 190 et suivantes, noté $(R, r)$,
\emph{est} cette catégorie, prise comme donnée, et c'est un champ. Les deux ne
sont pas de même nature, et le second est ce que la page 24 annonçait qu'on
serait infailliblement conduit à prendre. La casse les distingue partout dans
ce qui suit.

\subsection*{Le reste de la notation}

$\mathrm{Et}(S)$ est le site étale de $S$ et $\mathrm{Fl}_S$ la catégorie
fibrée des schémas relatifs au-dessus de lui ; le tapuscrit les souligne, ce
qui ne portait aucune information que la casse ne porte. On écrit
$\mathrm{Rev}^{\mathcal{R}}(S)$ toujours en exposant, là où le tapuscrit écrit
les deux. Et l'on écrit $\pi^1(S, \xi\,;\, G)$ pour le foncteur des fibrés
principaux pointés --- que la lettre de 1959 note $Z(S,a;G)$, que les feuilles
anciennes notent d'un $\mathfrak{Z}$ gothique, et que les pages 164 et
suivantes notent $\pi^1$ : trois notations, un seul foncteur, et c'est la
lecture du carton entier qui permet de le dire.

\subsection*{Ce qui n'est pas établi}

Deux énoncés sont annoncés dans le carton et n'y figurent pas. La page 196
écrit « Prop.\ 5 \dots » et rien après. La page 206 écrit « Prouvons que
$\alpha$ est essentiellement surjectif » et la démonstration ne suit pas, de
sorte que la reconstitution d'un domaine de ramification à partir d'une famille
génératrice reste \emph{conditionnelle} : la pleine fidélité est démontrée,
l'essentielle surjectivité ne l'est pas. Ces deux absences sont signalées à
leur place et ne sont pas comblées.

Enfin, une part importante du carton est illisible à la définition du
fac-similé --- les transcriptions portent mille cinq cent soixante-dix marques
d'illisibilité --- et les pages 61, 141, 155 et 164 à 169 sont des palimpsestes
où les formules se lisent et les phrases qui les joignent non. Là où c'est le
cas, cette lecture donne les énoncés et dit que les liaisons manquent, plutôt
que de les inventer.

\pagerange{1}{11}
\section*{Ce que Murre n'arrivait pas à voir (pages 1 à 11)}

La chemise porte, de la main de Grothendieck, deux mots : « Tame
ramification ». Dessous, quatre lettres de 1967, que le carton range dans le
désordre. Remises dans l'ordre où elles ont été écrites, elles racontent une
histoire complète.

\subsection*{29 mars 1967 : Murre envoie le manuscrit}

Murre rédige, pour le séminaire de 1960--61, l'exposé sur le groupe fondamental
modéré. Il envoie son manuscrit et sa propre liste de doutes : il a supposé
plus tôt que nécessaire l'existence des quotients ; il hésite sur le mot
\emph{covering}, qui traduit \emph{revêtement} mais entre en collision avec le
« recouvrement » d'une topologie ; il n'a pas trouvé de démonstration plus
simple pour deux de ses lemmes ; et il juge insatisfaisant, sans savoir
l'amender, l'énoncé de son théorème 9.4.

\subsection*{La réponse : cinq énoncés faux, et un diagnostic}

Grothendieck répond que les erreurs sont les siennes --- ses notes étaient trop
sommaires --- et il en donne la liste. La première est un contre-exemple, et
c'est le seul endroit de la correspondance où il calcule.

Soit $S$ un schéma, $b$ une section de $\mathcal{O}_{S}$, $n \geqslant 1$
inversible sur $S$, et posons $a = b^{n}$. Soit
\[
  Y' \;=\; \operatorname{Spec} \mathcal{O}_{S}[T]/(T^{n} - a) ,
\]
revêtement kummérien élémentaire, principal sous $H' = \mu_{n}$. Prenons pour
$H$ le groupe unité et $Y = S$.\footnote{Le tapuscrit écrit $Y\,T/(T^{n}-a)$ :
la machine ne portait pas de crochets, et $Y[T]/(T^{n}-a)$ est ce qu'il faut
lire. La transcription le signale ; c'est le premier des blancs mécaniques de
ce carbone, dont les flèches et les lettres grecques sont dans le même cas.}
Un morphisme de $Y$ dans $Y'$ compatible avec l'inclusion $H \to H'$ est
simplement une section de $Y'$ au-dessus de $S$ : elle existe, $T \mapsto b$.
Or $H \to H'$ n'est pas surjectif dès que $n > 1$. L'énoncé 1.16 des notes de
Murre --- qu'un tel morphisme force $H \to H'$ à être surjectif --- est donc
faux, et avec lui 3.6, puis les démonstrations de 3.7 et 3.8 qui s'en
servaient.

Le théorème 6.4 (le théorème d'Abhyankar) n'est correct, sous la forme donnée,
que si les diviseurs $D_{i}$ sont réguliers ; Grothendieck avertit au passage
que \emph{les groupes d'inertie ne sont déterminés qu'à automorphisme intérieur
près}, ce qui est la raison de la panne, et propose de remplacer la famille
$(D_{i})$ par un unique diviseur $D$ à croisements normaux, en demandant que
les revêtements soient modérément ramifiés localement pour la topologie étale
le long des composantes irréductibles locales de $D$. C'est la notion qu'il
appellera page 21 « modérément ramifié \emph{au sens large} », et c'est elle,
dit-il, qui convient au paragraphe 9.

La démonstration de 7.1 achoppe, elle, sur un point de principe : même pour
$D = 0$ elle aurait à faire intervenir des revêtements étales connexes
\emph{non galoisiens}. D'où le diagnostic, qui est l'origine de tout le reste :
\emph{il faut une notion de ramification modérée pour les revêtements non
galoisiens}, et les hypothèses de normalité et de réduction dont les notes se
servent sont artificielles.

\subsection*{29 avril 1967 : le plan en six paragraphes}

La lettre suivante accompagne l'esquisse que reproduisent les pages 12 à 20, et
propose de refondre les cinq premiers paragraphes de l'exposé en six :
revêtements kummériens relatifs à une famille de sections ; données de
ramification ; propriétés de la catégorie des revêtements de type
$\mathcal{R}$ et le groupe $\pi_{1}^{\mathcal{R}}(S,\;)$ ; fonctorialités ;
revêtements galoisiens de type $\mathcal{R}$ ; théorème d'Abhyankar pour les
diviseurs à croisements normaux.

Trois observations de cette lettre valent d'être retenues, parce que les pages
suivantes s'en servent sans y revenir.

D'abord, une remarque de régularité : si $A$ est un anneau local régulier,
$(x_{i})_{1 \leqslant i \leqslant r}$ une partie d'un système régulier de
paramètres et $(n_{i})$ des entiers $\geqslant 1$ \emph{non nécessairement
premiers à la caractéristique résiduelle}, alors
\[
  B \;=\; A[T_{1},\dots,T_{r}]\,/\,(T_{1}^{\,n_{1}} - x_{1},\ \dots,\
  T_{r}^{\,n_{r}} - x_{r})
\]
est encore régulier. C'est ce qui permet de ne pas regarder les
caractéristiques résiduelles des points non maximaux des $D_{i}$ : il en
résulte \emph{a posteriori} que si, en tout point maximal $s$ de chaque
$D_{i}$, le groupe d'inertie est d'ordre premier à la caractéristique de
$k(s)$, alors il l'est automatiquement en toute spécialisation de $s$ ---
énoncé vide sauf si $s$ est de caractéristique nulle.

Ensuite, une remarque de méthode : les démonstrations des paragraphes 7 et 8
« marchent pour n'importe quel type de données de ramification », pourvu qu'on
dispose d'un énoncé du type de son lemme 7.2. C'est la première indication que
la généralité coûte moins qu'elle ne rapporte, et c'est l'argument qu'il
opposera à Murre quand celui-ci proposera de se restreindre au cas kummérien.

Enfin, la lecture correcte des groupes de 10.2 et 10.3 : ce sont
$\pi_{1}(U,\;)$ et son plus grand quotient premier à $p$, et les théorèmes
devraient être énoncés en ces termes.

\subsection*{16 mai 1967 : Murre ne voit toujours pas}

La dernière lettre --- placée en tête du carton, écrite en dernier --- est la
plus utile pour lire l'esquisse, parce qu'elle pose deux questions dont les
réponses fixent des définitions que l'esquisse laisse incomplètes.

Murre demande si son intuition est juste : \emph{au lieu de donner directement
le revêtement $r(X)$ de $S$, on donne le comportement du revêtement au-dessus
d'autres revêtements, les $s(M)$, qui sont d'un type particulier} ; le
comportement au-dessus de tous les $s(M)$ ensemble en dit plus que $r(X)$
seul. C'est exactement cela, et c'est la meilleure phrase du dossier pour dire ce qu'est une donnée de ramification.

Il demande ensuite --- et c'est une lacune réelle --- si, dans la définition
d'un revêtement galoisien de type $\mathcal{R}$, on ne doit pas exiger que
l'action du groupe de Galois $G$ \emph{commute} à celle du groupe
d'automorphismes que les modèles $s(M)$ portent déjà. Grothendieck le confirme,
mais dans la lettre d'accompagnement, à la page 36 : la définition de la page
30 est incomplète telle qu'elle est écrite, et cette lecture-ci l'a complétée
là où elle se trouve.\footnote{La condition manquante est énoncée page 36 :
« the action of $G$ commutes to étale localization on $S$ \emph{and} to the
action of the groups of automorphisms given on the various $Z$'s ». Elle ne
figure pas au paragraphe 4 des pages 30 à 31, qui définit la notion. Elle est rétablie à sa place plus bas, avec le paragraphe 4 (pages 30 et 31).}

Il relève enfin, dans le cas kummérien, que la condition
$\underline{n}' \geqslant \underline{n}$ devrait s'écrire
$\underline{n}' = \underline{m} \cdot \underline{n}$. Grothendieck répond page
37 : « when I wrote $n' \geqslant n$, I meant of course the order relation of
divisibility ». La lecture ci-dessous écrit donc la divisibilité.

\pagerange{12}{14}
\section*{La donnée de ramification (pages 12 à 14)}

Soit $S$ un schéma et $\mathrm{Et}(S)$ son site étale. Une \emph{donnée de
ramification} sur $S$ est un couple
\[
  \mathcal{R} \;=\; (\mathcal{C},\, s)
\]
formé d'une catégorie fibrée $p : \mathcal{C} \to \mathrm{Et}(S)$ et d'un
foncteur cartésien
\[
  s : \mathcal{C} \longrightarrow \mathrm{Fl}_{S} ,
\]
où $\mathrm{Fl}_{S}$ est la catégorie fibrée sur $\mathrm{Et}(S)$ dont la fibre
en $S'$ est la catégorie des schémas relatifs sur $S'$.\footnote{Le tapuscrit
souligne les catégories : $\underline{C}$, $\underline{s}$, $\underline{Et}_S$,
$\underline{Fl}_S$, $\underline{R}$. Le soulignement est ici abandonné, comme
partout dans cette édition ; il ne portait aucune information que la casse ne
porte.} On ne suppose pas que $\mathcal{C}$ soit un \emph{champ} : ni les
objets ni les morphismes ne se recollent nécessairement. On ne suppose pas non
plus $s$ fidèle, encore moins pleinement fidèle --- et c'est le point de toute
l'affaire.

Un objet $M$ de la fibre $\mathcal{C}_{S'}$ s'appelle un \emph{modèle de
ramification}, et le $S'$-schéma $s(M)$ sa \emph{réalisation géométrique}. Dans
tous les exemples qui suivent, les $s(M)$ sont finis sur $S'$.

Composant $s$ avec le foncteur source, on obtient un foncteur
$\mathcal{C} \to (\mathrm{Sch})$, et l'on peut tirer en arrière le long de lui
la catégorie fibrée des schémas finis étales : on obtient une catégorie fibrée
$\mathcal{D}$ sur $\mathcal{C}$, dont la fibre en $M$ est la catégorie des
revêtements étales finis de $s(M)$.

\subsection*{Les revêtements de type $\mathcal{R}$}

Un \emph{crible couvrant} de $\mathcal{C}$ est une sous-catégorie pleine
$\mathcal{C}_{0}$ qui est un crible --- si l'extrémité d'une flèche est dans
$\mathcal{C}_{0}$, son origine y est --- et telle que les $p(M)$, pour $M$
parcourant $\mathcal{C}_{0}$, recouvrent $S$. Pour un tel crible, on pose
\[
  \mathrm{Rev}^{\mathcal{R}}_{\mathcal{C}_{0}}(S) \;=\;
  \text{catégorie des sections cartésiennes de } \mathcal{D}
  \text{ au-dessus de } \mathcal{C}_{0} ,
\]
et l'on définit la \emph{catégorie des revêtements de $S$ de type
$\mathcal{R}$} comme la limite inductive de ces catégories suivant les
restrictions :
\[
  \mathrm{Rev}^{\mathcal{R}}(S) \;=\;
  \varinjlim_{\mathcal{C}_{0}} \mathrm{Rev}^{\mathcal{R}}_{\mathcal{C}_{0}}(S) .
\]
Concrètement, un revêtement de type $\mathcal{R}$ est la donnée, pour chaque
modèle $M$ d'un crible couvrant, d'un revêtement étale fini $X(M)$ de $s(M)$,
et, pour chaque flèche $M' \to M$, d'un isomorphisme
$X(M') \simeq X(M) \times_{s(M)} s(M')$, ces isomorphismes se composant.
C'est la traduction exacte de la phrase de Murre : on ne donne pas le
revêtement, on donne son comportement au-dessus de tous les
modèles.\footnote{Le tapuscrit écrit tantôt
$\mathrm{Rev}^{\underline{R}}$, tantôt $\mathrm{Rev}_{\underline{R}}$, pour le
même objet, et parfois dans la même phrase (pages 13 et 14). L'exposant est
retenu ici pour tout le carton.}

Deux hypothèses accompagnent la définition, et il importe de les voir comme des
hypothèses.\footnote{Le tapuscrit écrit « Signalons que dans tous les cas que
nous aurons à utiliser, les foncteurs de transition \dots\ sont pleinement
fidèles », puis « Nous supposerons que \dots ». Ce sont bien deux conditions
posées sur les exemples, non deux propositions : la première est fausse pour
une donnée quelconque, et l'identification des
$\mathrm{Rev}^{\mathcal{R}}_{\mathcal{C}_{0}}(S)$ à des sous-catégories pleines
en dépend.}

\begin{enumerate}
  \item[(H1)] Les foncteurs de restriction
    $\mathrm{Rev}^{\mathcal{R}}_{\mathcal{C}_{0}}(S) \to
    \mathrm{Rev}^{\mathcal{R}}_{\mathcal{C}'_{0}}(S)$ sont pleinement fidèles.
    Alors les $\mathrm{Rev}^{\mathcal{R}}_{\mathcal{C}_{0}}(S)$ s'identifient,
    à équivalence près, à des sous-catégories pleines de
    $\mathrm{Rev}^{\mathcal{R}}(S)$, qui les épuise.
  \item[(H2)] Pour tout crible $\mathcal{C}_{0}$ et tout objet $X$ de
    $\mathrm{Rev}^{\mathcal{R}}_{\mathcal{C}_{0}}(S)$, le foncteur composé
    $\mathcal{C}_{0} \xrightarrow{\ X\ } \mathcal{D} \to (\mathrm{Sch})_{/S}$
    admet une limite inductive dans $(\mathrm{Sch})_{/S}$, notée $r(X)$, et
    cette limite ne change pas quand on restreint $X$ à un crible plus petit.
\end{enumerate}

Sous (H2) les $r$ partiels se raccordent en un unique foncteur
\[
  r : \mathrm{Rev}^{\mathcal{R}}(S) \longrightarrow (\mathrm{Sch})_{/S} ,
\]
la \emph{réalisation géométrique}. Les deux hypothèses sont vérifiées dans les
trois exemples ci-dessous, « de généralité croissante ».

Dans les cas qui intéressent, $r$ sera \emph{fidèle} sans être pleinement
fidèle. C'est là toute la mise : un revêtement de type $\mathcal{R}$ n'est pas
déterminé par sa réalisation géométrique, et il faut donc distinguer
soigneusement les deux. Quand $r$ est fidèle, on peut décrire
$\mathrm{Rev}^{\mathcal{R}}(S)$ comme la catégorie des « revêtements de $S$
munis d'une structure de type $\mathcal{R}$ », une telle structure sur un
revêtement $X$ étant la donnée d'une classe d'isomorphie d'objets
$X' \in \mathrm{Rev}^{\mathcal{R}}(S)$ et d'un isomorphisme
$r(X') \simeq X$.\footnote{La remarque est page 15, avec l'aveu qu'elle
« aurait dû figurer avant l'exemple ». Elle est remontée ici, où elle éclaire
la définition qu'elle commente.}

\pagerange{14}{16}
\section*{Exemple 1 : un revêtement à opérateurs (pages 14 à 16)}

C'est l'exemple à garder en tête, et Grothendieck le dira lui-même page 34 :
« the basic example in order to understand the idea is example 1 ».

Donnons-nous un $S$-schéma $Z$ muni d'une action à droite d'un groupe fini $G$
par $S$-automorphismes. Soit $\mathcal{C}_{0}$ la catégorie à un objet dont le
monoïde des endomorphismes est $G$, et prenons
$\mathcal{C} = \mathrm{Et}(S) \times \mathcal{C}_{0}$. Se donner un foncteur
cartésien de $\mathcal{C}$ dans une catégorie fibrée sur $\mathrm{Et}(S)$
revient à se donner un objet de la fibre au-dessus de $S$ muni d'une action de
$G$ ; le couple $(Z, G)$ fournit donc bien un $s$.

On trouve aussitôt une équivalence entre $\mathrm{Rev}^{\mathcal{R}}(S)$ et la
catégorie des revêtements étales finis $Z' \to Z$ munis d'une action à droite
de $G$ au-dessus de celle de $G$ sur $Z$ ; et la réalisation géométrique
existe si et seulement si le quotient $Z'/G$
existe, auquel cas $r(Z') = Z'/G$. Il suffit pour cela que toute orbite de $G$
dans $Z$ soit contenue dans un ouvert affine --- en particulier, il suffit que
$Z$ soit affine sur $S$, ce qui sera le cas puisque $Z$ sera un revêtement de
$S$.

\subsection*{Le sens intuitif}

En première approximation, les objets de $\mathrm{Rev}^{\mathcal{R}}(S)$ sont
les revêtements de $S$ qui \emph{deviennent étales} au-dessus du revêtement
ramifié donné $Z$. C'est ce qui donne son sens au mot « donnée de
ramification » : $Z$ est le patron de ramification que l'on s'autorise, et l'on
regarde tout ce qui n'est pas pire que lui.

Grothendieck le dit une seconde fois, page 34, dans le langage des corps de
fonctions, et cette formulation-là est la plus parlante. Supposons $S$ intègre
normal de corps de fonctions $K$, et $Z$ le normalisé de $S$ dans une extension
galoisienne finie $L$ de $K$. Un revêtement $S'$ de $S$, connexe, correspond à
une extension $K'$ de $K$. Dire que $S'$ n'a pas de ramification pire que
$Z/S$, c'est dire que l'image inverse normalisée de $S'$ au-dessus de $Z$ est
étale ; et lorsque $S$ est strictement local, cela revient tout simplement à
dire que \emph{$K'$ est isomorphe à une sous-extension de $L$}. La notion est
de nature locale pour la topologie étale, de sorte qu'on peut toujours se
ramener au cas strictement local, où elle est aussi simple que cela.

Alors, par transport de structure, $G$ opère sur $S'$, et l'on retrouve $S'$ à
partir du $G$-revêtement $Z'$ de $Z$ comme $S' = Z'/G$.

\subsection*{Quand $r$ est fidèle, quand il est pleinement fidèle}

Soit $U$ un ouvert de $S$ tel que $V = Z|U$ soit un revêtement principal de
groupe $G$, et supposons que toute partie ouverte et fermée de $Z$ stable par
$G$ et contenant $V$ soit égale à $Z$ --- c'est le cas si $V$ est dense dans
$Z$.

\begin{enumerate}
  \item[(a)] Alors $r$ est \emph{fidèle}. En effet le foncteur de restriction
    $Z' \mapsto Z'|V$, des $G$-revêtements étales de $Z$ dans ceux de $V$, est
    fidèle ; et le foncteur $Z'_{V} \mapsto Z'_{V}/G$, des $G$-revêtements
    étales de $V$ dans les revêtements étales de $U$, est une
    \emph{équivalence} par descente le long du revêtement principal
    $V \to U$ --- \emph{a fortiori} fidèle.
  \item[(b)] Si de plus $V$ est dense dans $Z$ et $Z$ est \emph{normal}, alors
    $Z' \mapsto Z'|V$ est pleinement fidèle, donc $r$ l'est aussi. Dans ce cas
    $\mathrm{Rev}^{\mathcal{R}}(S)$ est équivalente à une sous-catégorie
    \emph{pleine} de la catégorie des revêtements normaux de $S$ : ceux qui
    sont étales au-dessus de $U$ et dont toute composante irréductible domine
    une composante irréductible de $S$ rencontrant $U$.
\end{enumerate}

L'énoncé (b) est une équivalence sur une sous-catégorie pleine, non une
égalité : la sous-catégorie n'est pas décrite par une propriété fermée des
revêtements, et c'est précisément le théorème d'Abhyankar qui, dans le cas des
croisements normaux, la décrira ainsi.\footnote{Page 21 : « Le théorème
d'Abhyankar permet de caractériser ces revêtements par la condition d'être
étales sur $S - \operatorname{supp} D$, normaux en les points au-dessus de
$\operatorname{supp} D$ \dots ». C'est-à-dire : dans ce cas particulier, et
dans celui-là seulement, la structure devient une propriété. Voir plus bas,
page 21.}

Signalons enfin ce que la page 35 ajoute et qui n'est pas dans l'esquisse :
toutes les hypothèses de normalité sont \emph{superflues} pour la théorie
générale des revêtements étales de $(Z,G)$, qui a toutes les propriétés de
celle des revêtements étales d'un schéma, et est décrite par un groupe profini
--- extension de $G$ par $\pi_{1}(Z)$. Les hypothèses de normalité ne servent
qu'à assurer la pleine fidélité de $Z' \mapsto Z'/G$.

\pagerange{16}{18}
\section*{Exemple 2 : un système projectif de modèles (pages 16 à 18)}

Variante immédiate. Au lieu d'un seul couple $(Z, G)$, on se donne un système
projectif $(Z_{i}, G_{i})_{i \in I}$ de tels couples, $I$ ordonné filtrant. Le
système projectif des $G_{i}$ définit une catégorie fibrée $\mathcal{C}_{0}$
sur la catégorie associée à $I$, et l'on prend
$\mathcal{C} = \mathrm{Et}(S) \times \mathcal{C}_{0}$. Alors
\[
  \mathrm{Rev}^{\mathcal{R}}(S) \;\simeq\;
  \varinjlim_{i} \mathrm{Rev}^{\mathcal{R}_{i}}(S) ,
\]
limite inductive des catégories définies séparément par chaque $(Z_{i},G_{i})$.

Le cas intéressant est celui où les foncteurs de transition --- qui sont des
images inverses de revêtements étales à opérateurs --- sont pleinement fidèles,
de sorte que les $\mathrm{Rev}^{\mathcal{R}_{i}}(S)$ s'identifient à des
sous-catégories pleines qui épuisent la limite : c'est l'hypothèse (H1) sur cet
exemple. Il en est ainsi lorsque les morphismes de transition
$(Z_{i}, G_{i}) \to (Z_{j}, G_{j})$ sont tels que $G_{i} \to G_{j}$ soit
surjectif et que $Z_{i} \to Z_{j}$ induise un isomorphisme
$Z_{j} \simeq Z_{i}/N_{ij}$, $N_{ij}$ étant le noyau de $G_{i} \to G_{j}$. En
effet, tout revêtement étale $Z'_{j}$ de $Z_{j}$ vérifie alors
$Z'_{j} \simeq Z'_{i}/N_{ij}$, où $Z'_{i}$ est l'image inverse de $Z'_{j}$
sur $Z_{i}$.

La motivation est page 35 et elle est concrète : dans les applications on ne
dispose pas d'un seul modèle de ramification, mais d'une infinité --- tous les
$Z^{n}_{a}$ pour $a$ fixé et $n$ variable --- et la condition qu'on veut
imposer à un revêtement est d'avoir une ramification qui ne soit pire que celle
de \emph{l'un} des modèles.

\pagerange{18}{19}
\section*{Exemple 3, dans ses deux moutures (pages 18 et 19)}

L'exemple 3 est retapé, et le carton conserve les deux états. Ils ne sont pas
fusionnés ici : leur différence est la trace du travail, et elle va toujours
dans le même sens --- vers l'élimination d'une donnée arbitraire.

La situation qu'ils veulent couvrir est celle de la page 35 : les modèles de
ramification ne sont pas définis globalement --- dans le cas kummérien, les
diviseurs peuvent n'être pas principaux --- mais seulement localement, avec des
conditions de compatibilité disant que la restriction d'un modèle donné sur $U$
à $U \cap V$ n'a pas de ramification pire que celle d'un modèle donné sur $V$.

\subsection*{Première mouture : par une factorisation de $s$}

On suppose que $s$ se factorise
\[
  \mathcal{C} \xrightarrow{\ s'\ } \mathrm{Flop}_{S} \longrightarrow
  \mathrm{Fl}_{S} ,
\]
où $\mathrm{Flop}_{S}$ est la catégorie fibrée des schémas relatifs munis d'un
groupe fini d'opérateurs à droite, le second foncteur étant l'oubli des
opérations. Les morphismes de $\mathrm{Flop}_{S}$ sont les couples $(f, u)$ où
$f : Z \to Z'$ est un $S$-morphisme et $u : G_{S} \to G'_{S}$ un morphisme
\emph{localement constant, pas nécessairement constant}, de schémas en
groupes sur $S$, compatibles au sens évident.\footnote{Le tapuscrit écrit
« $f : X \to X'$ » là où ses propres objets sont $Z$ et $Z'$ ; c'est une faute
de frappe, non un changement d'objet, et elle est corrigée ici. La
transcription la conserve, comme elle doit.} La liberté laissée à $u$ est ce
qui permettra, au cas kummérien, de faire varier le groupe des racines de
l'unité d'un ouvert étale à l'autre.

On suppose alors deux choses. Pour toute flèche $M' \to M$ dans un
$\mathcal{C}_{S'}$, si $s(M') = (Z', G')$ et $s(M) = (Z,G)$, la flèche induit
un isomorphisme $Z \simeq Z'/H$, où $H$ est le noyau de $G'_{S} \to G_{S}$.
Et pour tout objet $M$, posant $s(M) = (Z,G)$, on se donne un homomorphisme
$G \to \operatorname{Aut}(M)$ tel que l'image $\bar{g}$ de $g$ vérifie
$s(\bar{g}) = (g_{Z}, \operatorname{int}(g))$, tout endomorphisme de $M$ étant
localement de cette forme ; le tout compatible aux images inverses et aux
flèches de $\mathcal{C}_{S'}$.

C'est cet homomorphisme $G \to \operatorname{Aut}(M)$ qui est la donnée
artificielle, et la seconde mouture s'en débarrasse.

\subsection*{Seconde mouture : par les groupes d'automorphismes}

On suppose maintenant, sur $\mathcal{R} = (\mathcal{C}, s)$ seule :

\begin{enumerate}
  \item[(1)] pour tout $S' \in \mathrm{Et}(S)$ et tout
    $M \in \mathcal{C}_{S'}$, le groupe $G = \operatorname{Aut}(M)$ est fini ;
  \item[(2)] pour toute flèche $M' \to M$ de $\mathcal{C}_{S'}$, il existe un
    unique morphisme de groupes $G' = \operatorname{Aut}(M') \to
    G = \operatorname{Aut}(M)$ rendant $M' \to M$ équivariant ; et si $H$ en est
    le noyau, $s(M') \to s(M)$ induit un isomorphisme
    $s(M')/H \simeq s(M)$ ;\footnote{Le sens de ce morphisme est fixé par deux
    flèches à l'encre portées au verso de la feuille (page 17) :
    $G' \to G$ au-dessus de $M' \to M$. Sans elles la phrase du recto est
    ambiguë. C'est le seul endroit du dossier où un verso décide d'une lecture.}
  \item[(3)] deux morphismes de $M'$ dans $M$ se déduisent l'un de l'autre par
    composition avec un élément de $G$ ;
  \item[(4)] deux objets $M$, $M'$ de $\mathcal{C}_{S'}$ sont localement
    majorés : il existe une famille surjective de morphismes étales
    $S'_{i} \to S'$ telle que $M_{S'_{i}}$ et $M'_{S'_{i}}$ soient tous deux
    majorés par un même objet $M''_{i}$ de $\mathcal{C}_{S'_{i}}$.
\end{enumerate}

Sous ces conditions, la théorie de la descente montre que la réalisation
géométrique $r : \mathrm{Rev}^{\mathcal{R}}(S) \to (\mathrm{Sch})_{/S}$ est
définie, et à valeurs dans les revêtements de $S$ dès que les $s(M)$ en sont
--- il suffirait même, pour appliquer la descente, que les $s(M)$ soient
affines sur $S'$. Et les critères de fidélité et de pleine fidélité de
l'exemple 1 se transposent : il suffit que chaque
$\bigl(s(M), \operatorname{Aut}(M)\bigr)$ satisfasse, sur son $S'$, les
conditions que $(Z,G)$ y satisfaisait.

La condition (4) est celle qui fait le travail de recollement : c'est elle qui
dit que les modèles locaux, donnés sur des ouverts étales différents, sont
comparables. Le paragraphe des « Remarques » de la page 32 remplacera ces
conditions par cinq conditions a) à e) portant sur $\mathcal{R}$ seule, où la
donnée de $s$ elle-même a fondu ; c'est le troisième état de l'exemple 3, et il
est lu plus bas, avec les pages 32 et 33.

\pagerange{19}{20}
\section*{Le cas kummérien (pages 19 et 20)}

Soit $(D_{i})_{i \in I}$ une famille de diviseurs de Cartier sur $S$. Pour
$S'$ étale sur $S$, on prend pour $\mathcal{C}_{S'}$ la catégorie dont les
objets sont les quadruples
\[
  M = (\underline{a},\, \underline{n},\, G,\, u) ,
\]
où $\underline{a} = (a_{i})_{i \in I}$ est une famille d'équations des
$D_{i\,S'}$, $\underline{n} = (n_{i})_{i \in I}$ une famille d'entiers
$\geqslant 1$ premiers aux caractéristiques résiduelles de $S'$, $G$ un groupe
fini, et $u : G_{S'} \xrightarrow{\ \sim\ } \mu_{\underline{n},\,S'}$ un
isomorphisme de schémas en groupes.

Les morphismes sont donnés par
\[
  \operatorname{Hom}_{S'}(M', M) \;=\;
  \begin{cases}
    \bigl\{\, \underline{b} = (b_{i})_{i \in I} \ :\
      \underline{b}^{\,\underline{n}} =
      \underline{a}'\,\underline{a}^{-1} \,\bigr\}
      & \text{si } \underline{n} \ \text{divise}\ \underline{n}' , \\[2pt]
    \varnothing & \text{sinon} ,
  \end{cases}
\]
la divisibilité s'entendant composante par composante.\footnote{Le tapuscrit
écrit « sauf si $\underline{n}' \geqslant \underline{n}$ ». Murre relève page 2
qu'il faudrait $\underline{n}' = \underline{m}\cdot\underline{n}$, et
Grothendieck répond page 37 : « when I wrote $n' \geqslant n$, I meant of
course the order relation of divisibility ». Le carton corrige donc lui-même,
à quarante pages de distance ; c'est la divisibilité qui est écrite ici.}

Le foncteur cartésien associe à $M$ le \emph{revêtement kummérien élémentaire}
\[
  Z^{\underline{n}}_{\underline{a}} \;=\;
  \operatorname{Spec} \mathcal{O}_{S'}[T_{i}\,;\ i \in I]\,/\,
  (T_{i}^{\,n_{i}} - a_{i}) ,
\]
considéré comme revêtement à groupe d'opérateurs $G$ grâce à
$u : G \simeq \mu_{\underline{n}}$. On vérifie qu'on est alors dans les
conditions de l'exemple 3.

Le foncteur $r$ est \emph{fidèle} ; il est \emph{pleinement fidèle} lorsque
$S$ est localement noethérien, normal aux points des
$\operatorname{supp} D_{i}$, les $D_{i}$ réduits, et toute composante
irréductible d'un $D_{i}$ distincte de toute composante d'un $D_{j}$ pour
$j \neq i$.

Le groupe $G$ et l'isomorphisme $u$ sont ici encore des données arbitraires :
Murre demande page 2 si les objets ne sont pas plutôt les couples
$(\underline{a}, \underline{n})$ avec $G(M) = \mu_{\underline{n}}$, et
Grothendieck le confirme page 37. La présentation intrinsèque de la page 33
permet en effet de « se dispenser d'introduire aussi des isomorphismes
$u : G_{S'} \to \mu_{\underline{n},\,S'}$ », et l'objet se réduit au couple
$(\underline{a}, \underline{n})$. C'est l'état auquel le carton reviendra vingt
fois ; sous cette forme, la catégorie fibrée $\mathcal{C}$ est un groupoïde,
et son champ associé est une \emph{gerbe} sous $\mu_{\underline{n}}$ --- ce que
la page 200 dira explicitement, et ce qu'on appelle aujourd'hui le champ des
racines $\sqrt[\underline{n}]{\underline{D}/S}$.\footnote{Aucun de ces mots
n'est disponible en 1967 : les gerbes sont de Giraud (1971), et les champs des
racines de Cadman et d'Abramovich--Graber--Vistoli (années 2000). L'objet, lui,
est ici, avec sa classe d'obstruction dans $H^{2}$ ; voir plus bas, pages 199
et 200, où Grothendieck l'écrit.}

Le tapuscrit continue à la page 21, avec le second cas kummérien
--- celui d'un diviseur $D$ à croisements normaux dont les composantes ne sont
définies que localement pour la topologie étale --- et c'est là que le mot
« modérément ramifié au sens large » est prononcé.

\pagerange{21}{22}
\section*{Le cas kummérien local, et la ramification modérée au sens large (pages 21 et 22)}

La page 20 s'arrêtait au milieu du second cas kummérien. Le voici entier.

Soit $D$ un diviseur sur $S$. On considère le faisceau, sur $\mathrm{Et}(S)$,
associé au préfaisceau qui à $S'$ associe l'ensemble des familles $(D_{i})$ de
diviseurs sur $S'$ telles que $D_{S'} = \sum_{i} D_{i}$ et que, pour tout
$s' \in S'$, les restrictions à $\operatorname{Spec}(\mathcal{O}_{S',s'})$ des
$D_{i}$ non nulles soient deux à deux disjointes. Une section de ce faisceau
donne, localement pour la topologie étale, une situation comme au cas
précédent, ce qui suffit à définir une donnée de ramification sur
$S$.\footnote{La page écrit « une donnée de ramification $\mathcal{R}$ sur
$B$ », pour « sur $S$ » ; c'est une coquille de frappe, corrigée ici. Le
tapuscrit appelle d'ailleurs « canularesque » la condition locale sur les
$(D_{i})$ --- elle n'est là que pour qu'on ait bien un préfaisceau.}

Deux cas particuliers méritent un nom, car $\mathcal{C}$ s'y décrit en termes
du seul $D$ :

\begin{enumerate}
  \item[1\textsuperscript{o})] $S$ localement noethérien, $D$ à croisements
    normaux ; on prend pour $D_{i}$ les composantes irréductibles locales
    régulières de $D$, qui ne sont définies que localement pour la topologie
    étale ;
  \item[2\textsuperscript{o})] $S$ au-dessus d'un préschéma $T$, $D$ à
    croisements normaux \emph{relativement à $T$} ; on prend pour $D_{i}$ les
    composantes locales lisses sur $T$.
\end{enumerate}

Dans le cas 1\textsuperscript{o}, $\mathrm{Rev}^{\mathcal{R}}(S)$ est
équivalente à une sous-catégorie pleine de la catégorie des revêtements de $S$,
et l'on appelle ses objets les \emph{revêtements modérément ramifiés au sens
large} relativement à $D$. C'est une notion strictement plus large que celle de
revêtement modérément ramifié par rapport à la famille réduite au seul $D$, et
c'est le \emph{théorème d'Abhyankar} qui la caractérise : ce sont les
revêtements étales sur $S - \operatorname{supp} D$, normaux aux points
au-dessus de $\operatorname{supp} D$, dont toute composante irréductible en
domine une de $S$, et qui sont modérément ramifiés au sens classique aux points
maximaux de $D$.

C'est le seul endroit du carton où la \emph{structure} redevient une
\emph{propriété}, et il vaut d'être noté comme tel : la caractérisation
d'Abhyankar est ce qui rend le cas des croisements normaux exceptionnellement
maniable, et son absence ailleurs est la raison d'être de tout l'appareil.

Le cas 2\textsuperscript{o} attend encore une vérification --- « sauf erreur
(il faudrait le vérifier) » ---, et l'on parlerait alors de revêtements
« modérément ramifiés pour $D$, relativement à la base $T$ ». Lorsque $T$ est
régulier les deux notions coïncident. L'intérêt du cas relatif est indiqué en
une phrase : c'est avec la ramification modérée au sens large que le théorème
de spécialisation pour le groupe fondamental modéré se démontre par la
technique de SGA 1, pour $S$ propre et lisse sur $T$ et $D$ à croisements
normaux relativement à $T$.\footnote{Ce théorème de spécialisation est celui de
Michèle Raynaud, SGA 1 XIII ; le tapuscrit y renvoie explicitement page 114 du
carton, « exposé de Mme Raynaud ».}

\pagerange{22}{24}
\section*{La donnée est auxiliaire : l'analogie du site et du topos (pages 22 à 24)}

Pour $S'$ variable dans $\mathrm{Et}(S)$ on définit
$\mathrm{Rev}^{\mathcal{R}}(S')$ en restreignant $\mathcal{R}$ à
$\mathrm{Et}(S') = \mathrm{Et}(S)_{/S'}$, d'où une catégorie fibrée
$\underline{\mathrm{Rev}}^{\mathcal{R}}$ sur $\mathrm{Et}(S)$ --- qui est même
un \emph{champ} --- et un foncteur cartésien
\[
  (\ast) \qquad
  r : \underline{\mathrm{Rev}}^{\mathcal{R}} \longrightarrow
  \underline{\mathrm{Rev}} ,
\]
où $\underline{\mathrm{Rev}}$ désigne la catégorie fibrée des revêtements des
objets de $\mathrm{Et}(S)$.\footnote{On suppose ici, pour tout $S'$, que la
réalisation géométrique sur $\mathcal{C}_{S'}$ est à valeurs dans les
revêtements de $S'$, et que le foncteur $r$ sur
$\mathrm{Rev}^{\mathcal{R}}(S')$ est bien défini --- c'est l'hypothèse (H2)
ci-dessus, prise fibre à fibre.}

C'est ici que le tapuscrit dit ce qu'il faut penser de tout ce qui précède, et
la phrase doit être citée :

\begin{quote}
Il faut insister sur le fait que pour nous, la « donnée de ramification »
$\mathcal{R}$ est un objet de nature essentiellement auxiliaire, destiné à
définir la catégorie fibrée $\underline{\mathrm{Rev}}^{\mathcal{R}}$ (qui est
même un champ) et le foncteur cartésien $r$ de $(\ast)$ ; c'est $(\ast)$ qui
est l'objet véritable de l'étude.
\end{quote}

On dira donc que deux données $\mathcal{R}$, $\mathcal{R}'$ sont
\emph{essentiellement équivalentes} si elles définissent des foncteurs $(\ast)$
équivalents, par une équivalence induisant l'identité sur
$\underline{\mathrm{Rev}}$.

\subsection*{Trois exemples d'équivalence}

Le premier est formel. Sous les conditions de l'exemple 3, soit $\mathcal{C}'$
une sous-catégorie fibrée de $\mathcal{C}$ telle que tout objet $M$ d'un
$\mathcal{C}_{S'}$ soit localement majoré par un objet de $\mathcal{C}'$, et
munissons $\mathcal{R}'$ de la réalisation géométrique induite : alors
$\mathcal{R}$ et $\mathcal{R}'$ sont équivalentes. Autrement dit,
\emph{seul le crible compte, pas le catalogue}.

Le deuxième compare les exemples 2 et 3 dans le cas kummérien. Si les $D_{i}$
admettent des équations globales $c_{i}$, on peut considérer le système
projectif des $Z^{\underline{n}}_{\underline{c}}$ pour $\underline{n}$
variable, premier à la caractéristique de $S$ (supposée exister), avec le
système projectif des groupes $\mathbb{Z}/n_{i}\mathbb{Z}$ ; si l'on dispose
d'un isomorphisme entre ce dernier et le système des $\mu_{\underline{n}}$ ---
il en existe lorsque $S$ est un schéma sur un corps algébriquement clos --- on
est dans les conditions de l'exemple 2, et la donnée obtenue est équivalente à
la donnée kummérienne de l'exemple 3. Les deux exemples décrivent donc, dans ce
cas, la même chose ; ce qui les distingue est l'existence globale des racines
de l'unité.

Le troisième est une torsion. Soit $\mathcal{R}$ la donnée de l'exemple 1,
définie par $(Z,G)$. Donnons-nous un revêtement principal $P$ de $S$ de groupe
$G'$ opérant à droite, et supposons que $G$ opère à gauche sur $P$, librement,
en commutant à $G'$. Alors $Z' = Z \times^{G} P$, muni de $G'$, définit une
donnée $\mathcal{R}'$ équivalente à $\mathcal{R}$ : on le voit d'abord quand
$P = G'_{S}$ avec l'opération régulière droite et $G \hookrightarrow G'$, puis
par localisation étale.

\subsection*{« On soit infailliblement conduit \dots »}

Et vient alors la phrase qui donne au dossier son programme :

\begin{quote}
On peut se proposer de normaliser la notion de donnée de ramification, pour
éliminer l'arbitraire dans le choix d'un $\mathcal{R}$. Mais alors il semble
qu'on soit infailliblement conduit à prendre comme catégorie $\mathcal{C}$ la
catégorie $\mathrm{Rev}^{\mathcal{R}}$ elle-même qu'on se propose de construire
et d'étudier ! (Comparer avec la situation analogue pour les relations entre
sites et topos, l'objet auxiliaire étant le site, l'objet intrinsèque de nature
géométrique étant le topos associé à ce site.)
\end{quote}

Il faut lire cela comme un problème ouvert posé à soi-même, et il est résolu
dans le même carton, deux cents pages plus loin : le \emph{domaine de
ramification} des pages 190 à 200 est un champ de catégories multigaloisiennes
au-dessus de $\mathrm{Et}(S)$, muni d'une réalisation géométrique continue,
c'est-à-dire exactement la catégorie elle-même prise comme donnée. Le carton ne
date pas ces pages-là et cette lecture ne le fera pas ; mais la boucle est
close, et elle l'est dans un seul carton.

\pagerange{24}{25}
\section*{La catégorie et son groupe fondamental (pages 24 et 25)}

\subsection*{Exactitude}

La catégorie $\mathrm{Rev}^{\mathcal{R}}(S)$ hérite de toutes les propriétés
d'exactitude établies pour les revêtements étales d'un schéma. De façon
précise : les limites projectives finies et les limites inductives finies y
existent et satisfont les mêmes propriétés d'exactitude que dans la catégorie
des ensembles ; tout morphisme s'y factorise en un épimorphisme suivi d'un
monomorphisme ; épimorphismes et monomorphismes sont effectifs ; et tout
monomorphisme est un isomorphisme sur un facteur direct.

Ce dernier énoncé est exactement l'axiome que la page 209 numérotera
\textbf{G3}, et que SGA 1 V numérote de même : \emph{tout morphisme se
factorise en un épimorphisme suivi d'un isomorphisme sur un sommande direct}.
Il est ici obtenu comme une propriété, et il sera là-bas posé comme un axiome ;
c'est le même énoncé, à onze ans et cent quatre-vingts pages de distance.

Un avertissement, qui n'est pas une remarque de détail : la réalisation
géométrique $r$ commute aux limites inductives finies --- en particulier aux
sommes --- mais \emph{pas} aux limites projectives finies, et en particulier
pas au produit de deux objets. Le produit de deux revêtements de type
$\mathcal{R}$ n'est donc pas le produit fibré de leurs réalisations au-dessus
de $S$. C'est le fait que la page 31 signalera comme la cause de la
complication de l'un des lemmes de Murre.

\subsection*{Le foncteur fibre}

Pour appliquer la théorie de Galois abstraite, il faut de plus un
\emph{foncteur fibre}
\[
  F : \mathrm{Rev}^{\mathcal{R}}(S) \longrightarrow (\text{ensembles finis}) ,
\]
c'est-à-dire un foncteur exact.

Dans l'exemple 1, tout point géométrique de $Z$ en fournit un ; il est
\emph{conservatif} si toute partie ouverte et fermée de $Z$ stable par $G$ est
vide ou égale à $Z$. Dans l'exemple 2, on prend un point géométrique de
$Z = \varprojlim Z_{i}$.

Dans le cas général de l'exemple 3, il faut travailler un peu. On choisit un
point géométrique $s$ de $S$ qui n'appartienne pas à \emph{l'ensemble de
ramification} de $\mathcal{R}$ : pour tout $S'$ dans $\mathrm{Et}(S)$ et tout
$M$ de $\mathcal{C}_{S'}$, le revêtement à opérateurs $\bigl(s(M),
\operatorname{Aut}(M)\bigr)$ de $S'$ doit être \emph{principal} aux points de
$S'$ au-dessus de l'image de $s$. Alors, pour tout objet $T$ de
$\mathrm{Rev}^{\mathcal{R}}(S)$, $r(T)$ est étale sur $S$ en $s$, et
\[
  F : T \longmapsto r(T)(s)
\]
est exact. Il est conservatif si et seulement si le site $\mathcal{C}$, muni de
la topologie image inverse de celle de $\mathrm{Et}(S)$, est \emph{connexe} ;
et si l'on suppose que $s(M)/\!\operatorname{Aut}(M) \to S'$ est un isomorphisme
pour tout $M$ --- « ce qui sera sans doute le cas dans tous les cas
intéressants » --- cela revient à dire que $S$ est connexe.

\subsection*{Le groupe}

Sous ces conditions, la théorie de Galois abstraite fournit un groupe profini
\[
  \pi_{1}^{\mathcal{R}}(S, s)
\]
et une équivalence entre $\mathrm{Rev}^{\mathcal{R}}(S)$ et la catégorie des
ensembles finis munis d'une action continue de ce groupe.

La connexité n'est pas un détail technique : c'est elle, et elle seule, qui
permet de remplacer un groupoïde par un groupe. Sans elle,
$\mathrm{Rev}^{\mathcal{R}}(S)$ n'est pas galoisienne mais seulement, comme la
page 30 le dira, \emph{multigaloisienne} --- un produit de catégories
galoisiennes indexé par les composantes connexes de l'objet final, donc un
\emph{groupoïde} profini et non un groupe.\footnote{Le mot apparaît page 30
sans définition, comme un terme déjà disponible : « Comme
$\mathrm{Rev}^{\mathcal{R}}(S)$ est en tous cas une catégorie
multigaloisienne \dots ». Sa définition axiomatique est aux pages 208 à 216 du
même carton, non datées ; ce sont les conditions A, B et C, et la
caractérisation (iii) y est précisément « $C$ équivalente à un produit
$\prod_{i} C_{i}$, les $C_{i}$ galoisiennes ». Le carton contient donc l'usage
du mot et son axiomatique, sans dire lequel des deux est venu d'abord, et cette
lecture ne le décide pas.}

\pagerange{25}{29}
\section*{Fonctorialité en $S$ et en $\mathcal{R}$ (pages 25 à 29)}

Soit $f : T \to S$ un morphisme de schémas, d'où un foncteur continu
$f^{\circ} : \mathrm{Et}(S) \to \mathrm{Et}(T)$. Donnons-nous une donnée
$\mathcal{R}_{S} = (\mathcal{C}_{S}, s_{S})$ sur $S$, une donnée
$\mathcal{R}_{T} = (\mathcal{C}_{T}, s_{T})$ sur $T$, un foncteur
$M \mapsto M \times_{S} T$ de $\mathcal{C}_{S}$ dans $\mathcal{C}_{T}$
transformant morphismes cartésiens en morphismes cartésiens et relevant
$f^{\circ}$, et enfin un morphisme entre les deux composés

\begin{tikzcd}[column sep=huge]
\mathcal{C}_{S} \arrow[r] \arrow[d, "s_{S}"'] & \mathcal{C}_{T} \arrow[d, "s_{T}"] \\
(\mathrm{Sch})_{/S} \arrow[r, "X \mapsto X \times_{S} T"'] & (\mathrm{Sch})_{/T}
\end{tikzcd}

soit
\[
  u(M) : s_{T}(M \times_{S} T) \longrightarrow s_{S}(M) \times_{S} T ,
\]
morphisme au-dessus de $f^{\circ}(S')$ pour chaque $M$ sur $S'$.

On en déduit un foncteur
$\mathrm{Rev}^{\mathcal{R}_{S}}(S) \to \mathrm{Rev}^{\mathcal{R}_{T}}(T)$ et un
homomorphisme de foncteurs composés

\begin{tikzcd}[column sep=huge]
\mathrm{Rev}^{\mathcal{R}_{S}}(S) \arrow[r] \arrow[d, "r"'] & \mathrm{Rev}^{\mathcal{R}_{T}}(T) \arrow[d, "r"] \\
\mathrm{Rev}(S) \arrow[r, "X \mapsto X \times_{S} T"'] & \mathrm{Rev}(T)
\end{tikzcd}

Cet homomorphisme est un isomorphisme dans certains cas : sous les conditions
de l'exemple 3, lorsque, pour tout sous-groupe $H$ de
$G = \operatorname{Aut}(M)$, le passage au quotient $s(M)/H$ commute au
changement de base $T \to S$ --- ce qui est vérifié si $T \to S$ est plat, ou
dans les cas kummériens a) pourvu que les images inverses des $D_{i}$ existent
et servent à définir la donnée sur $T$.

Le point de la construction est que \emph{le foncteur ainsi obtenu est exact}.
Il en résulte, par les sorites de la théorie de Galois abstraite, que si $t$
est un point géométrique de $T$ non ramifié pour $\mathcal{R}_{T}$ dont l'image
$s$ dans $S$ est non ramifiée pour $\mathcal{R}_{S}$, et si $u$ induit un
isomorphisme sur les fibres géométriques en $t$, on a un homomorphisme
canonique
\[
  \pi_{1}^{\mathcal{R}_{T}}(T, t) \longrightarrow
  \pi_{1}^{\mathcal{R}_{S}}(S, s) .
\]
Cela s'applique aux données kummériennes des trois types dès que $T \to S$ est
régulier, respectivement lisse, de sorte que les images inverses des $D_{i}$
vérifient encore les conditions de croisements normaux.

Sur la transitivité pour un composé $U \to T \to S$, le tapuscrit renonce ---
« on recule devant la tâche \dots , jetant provisoirement un voile pudique sur
ces abîmes » --- et cette lecture n'y supplée pas.\footnote{Le renoncement est
significatif et non anecdotique : c'est la 2-catégorialité de la situation qui
rend l'écriture pénible, et c'est exactement ce que les pages 196 à 198
finiront par assumer en posant que les domaines de ramification sur $S$ forment
une \emph{2-catégorie}, et que celles-ci forment pour $S$ variable une
2-catégorie fibrée.}

\subsection*{Deux cas particuliers}

Le premier est le changement de base \emph{étale}, avec $\mathcal{R}_{T}$
induite par $\mathcal{R}_{S}$ : le carré est déjà connu, puisqu'on dispose du
foncteur cartésien $(\ast)$ sur $\mathrm{Et}(S)$, et l'homomorphisme de
foncteurs composés y est un isomorphisme. Si de plus $T$ est fini sur $S$, avec
$S$ et $T$ connexes, et sous les conditions de l'exemple 3, l'homomorphisme sur
les groupes fondamentaux est un \emph{monomorphisme}, d'image le sous-groupe de
$\pi_{1}^{\mathcal{R}}(S,s)$ défini par le revêtement pointé $T$ de $S$.

Le second est le changement de base \emph{entier, radiciel et surjectif}, avec
$\mathcal{R}_{T}$ image inverse de $\mathcal{R}_{S}$. Alors le foncteur est
\emph{trivialement} une équivalence, par le résultat analogue pour les
revêtements étales de $X$ et $X \times_{S} T$ ; en particulier, sur un point
géométrique non ramifié de $T$, on obtient un isomorphisme des groupes
fondamentaux.

\subsection*{L'image inverse d'une donnée, et pourquoi elle est encombrante}

Grothendieck signale une difficulté qu'il ne résout pas. Si $\mathcal{R}$ est
définie par $(Z,G)$ comme à l'exemple 1, son image inverse par $T \to S$ est
définie par $(Z \times_{S} T,\, G)$ ; c'est immédiat. Mais si $\mathcal{R}$ est
la donnée kummérienne associée à un système $(D_{i})$, l'image inverse au sens
strict n'est \emph{pas} la donnée associée au système $(D_{i\,T})$, même à
équivalence de catégories fibrées près : elle lui est seulement
« essentiellement équivalente » au sens ci-dessus.

Il en tire deux propositions : ou bien introduire la notion d'image inverse
avec tout son poids, ou bien --- c'est le parti qu'il juge le plus raisonnable
--- donner des conditions axiomatiques sur le carré $(\ast)$, sous les
conditions de l'exemple 3, qui assurent \emph{moralement} que
$\mathcal{R}_{T}$ est essentiellement équivalente à l'image inverse de
$\mathcal{R}_{S}$, et qui s'appliquent directement aux cas kummériens. Il
ajoute qu'il serait raisonnable, « assez rapidement après les définitions
générales », de se borner aux données du type de l'exemple 3 : ce sont les
seules qu'on rencontrera jamais et les seules où l'on sache construire $r$.

\subsection*{Descente et passage à la limite}

Deux énoncés sont annoncés et non démontrés, avec l'usage qu'on veut en faire.

Un morphisme fidèlement plat et quasi-compact $T \to S$ est de \emph{descente
effective} pour la catégorie fibrée des revêtements de type $\mathcal{R}_{S}$
sur le diagramme $(S, T, T\times_{S}T, T\times_{S}T\times_{S}T)$. En
particulier la descente fidèlement plate vaut pour les revêtements kummériens,
ce qui est utilisé sans démonstration dans les notes de Murre pour la descente
à partir d'un anneau complété.

Il y a lieu ensuite d'énoncer un \emph{passage à la limite} pour un système
projectif $(X_{i})$ de schémas quasi-compacts et quasi-séparés à morphismes de
transition affines : cela ramène l'étude au voisinage d'un point $s$ à celle de
$\operatorname{Spec}(\mathcal{O}_{S,s})$. Ainsi, pour $S$ localement noethérien
normal muni d'une famille $(D_{i})$ de diviseurs réduits --- où la catégorie
des revêtements kummériens est une sous-catégorie de celle des revêtements ---
un revêtement $X$ de $S$ est kummérien si et seulement si le revêtement induit
sur $\operatorname{Spec}(\mathcal{O}_{S,s})$ l'est pour tout $s$, et l'on peut
y remplacer $\mathcal{O}_{S,s}$ par son hensélisé, par son hensélisé strict,
voire par son complété si les hypothèses restent stables par complétion --- par
exemple si $\mathcal{O}_{S,s}$ est excellent. Sans hypothèse de normalité ni de
réduction, les mêmes observations valent à condition de se borner aux
revêtements modérément ramifiés \emph{galoisiens}, ce qui est le sujet du
paragraphe suivant.

\pagerange{30}{31}
\section*{Revêtements galoisiens, et pourquoi la catégorie est indispensable (pages 30 et 31)}

\subsection*{La définition, complétée}

Un \emph{revêtement galoisien de type $\mathcal{R}$ de $S$, à groupe de Galois
$G$}, est un objet $X$ de $\mathrm{Rev}^{\mathcal{R}}(S)$ sur lequel le groupe
fini $G$ opère à droite de telle façon que, \emph{dans la catégorie}
$\mathrm{Rev}^{\mathcal{R}}(S)$, on obtienne un revêtement principal de l'objet
final : le morphisme $X \times G \to X \times X$ est un isomorphisme, et
$X \to e$ est couvrant.

Il faut y ajouter une condition que la page ne porte pas, sans quoi la
définition explicite qui suit ne coïncide pas avec celle-ci : \emph{l'action de
$G$ doit commuter aux actions des groupes d'automorphismes dont les modèles
sont déjà munis}, aussi bien qu'à la localisation étale sur $S$.\footnote{C'est
la question de Murre, page 2 : « It seems to me that one must require that the
actions of $G$ and $\mathcal{G}$ commute; is this correct? » La réponse, page
36, est oui : « the action of $G$ commutes to étale localization on $S$
\emph{and} to the action of the groups of automorphisms given on the various
$Z$'s ». Le paragraphe 4 lui-même ne la porte pas ; elle est rétablie ici, à sa
place logique, parce que sans elle l'explicitation qui suit --- et le critère
de pleine fidélité de la page 31 --- ne veut rien dire. Voir aussi les
pages 1 à 11.}

Sous cette condition, la notion s'explicite de trois façons équivalentes :
comme ci-dessus ; en termes des groupes fondamentaux, puisque
$\mathrm{Rev}^{\mathcal{R}}(S)$ est en tous cas une catégorie
multigaloisienne ; et en revenant à la définition, en exigeant que, pour tout
$M \in \mathcal{C}_{S'}$, le revêtement $X(M)$ de $s(M)$ soit \emph{principal
de groupe $G$}.

Par abus de langage, on appellera aussi revêtement galoisien de type
$\mathcal{R}$ le revêtement à opérateurs $\bigl(r(X), G\bigr)$. L'abus est sans
danger lorsque le foncteur de la catégorie des revêtements galoisiens de type
$\mathcal{R}$ --- morphismes restreints aux \emph{isomorphismes} --- dans celle
des revêtements à opérateurs de $S$ est pleinement fidèle : alors, pour un
revêtement à opérateurs $(T,G)$ donné, il existe au plus un $(X,G)$ galoisien
de type $\mathcal{R}$ dont il provienne, à isomorphisme canonique près.

\subsection*{La condition nécessaire et suffisante}

Voici la condition, d'abord dans le cas de l'exemple 1. Pour tout localisé
strict $S'$ d'un point de $S$, soit $Z'_{0}$ une composante connexe de
$Z' = Z \times_{S} S'$ et $H$ son stabilisateur dans $G$. On demande : pour
deux sous-groupes distingués $H_{1}, H_{2}$ de $H$ et tout isomorphisme
$u : Z'_{0}/H_{1} \to Z'_{0}/H_{2}$ compatible avec un isomorphisme
$\varphi : H/H_{1} \to H/H_{2}$, il existe $g \in H$ tel que
$\operatorname{int}(g)$ applique $H_{1}$ sur $H_{2}$, induise $\varphi$ par
passage au quotient, et tel que $u$ soit induit par $g_{Z'_{0}}$.

Dans le cas général de l'exemple 3, on exige cela pour tous les couples
$\bigl(s(M), \operatorname{Aut}(M)\bigr)$ sur les divers $S'$ de
$\mathrm{Et}(S)$.

Cette condition est satisfaite dans le cas kummérien --- Grothendieck renvoie
aux énoncés 1.10 et 1.15 des notes de Murre --- et la démonstration du critère
est esquissée page 37 : la question est locale pour la topologie étale, donc on
se ramène au cas strictement local, où la classification est explicite (voir
plus bas).

\subsection*{Le passage en anglais}

C'est ici que le tapuscrit change de langue, en pleine page, parce que
Grothendieck est en train de répondre à Murre. Le passage est le centre du
dossier et il tient en quatre affirmations.

\begin{enumerate}
  \item Ce phénomène explique pourquoi il a été si réticent --- à tort ---
    à définir les revêtements modérément ramifiés hors du cas galoisien :
    \emph{dans le cas non galoisien, ce ne sera pas un revêtement satisfaisant
    des conditions supplémentaires, mais un revêtement muni d'une structure
    supplémentaire} ; alors que dans le cas galoisien, la structure est
    uniquement déterminée si elle existe, et l'on peut ne pas la voir.
  \item Tous les morphismes au sens naïf entre revêtements modérés, même
    galoisiens, ne doivent pas être considérés : seuls ceux qui viennent de
    morphismes de $\mathrm{Rev}^{\mathcal{R}}(S)$. Le contre-exemple à 1.16 (pages 7 et 8) montre qu'il faut y prendre garde.
  \item La notion naturelle de \emph{produit} de deux revêtements modérés,
    galoisiens ou non, n'est pas le produit fibré naïf au-dessus de $S$ ---
    ce qui, ajoute-t-il, explique la complication de la démonstration d'un des
    lemmes de Murre.
  \item Donc : « This shows rather clearly that even in the Galois case, the
    consideration of the category $\mathrm{Rev}^{\mathcal{R}}(S)$ is
    undispensable for a good comprehension of what is really going on. »
\end{enumerate}

\pagerange{32}{33}
\section*{Le troisième état de l'exemple 3 : cinq conditions sur $\mathcal{R}$ (pages 32 et 33)}

Deux feuillets séparés reprennent l'exemple 3 et l'ouvrent par un jugement :
« La présentation de cet exemple n'est ni élégante ni naturelle, la donnée de
$s$ et de \dots\ étant un élément de structure artificiel. » Suit la
présentation intrinsèque, en cinq conditions sur $\mathcal{R}$ seule.

\begin{enumerate}
  \item[a)] Pour tout objet $M$ d'un $\mathcal{C}_{S'}$, le faisceau étale
    associé au préfaisceau des automorphismes de $M$ sur $\mathrm{Et}(S')$ est
    représentable par un schéma en groupes $G(M)$ fini et étale sur $S'$,
    déterminé à isomorphisme unique près. Il opère sur le schéma relatif
    $s(M)$.
  \item[b)] Pour tout $S'$, tout morphisme de $\mathcal{C}_{S'}$ est un
    épimorphisme et tout endomorphisme un automorphisme ; et pour tout
    $u : M' \to M$ et tout automorphisme $g'$ de $M'$, il existe un
    automorphisme $g$ de $M$ tel que $gu = ug'$. Cet automorphisme est alors
    unique, et $u$ détermine un unique homomorphisme de schémas en groupes
    $G(M') \to G(M)$ le rendant équivariant.
  \item[c)] Cet homomorphisme $G(M') \to G(M)$ est surjectif, et si $H$ en est
    le noyau, $s(M') \to s(M)$ induit un isomorphisme $s(M')/H \simeq s(M)$ ---
    les $s(M)$ étant supposés relativement affines, de sorte que le quotient a
    un sens.
  \item[d)] Deux homomorphismes de $M'$ dans $M$ se déduisent l'un de l'autre,
    localement pour la topologie étale, par composition avec un automorphisme
    de $M$.
  \item[e)] Étant donnés $M$, $M'$ dans $\mathcal{C}_{S'}$, il existe une
    famille couvrante $S'_{i} \to S'$ dans $\mathrm{Et}(S)$ telle que les
    images de $M$ et $M'$ sur $S'_{i}$ soient dominées par un même objet
    $M''_{i}$ de $\mathcal{C}_{S'_{i}}$.
\end{enumerate}

Trois choses sont à voir. La première est le mouvement : le groupe
d'automorphismes, qui n'était qu'un groupe abstrait à la première mouture et un
groupe fini à la seconde, est ici un \emph{schéma en groupes}, obtenu par
faisceautisation du préfaisceau des automorphismes --- donc extrait de
$\mathcal{C}$ et non ajouté à elle. La condition c) fait alors de $s(M') \to
s(M)$ un torseur sous $H$, ce que la page ébauche avant de biffer.

La deuxième est ce que cela permet : « on peut, dans la description des données
de ramification kummériennes, se dispenser d'introduire aussi des isomorphismes
$u : G_{S'} \to \mu_{\underline{n},\,S'}$ ». Le quadruple
$(\underline{a}, \underline{n}, G, u)$ des pages 19 et 20 se réduit au couple
$(\underline{a}, \underline{n})$, et c'est la réponse à la question de Murre.

La troisième est la conditions d) et e) prises ensemble : $\mathcal{C}_{S'}$
est localement un groupoïde connexe. C'est la structure que les pages 198 à 200
appelleront une \emph{gerbe}, et sous cette forme la donnée kummérienne est la
gerbe des racines $\underline{n}$-ièmes des équations des $D_{i}$.

\pagerange{34}{38}
\section*{La lettre d'accompagnement : le cas strictement local (pages 34 à 38)}

La lettre qui accompagne le tapuscrit est une explication, non un complément :
elle reprend l'exemple 1 dans le langage des corps de fonctions (lu plus haut), puis fait le calcul strictement local, qui est ce que l'esquisse omet.

Plaçons-nous sur $S$ strictement local. Pour des revêtements définis sur $S$
tout entier, la situation générale de l'exemple 3 se ramène à celle de
l'exemple 2 ; et, quitte à remplacer les $Z_{i}$ par des composantes connexes,
on peut les supposer connexes. Pour $Z_{i}$ fixé, les revêtements étales de
$Z_{i}$ sont triviaux, puisque $Z_{i}$ est lui aussi strictement local. Donc :
\[
  \bigl\{\text{revêtements étales de } (Z_{i}, G_{i})\bigr\}
  \;\simeq\;
  \bigl\{\text{ensembles finis } E \text{ munis d'une action de } G_{i}\bigr\} ,
\]
et la réalisation géométrique sur $S$ de l'objet associé à $E$ est
$Z_{i} \times^{G_{i}} E$.

Cela donne, pour $S$ quelconque et localement pour la topologie étale, une
description concrète de ce que signifie, pour un revêtement $S'$ de $S$, d'être
associé à un modèle $M$ : \emph{cela signifie qu'il est localement décrit par
une représentation d'un groupe d'inertie sur un ensemble fini} ; et la
structure supplémentaire consiste à se donner de telles représentations locales
qui se recollent. La description a un sens dès que $r$ est fidèle, sans qu'on
ait à le supposer pleinement fidèle.

Pour les revêtements galoisiens de groupe $H$, le calcul est le suivant, et il
est joli. Dans le cas strictement local, pour $Z_{i}$ fixé, la catégorie de
ces revêtements est équivalente à celle des ensembles finis $E$ munis d'une
action à gauche de $G_{i}$ et d'une action à droite de $H$ qui commutent, $E$
étant un torseur sous $H$. À isomorphisme près, ils sont classés par
\[
  \operatorname{Hom}(G_{i}, H) \big/ H ,
\]
homomorphismes de groupes modulo automorphisme intérieur.\footnote{Le calcul
est immédiat et vaut d'être refait : $E$ étant un $H$-torseur, le choix d'un
point l'identifie à $H$ opérant sur lui-même à droite ; les automorphismes de
$H$ comme $H$-ensemble à droite sont les translations à gauche ; une action à
gauche de $G_{i}$ qui commute aux translations à droite est donc donnée par un
homomorphisme $G_{i} \to H$, et le changement du point choisi le conjugue.}

C'est de cette description que sort le critère de pleine fidélité du foncteur
« réalisation géométrique » sur les revêtements galoisiens de type
$\mathcal{R}$ : la question est locale pour la topologie étale, ce qui ramène
au cas strictement local, où l'on vient de tout calculer.

La lettre porte enfin, en passant, une remarque de vocabulaire qui n'est pas
anodine : « By the way, one should rather say \emph{category of ramification
data} ». Le carton finira par lui donner raison.

\pagerange{39}{40}
\section*{Un dossier s'ouvre : « Finitude (Pbs et conjectures) » (pages 39 et 40)}

Les pages 39 et 40 ouvrent une chemise nouvelle, de la main de
Grothendieck, et sans rapport apparent avec ce qui précède --- sauf qu'elle
demande, du groupe fondamental, ce que le tapuscrit vient de construire.

Trois énoncés, sur un corps $k$ algébriquement clos.

\begin{enumerate}
  \item[a)] $X$ propre sur $k$ $\Longrightarrow$ $\pi_{1}(X)$ topologiquement
    de type fini.
  \item[b)] $X$ de type fini sur $k$ $\Longrightarrow$ $\pi'_{1}(X)$ de
    \emph{présentation finie}, où $\pi'_{1}$ désigne le plus grand quotient
    premier à la caractéristique.
  \item[c)] $X$ de type fini sur $k$ $\Longrightarrow$ $\pi^{t}_{1}(X)$
    topologiquement de type fini, où $\pi^{t}_{1}$ classifie les revêtements
    modérément ramifiés le long du bord d'une compactification.
\end{enumerate}

Et une note en marge : « c) plus général que a) ». En effet, pour $X$ propre il
n'y a pas de bord, donc $\pi^{t}_{1}(X) = \pi_{1}(X)$, et c) redonne
a).\footnote{Le sens précis de $\pi_{1}^{t}$ est indiqué en marge de la page
d'une écriture très abîmée ; la transcription y laisse plusieurs mots
illisibles. La lecture retenue ici est celle que le reste du carton impose ---
notamment les pages 112 à 115, où $\pi_{1}^{t}(U)$ est défini comme le quotient
de $\pi_{1}(U)$ classifiant les revêtements modérément ramifiés le long des
composantes de codimension 1 du bord.}

L'état des démonstrations est noté en regard, et il faut le lire comme un
tableau de bord. L'énoncé a) est \emph{démontré} : réduction au cas normal,
Bertini pour se ramener à une courbe, puis théorie de la spécialisation du
groupe fondamental et méthode transcendante. Pour b), « le cas type fini est
prouvé » : réduction au cas normal par descente, puis au cas où les
singularités sont « jolies au sens Artin », suite exacte des $\pi_{1}$,
c'est-à-dire théorème de spécialisation, puis passage à une courbe normale.
Reste ce que la résolution des singularités permettrait de faire, et le cas
d'une surface projective avec un diviseur à croisements normaux.

Ces trois questions occupent les pages 41 à 76, où l'on trouvera le théorème de
finitude relatif, le groupe fondamental birationnel, et l'énoncé --- démontré
là --- que $\pi'_{1}(X \times Y) = \pi'_{1}(X) \times \pi'_{1}(Y)$.

\pagerange{41}{42}
\section*{Deux complétions d'un croisement normal, et le groupe fondamental local (pages 41 et 42)}

\subsection*{Le feuillet de calcul}

Soit $A$ un anneau noethérien, $x, y \in A$, tels que $A/x$ et $A/y$ soient
réguliers intègres de dimension 1, et supposons $A$ complet pour la topologie
définie par $xy$. C'est le modèle local d'un croisement normal : deux branches
régulières se coupant en un point.

On dispose alors de trois complétions,
\[
  \hat{A}^{(x)} \;=\; k[y][[x]] , \qquad
  \hat{A}^{(y)} \;=\; k[x][[y]] , \qquad
  \hat{A}^{(x,y)} \;=\; k[[x,y]] ,
\]
et $A$ s'envoie dans les deux premières, qui s'envoient toutes deux dans la
troisième. En termes de la région d'indices $(\alpha, \beta)$ des séries
$\sum a_{\alpha\beta}\, x^{\alpha} y^{\beta}$, que la marge de la page
dessine : $\hat{A}^{(x)}$ n'autorise, pour chaque $\alpha$ borné, qu'un nombre
fini de $\beta$ ; $\hat{A}^{(y)}$ l'inverse ; et $\hat{A}^{(x,y)}$ tout.
Le croquis à l'encre, deux axes $\alpha$ et $\beta$ coupés par une courbe,
figure exactement le découpage de cette région.\footnote{La ligne médiane de la
page décrit $A$ par « seul un nombre fini des $a_{\alpha\beta}$, pour
$\alpha + \beta \leqslant N$ donné, sont non nuls ». Ce n'est pas une
condition : il n'y a qu'un nombre fini de couples $(\alpha,\beta)$ avec
$\alpha+\beta \leqslant N$. La ligne enregistre l'anneau $A$ lui-même, entre
les deux complétions, et non une condition qui le caractériserait. Le feuillet
est par ailleurs traversé de traits de crayon gras et d'un griffonnage qui
recouvrent une partie de l'écriture, et l'hypothèse b) elle-même --- le mot
porté en interligne au-dessus de « complet » --- reste indéterminée.}

C'est la situation locale que la page 106 dessinera --- deux branches et le
petit disque $Z = \operatorname{Spec}\hat{\mathcal{O}}_{X,a}$ autour de leur
croisement --- et celle que les pages 87 à 97 traiteront le long d'une fibre
spéciale.

\subsection*{Invariance birationnelle, et le groupe fondamental local}

La page 42 achève la liste des conjectures ouverte page 40. Son point c) est
l'\emph{invariance birationnelle de la ramification modérée}, marqué « Prouvé »
en marge : c'est ce qui donne un sens au problème, puisqu'il permet ensuite
de procéder par normalisation et par Bertini pour se ramener à une courbe munie
d'un diviseur à croisements normaux, où l'on conclut.

Puis s'ouvre une deuxième question, locale cette fois. Soit $X$ un schéma
noethérien --- complet, ou strictement local et « bon » --- et $U$ un ouvert.
On conjecture :

\begin{enumerate}
  \item[a)] $\pi^{t}_{1}(U)$ topologiquement de type fini ;
  \item[b)] $\pi'_{1}(U)$ de présentation finie.
\end{enumerate}

Par la résolution des singularités et l'invariance birationnelle du cas a), on
doit s'en tirer en se ramenant aux énoncés globaux de la page 40. Et pour a),
par descente on se ramène au cas $X$ normal, puis par section hyperplane ---
du moins si $U = X$ privé de l'origine --- au cas $\dim X = 2$, où l'on dispose
du théorème d'Abhyankar.

En marge, deux notes qui sont des notes à soi-même : « vérifier si c'est
\emph{nécessaire} », et « cela doit marcher \dots ».

\pagerange{43}{45}
\section*{Un schéma formel modulo un sous-schéma de définition (pages 43 et 45)}

Deux feuillets, chacun barré d'un long trait de crayon, construisent deux
catégories quotient. Ils sont hors du fil des conjectures et ils annoncent les
pages 87 à 97.

Soit $\mathfrak{X}$ un préschéma formel localement noethérien et
$\mathfrak{X}_{0} = V(\mathcal{J})$ un sous-schéma de définition.

\subsection*{Les modules cohérents}

Soit $\mathrm{Coh}(\mathfrak{X})$ la catégorie des modules cohérents sur
$\mathfrak{X}$, et $\mathfrak{S}$ l'ensemble des flèches $u$ dont le noyau et
le conoyau sont tués par une puissance de $\mathcal{J}$. Cet ensemble admet un
calcul de fractions, et l'on pose
\[
  \mathrm{Coh}(\mathfrak{X}/\mathfrak{X}_{0}) \;=\;
  \mathrm{Coh}(\mathfrak{X})\bigl[\mathfrak{S}^{-1}\bigr] .
\]
En pratique --- c'est la remarque de la page --- on prend la sous-catégorie
épaisse des modules cohérents tués par une puissance de $\mathcal{J}$ et l'on
passe au quotient par elle.\footnote{C'est-à-dire, dans le vocabulaire
d'aujourd'hui, le quotient de Serre de $\mathrm{Coh}(\mathfrak{X})$ par la
sous-catégorie de Serre des modules $\mathcal{J}$-de torsion. Le folio 19 du
même fonds contient la théorie générale de ces quotients, sous le nom de
Gabriel ; ici c'est le cas particulier qui sert.}

\subsection*{L'exemple, et ce qu'il calcule}

Soit $X$ propre sur $Y = \operatorname{Spec} A$, $A$ noethérien séparé et
complet pour un idéal $I$, $Y_{0} = V(I)$, $X_{0} = X \times_{Y} Y_{0}$, et
$\mathfrak{X}$ le complété formel de $X$ le long de $X_{0}$. Le théorème
d'existence donne une équivalence
$\mathrm{Coh}(X) \xrightarrow{\ \sim\ } \mathrm{Coh}(\mathfrak{X})$, compatible
avec les sous-catégories des modules à support dans $X_{0}$ :

\begin{tikzcd}
\mathrm{Coh}(\mathfrak{X}) & \mathrm{Coh}(X) \arrow[l, "\approx"'] \\
\mathrm{Coh}(\mathfrak{X})_{\mathrm{nilp}} \arrow[u] & \mathrm{Coh}(X)_{\mathrm{nilp}} \arrow[u]
\end{tikzcd}

d'où, en passant aux quotients,
\[
  \mathrm{Coh}(\mathfrak{X}/\mathfrak{X}_{0}) \;\approx\;
  \mathrm{Coh}(X - X_{0}) .
\]
C'est le point d'exclamation de la page : ce qui reste du schéma formel une
fois qu'on a divisé par ce qui vit sur la fibre spéciale, c'est le
\emph{complémentaire} de cette fibre.

\subsection*{Les revêtements}

La même construction se fait avec les revêtements. Soit $R(\mathfrak{X})$ la
catégorie des préschémas formels finis sur $\mathfrak{X}$, et disons qu'un
$u : \mathfrak{X}' \to \mathfrak{X}''$ est un \emph{isomorphisme modulo
$\mathfrak{X}_{0}$} si l'homomorphisme d'algèbres correspondant
$\mathcal{A}'' \to \mathcal{A}'$ en est un au sens précédent. On vérifie qu'on
obtient un calcul de fractions à droite, et l'on note
$R(\mathfrak{X}/\mathfrak{X}_{0})$ la catégorie quotient.

Sur le même exemple, le théorème de comparaison donne
$R(\mathfrak{X}) \approx R(X)$, donc
\[
  R(\mathfrak{X}/\mathfrak{X}_{0}) \;\approx\; R(X/X_{0}) \;\approx\;
  R(X - X_{0}) ,
\]
avec, en marge, « à vérifier » : c'est le second isomorphisme qui n'est pas
formel, et la page indique qu'il faudrait passer par le \emph{Main Theorem}
global. Cette question --- \emph{quand un revêtement formel est-il
algébrisable} --- est exactement celle que Murre et Grothendieck traiteront en
1969 sous le nom de théorème de Mumford, aux pages 77 à 85 du même
carton.\footnote{Les deux feuillets sont barrés chacun d'un long trait de
crayon. La biffure porte sur la page entière et non sur telle ou telle ligne ;
ce que le trait signifie --- abandon, ou report ailleurs --- n'est pas dit, et
cette lecture ne le devine pas.}

\pagerange{44}{44}
\section*{Finitude en famille (page 44)}

Un troisième point de la chemise, sous le titre « Théorèmes de finitude et
simultanéité », et marqué « c'est fait » en marge. La question est de savoir
comment les groupes fondamentaux des fibres d'un morphisme de type fini
dépendent du point de la base, et l'on veut des énoncés du type
\emph{constructibilité}. Trois formes, calquées sur a), b), c) de la page 40 :

\begin{enumerate}
  \item[a)] majoration \emph{uniforme} du $\pi_{1}(X_{\bar{s}})$ dans les cas
    de finitude, pour $X/S$ propre ;
  \item[b)] majoration uniforme pour les $\pi'_{1}(X_{\bar{s}})$, ou mieux :
    \emph{la structure de $\pi'_{1}(X_{\bar{s}})$ est fonction constructible de
    $s$} ;
  \item[c)] majoration uniforme pour les $\pi^{t}_{1}(X_{\bar{s}})$.
\end{enumerate}

Et l'estimation : « il semble que les méthodes qui résoudront 1) résoudront
également 3) ». La forme b) --- la constructibilité, et non seulement une
borne --- est la plus forte des trois, et c'est celle que le lemme 1 ci-dessous
prépare.

\pagerange{46}{57}
\section*{« Théorème de finitude relatif pour $\pi_{1}$ » (pages 46 à 57)}

Une suite de six feuillets que Grothendieck pagine lui-même I à VI. Elle est
écrite vite et lourdement raturée ; les énoncés se lisent, les phrases qui les
relient souvent pas. Ce qui suit est ce que les énoncés disent.

\subsection*{Le pont entre $H^{1}$ et $\pi_{1}$}

Il n'est pas écrit sur les pages et il faut l'avoir en tête pour les lire. Pour
$X$ connexe pointé en $x$ et $G$ un groupe fini, l'ensemble $H^{1}(X,G)$ des
classes de revêtements principaux de groupe $G$ s'identifie à
\[
  \operatorname{Hom}_{\mathrm{cont}}\bigl(\pi_{1}(X,x),\, G\bigr) \big/
  \operatorname{int}(G) .
\]
Or un groupe profini est topologiquement de type fini si et seulement s'il n'a,
pour chaque entier $n$, qu'un nombre fini de sous-groupes ouverts d'indice $n$.
Donc : \emph{$\pi_{1}(X,x)$ est topologiquement de type fini si et seulement si
$H^{1}(X,G)$ est fini pour tout groupe fini $G$}. C'est pourquoi la suite,
intitulée « théorème de finitude pour $\pi_{1}$ », ne parle que de $H^{1}$.

\subsection*{Lemme 1 : un isomorphisme générique géométrique descend}

Soit $f : X \to Y$ un morphisme de préschémas localement noethériens, $Y$
intègre de point générique $y$, et soient $X'$, $X''$ deux schémas sur $X$ tels
que $X'_{\bar{y}}$ et $X''_{\bar{y}}$ soient isomorphes au-dessus de
$k(\bar{y})$. Alors il existe un revêtement fini étale $Y_{1}$ de $Y$ tel que
$X' \times_{Y} Y_{1}$ et $X'' \times_{Y} Y_{1}$ soient isomorphes au-dessus de
$X_{Y_{1}}$.

Plus précisément : si $\bar{u} : X'_{\bar{y}} \to X''_{\bar{y}}$ est donné et
si $K_{1}$ est la plus petite sous-extension de $k(\bar{y})$ sur laquelle il
est défini, alors $K_{1}$ est \emph{étale} sur $Y$, et sur le $Y_{1}$
correspondant $\bar{u}$ provient d'un morphisme
$X' \times_{Y} Y_{1} \to X'' \times_{Y} Y_{1}$.

La démonstration introduit $\underline{\mathrm{Hom}}_{X}(X', X'')$ --- ou son
sous-schéma des isomorphismes --- pour se ramener au cas $X' = X$, puis invoque
SGA IX 6.2 pour identifier la fibre en $y$ et conclure par un argument de
profondeur.\footnote{Le corps du raisonnement, aux pages 47 et 49, est presque
entièrement illisible à la définition du fac-similé et la transcription le dit.
Ce qui est restitué ici est ce que les phrases portantes autorisent : le
schéma des morphismes, la réduction à $X' = X$, le renvoi à SGA IX 6.2, et
l'isomorphisme final $X_{1} \approx Z$. Aucune étape n'est reconstruite.}

Deux notes accompagnent le lemme, et elles valent l'énoncé. La première rappelle
la suite exacte d'homotopie
\[
  \pi_{1}(X_{\bar{y}}, \xi) \longrightarrow \pi_{1}(X, \xi) \longrightarrow
  \pi_{1}(Y, \zeta) \longrightarrow 0 ,
\]
que le lemme généralise. La seconde avertit : « ne pas confondre avec les
résultats de SGA X, où on prend une fibre générique \emph{vraie} » --- et où il
faut alors supposer $f$ propre. C'est bien la fibre générique
\emph{géométrique} qui intervient ici, et c'est ce qui rend l'énoncé utile en
famille : il dit que l'isomorphisme, qui semblait n'exister qu'après passage à
une clôture algébrique, est déjà défini sur une extension finie séparable.

\subsection*{Lemme 2 : constructibilité de $f_{*}$}

Soit $f : X \to Y$ de type fini, $Y$ localement noethérien et universellement
japonais. Alors, pour tout faisceau d'ensembles constructible $F$ sur $X$,
$f_{*}(F)$ est constructible.

La marge dit « standard », et le lemme n'est pas démontré.\footnote{L'énoncé
moderne correspondant --- la constructibilité de $R f_{*}$ pour $f$ de type
fini --- se démontre sous l'hypothèse que les schémas soient \emph{excellents}
(le théorème de finitude de Deligne, SGA 4$\frac{1}{2}$). « Universellement
japonais » est la condition de Nagata, plus faible : elle ne porte que sur la
finitude des clôtures intégrales. La lecture retient l'hypothèse de la page et
signale l'écart plutôt que de la corriger, parce que l'usage qui en est fait
ci-dessous --- le lemme 3 --- ne demande que le cas noethérien.}

\subsection*{Lemme 3 : $H^{1}$ le long d'une partition}

Soit $X$ un schéma connexe noethérien, $(X_{i})$ une partition finie de $X$ en
parties localement fermées, et $G$ un groupe fini. Alors l'application
\[
  H^{1}(X, G) \longrightarrow \prod_{i} H^{1}(X_{i}, G)
\]
est à \emph{fibres finies}. En particulier, si $H^{1}(X_{i}, G)$ est fini pour
tout $i$, il en est de même de $H^{1}(X, G)$ --- « kif kif », dit la
page.\footnote{Les deux mots sont bien ceux de la page. Ils reviennent au
folio 152 du même carton (« Kif -- kif », pour dire que deux conditions d'une
proposition disent la même chose) et à la page 190 (« grâce à R1 c'est Kif
Kif »). C'est une locution de travail, non une notation.}

La démonstration passe par la théorie du recollement des faisceaux étales, en
utilisant le lemme 2, et par le fait suivant : si $F$ et $G$ sont des faisceaux
étales constructibles sur $X$ noethérien, alors $\operatorname{Hom}(F, G)$ est
un ensemble fini.

\subsection*{La proposition principale}

Soit $S$ un schéma noethérien universellement japonais et $G$ un groupe fini.
Quatre conditions sont mises en regard, dont la page ne laisse lire que la
forme.\footnote{Les conditions (i) à (iv) des pages 53 et 55 sont écrites en
abrégé et lourdement raturées ; la transcription en laisse une part importante
illisible, et notamment les quantificateurs de (ii) et de (iii) et une partie
de (iv). Ce qui suit est la forme des quatre conditions, non leur énoncé
complet ; le lecteur qui en a besoin doit aller à la transcription et, de là, à
la page.}

\begin{enumerate}
  \item[(i)] Pour tout $X$ de type fini sur $S$, $H^{1}(X, G)$ est fini.
  \item[(ii)] La même chose, restreinte aux $X$ affines lisses de type fini
    au-dessus d'un ouvert dense $S'$ d'un sous-schéma affine intègre non vide
    $S_{0}$ de $S$, à fibres géométriques connexes non vides.
  \item[(iii)] La même chose affaiblie : pour de tels $S'$, $X$, et tout ouvert
    non vide $U$ de $X$, il existe un ouvert non vide $U'$ de $U$ tel que
    $H^{1}(U', G)$ soit fini.
  \item[(iv)] Une condition en deux moitiés : a) une condition fibre à fibre,
    pour $s \in S$ et $\bar{X}$ connexe de type fini sur $s$ ; b) une condition
    sur la base seule, la finitude de $H^{1}(S', H)$ pour $S'$ comme en (ii) et
    $H$ un groupe auxiliaire. Lorsque $H$ est commutatif on peut affaiblir b)
    en y remplaçant $H$ par $G$.
\end{enumerate}

De la démonstration, une seule phrase se lit, et elle suffit à répartir le
travail : \emph{(ii) et (iii) résultent aussitôt du lemme 3 ; c'est pour (iv)
qu'il faut le lemme 1}. Le lemme 3 fait le découpage, le lemme 1 fait le
passage de la fibre géométrique générique à la base.

\subsection*{Les deux corollaires}

\textbf{Corollaire 1.} Supposons $G$ d'ordre premier aux caractéristiques
résiduelles. Supposons que pour tout sous-schéma affine intègre non vide
$S_{0}$ de $S$, tout ouvert dense $S'$ de $S_{0}$ et tout sous-groupe $H$ de
$G \times G$, l'ensemble $H^{1}(S', H)$ soit fini. Alors $H^{1}(X, G)$ est fini
pour tout $X$ de type fini sur $S$.

C'est le passage par $G \times G$ qui mérite l'attention : le lemme 1
\emph{compare deux} revêtements, et les couples qu'il fait intervenir ont pour
groupe structural un sous-groupe de $G \times G$. La condition sur la base est
donc plus forte que celle qu'on croirait avoir à imposer.

\textbf{Corollaire 2.} Deux conditions équivalentes : la finitude de
$H^{1}(X, G)$ pour tout $X$ de type fini et tout $G$ premier aux
caractéristiques résiduelles, et la même chose restreinte comme en
(ii).\footnote{L'énoncé du corollaire 2 est en grande partie illisible sur la
page 55, de même que la note qui le suit ; seule la structure --- deux
conditions dont la seconde est la restriction de la première --- se laisse
lire.}

\subsection*{La récurrence sur la dimension}

La page 57 esquisse comment ce qui précède se démontrerait. Récurrence sur
$\dim X$. Par le lemme 3, on se ramène au cas connexe. Si $\dim X = 1$, $X$ est
lisse, et c'est facile : réduction au cas d'Abhyankar, avec un lemme de
Serre--Lang. Si $\dim X \geqslant 2$, on construit une fibration \emph{lisse}
$X \to Y$ avec $\dim Y < \dim X$, et l'on conclut par l'hypothèse de
récurrence.\footnote{C'est la fin de la suite paginée I à VI, et la page 57 est la plus abîmée de la suite ; la marge porte « finir si le \dots », et la remarque
finale n'est pas lisible. La suite reprend aux pages 61 à 63, avec les
deux corollaires sur les produits.}

\pagerange{59}{60}
\section*{« $\pi_{1}$ et $\mathrm{Pic}^{\tau}$ birationnels » : la première page (pages 59 et 60)}

Une chemise nouvelle s'ouvre, dont la suite est aux pages 61 à 76. Sa première
page --- Grothendieck la numérote 1 --- porte un théorème et sa démonstration
par récurrence sur la dimension.

\textbf{Théorème.} Soit $X$ un préschéma algébrique sur un corps $k$
algébriquement clos. Alors $\pi'_{1}(X)$, limite projective des quotients
discrets de $\pi_{1}(X)$ d'ordre premier à $p$, est topologiquement de type
fini.

La démonstration procède par réductions successives.

Grâce au théorème de van Kampen on peut supposer $X$ affine, puis affine et
\emph{normal} : sinon, on considère $X' \to X$ un revêtement convenable et le
diagramme $X' \times_{X} X' \rightrightarrows X' \to X$, et $\pi_{1}(X)$ est
engendré par $\pi'_{1}(X')$ et par un nombre fini de générateurs correspondant
aux composantes irréductibles de $X' \times_{X} X'$.

Si $\dim X = 0$, c'est trivial. Si $\dim X = 1$, la théorie de la
spécialisation du groupe fondamental ramène à la caractéristique zéro, où l'on
conclut par voie transcendante.

Si $\dim X \geqslant 2$, on considère un revêtement
$\bar{X}^{n} \to \mathbf{P}^{r}_{k}$ et un sous-espace linéaire $L^{r-n+1}$
défini sur une extension algébriquement close $K$ de $k$, et l'on pose
\[
  C \;=\; \bar{X}^{n} \times_{\mathbf{P}^{r}_{k}} L^{r-n+1} ,
\]
courbe normale sur $K$, « générique ». On montre alors que
\[
  \pi_{1}(C) \longrightarrow \pi_{1}(X)
\]
est \emph{surjectif}, d'où la surjectivité de $\pi'_{1}(C) \to \pi'_{1}(X)$, ce
qui conclut puisque le cas des courbes est acquis.

C'est le théorème de Bertini de Zariski qui fournit cette surjectivité, et il
est énoncé et démontré à la page 66.

\subsection*{La question, et un mot anglais}

La page se termine sur une question, qui est celle vers laquelle tout le
dossier des finitudes tend. Soit $X$ normale complète et $Y \subset X$ fermé.
Le quotient de $\pi_{1}(X - Y)$ par le sous-groupe engendré par les groupes
d'inertie \emph{sauvages} aux composantes de codimension 1 de $Y$ est-il
topologiquement de type fini ?\footnote{Le mot \emph{wild} est en anglais sur
la page, et il y reste : c'est la partie $p$-primaire de l'inertie, celle
qu'aucune extension modérée ne contrôle, et le carton n'a pas de mot français
pour elle. C'est exactement ce que les énoncés a), b), c) de la page 40
cherchaient à écarter en passant à $\pi'_{1}$ ou à $\pi^{t}_{1}$ ; la question
demande si on peut le faire sans jeter le reste.}

\pagerange{61}{64}
\section*{Fin du théorème de finitude relatif : deux corollaires sur les produits (pages 61 à 64)}

Les pages 61 à 63 portent, dans l'angle supérieur gauche, les numéros 2, 3 et 4
que Grothendieck donne à la suite ouverte page 46. La page 61 est un
palimpseste --- un premier jet au crayon repris à l'encre, puis biffé --- et sa
plus grande partie est perdue ; cette lecture n'y supplée pas. Ce qui en reste
lisible est la réduction à une courbe : $X'$ le normalisé de $X_{\mathrm{red}}$,
$D_{i}$ les composantes irréductibles de codimension 1 de l'image inverse d'un
diviseur $D$, $U = X - D$, et le passage du triplet $(k, X, U)$ au triplet
$(K, C, V)$ où $C$ est une \emph{courbe}.

Deux corollaires en sortent, et ils sont propres.

\textbf{Corollaire (produits de paires).} Soient $X'$, $X''$ propres sur $k$,
et $Z' \subset X'$, $Z'' \subset X''$ des parties fermées. Posons
\[
  X = X' \times X'' , \qquad
  Z = (Z' \times X'') \cup (X' \times Z'') .
\]
Alors, pour des points-base compatibles,
\[
  \pi^{t}_{1}(X ; Z) \;\xrightarrow{\ \sim\ }\;
  \pi^{t}_{1}(X' ; Z') \times \pi^{t}_{1}(X'' ; Z'')
\]
est un \emph{isomorphisme}. Autrement dit : le groupe fondamental modéré d'un
produit de paires est le produit des groupes fondamentaux modérés.

\textbf{Corollaire (produits, partie première à $p$).} Soient $X$, $Y$ de type
fini sur $k$. Alors
\[
  \pi'_{1}(X \times Y) \;\xrightarrow{\ \sim\ }\;
  \pi'_{1}(X) \times \pi'_{1}(Y) .
\]

Cette dernière est l'énoncé que la lettre à Serre de 1959 conjecturait, à la
page 128 du même carton, comme conséquence de sa suite exacte d'homotopie ---
« elle impliquerait que pour des schémas complets $X$, $Y$ au-dessus d'un
corps algébriquement clos $k$, on a
$\pi_{1}(X\times Y/k) = \pi_{1}(X/k)\times\pi_{1}(Y/k)$ ». Dix ans, et une
centaine de pages du carton, séparent la conjecture du
corollaire.\footnote{À une différence près, qui est celle de tout le carton :
la lettre de 1959 énonce le produit pour le groupe fondamental \emph{total},
partie infinitésimale comprise --- un pro-schéma en groupes --- et la page 62
l'énonce pour le quotient premier à la caractéristique. Ce n'est pas le même
énoncé, et c'est la restriction que la lettre de 1959 (pages 124 à 129) rend intelligible.}

La restriction à $\pi'_{1}$ est signalée sur la page même : « mais elles \dots\
la car. ». C'est l'aveu que le $p$-groupe fondamental d'un produit n'est pas le
produit des $p$-groupes fondamentaux --- ce qui est exact, et c'est l'un des
phénomènes propres à la caractéristique $p$.

La page 63 énonce une variante de la question. Soit $X'$ un normalisé de $X$,
$Y \subset X$ fermé, $U = X - Y$ ; soient $Y_{1}, \dots, Y_{r}$ les composantes
de $Y$ de codimension 1 dans $X$ et $Y'_{ij}$ leurs images inverses dans $X'$.
Considérons les revêtements étales de $X - Y$ qui deviennent étales au-dessus
des anneaux locaux des $Y'_{ij}$ : ils sont classifiés par un groupe de type
galoisien $\pi$.

\textbf{Théorème.} Le groupe $\pi$ est topologiquement de génération finie.

Deux corollaires : pour $X$ complète et $Y$ de codimension $\geqslant 2$, un
quotient convenable de $\pi_{1}(X - Y)$ est topologiquement de génération
finie ; et pour $X$ algébrique, $\pi'_{1}(X)$ est topologiquement de type
fini.\footnote{Le corollaire 1 porte, à l'endroit du quotient, un mot que la
transcription laisse illisible. Cette lecture ne le devine pas : le groupe
quotient n'est pas nommé, et l'énoncé reste incomplet tel que la page le
laisse.}

La page 64 est un brouillon : des carrés de schémas, la surjectivité de
$\pi_{1}(X\times Y)\to\pi_{1}(X)\times\pi_{1}(Y)$, et en marge « surj ?? » en
regard de $\pi_{1}(Y)\to\pi_{1}(X)$ pour $X$ normal. Elle n'est pas reprise
ici : rien n'y est énoncé.

\pagerange{65}{69}
\section*{« $\pi_{1}$ birationnel » (pages 65 à 69)}

Grothendieck pagine cette suite A à E, et la suivante F à L.

\subsection*{Le théorème}

\textbf{Théorème.} Soient $k$ algébriquement clos et $X$ un schéma connexe sur
$k$ dont les composantes irréductibles $X_{i}$ soient isomorphes --- avec leur
structure réduite induite --- au complémentaire d'une partie fermée de
codimension $\geqslant 2$ dans un schéma propre sur $k$. Alors $\pi_{1}(X)$ est
topologiquement de type fini.

\emph{Démonstration.} Par SGA IX 6.2 on se ramène au cas où $X$ est intègre,
donc de la forme $X = Z - Y$ avec $Z$ propre et $Y$ de codimension $\geqslant 2$
dans $Z$ ; on peut supposer $Z$ intègre. Si $Z' \to Z$ est propre surjectif avec
$Z'$ intègre, et $X'$ l'image inverse de $X$, il suffit de prouver l'énoncé pour
$X'$ ; par le lemme de Chow on peut donc supposer $Z$ \emph{projectif}.

Récurrence sur $\dim Z$. Si $\dim Z \leqslant 0$, trivial. Si $\dim Z = 1$,
alors $Y = \varnothing$ puisque $Y$ est de codimension $\geqslant 2$, donc
$X = Z$ et le théorème est connu. Si $\dim Z \geqslant 2$, on utilise le lemme
suivant.

\subsection*{Le lemme de Bertini de Zariski}

\textbf{Lemme.} Soient $k$ un corps, $f : X \to \mathbf{P}^{r}_{k}$ un
morphisme avec $X$ \emph{irréductible} et $\dim \overline{f(X)} \geqslant n$.
Soit $1 \leqslant s \leqslant n-1$, soit $u$ le point générique d'un produit de
$s$ facteurs paramétrant les hyperplans, $K$ une clôture algébrique de $k(u)$,
$L_{u}$ la variété linéaire de $\mathbf{P}^{r}_{K}$ définie par $u$, et
\[
  X_{u} \;=\; X \times_{\mathbf{P}^{r}_{k}} L_{u}
        \;=\; X_{K} \times_{\mathbf{P}^{r}_{K}} L_{u} .
\]
Alors $X_{u}$ est irréductible.

En marge, la question et sa réponse, de sa main : « peut-on remplacer
\emph{irréductible} par \emph{connexe} ? --- non ». C'est un avertissement, non
une remarque : c'est l'irréductibilité, non la connexité, qui se transmet à une
section générique, et l'argument suivant repose entièrement sur cette
différence.

\textbf{Corollaire.} Si $X$ est \emph{normal}, alors
$\pi_{1}(X_{u}) \to \pi_{1}(X)$ est \emph{surjectif}.

\emph{Démonstration.} Soit $X' \to X$ un revêtement étale connexe. Puisque $X$
est normal, $X'$ est normal, donc connexe implique irréductible ; par le lemme,
$X'_{u} = X' \times_{X} X_{u}$ est irréductible, donc connexe. Un revêtement
étale connexe de $X$ reste donc connexe au-dessus de la section générique, ce
qui est exactement la surjectivité annoncée.

\begin{tikzcd}[column sep=small, row sep=small, nodes={font=\scriptsize}]
X'_u \arrow[r] \arrow[d] & X' \arrow[d] \\
X_u \arrow[r] \arrow[d] & X_K \arrow[r] \arrow[d] & X \arrow[d] \\
L_u \arrow[r] & \mathbf{P}^r_K \arrow[r] & \mathbf{P}^r_k
\end{tikzcd}

\subsection*{La conclusion : $2g$ générateurs}

Reprenons $Z$ normal de dimension $n \geqslant 2$ et $X = Z - Y$. Plongeons $Z$
dans $\mathbf{P}^{r}$, et prenons $s = n-1$, de sorte que $L_{u}$ soit de
codimension $n-1$. Alors
\[
  X_{u} \;=\; X_{K} \cap L_{u} \;=\; Z_{K} \cap L_{u} \;\neq\; \varnothing ,
\]
la seconde égalité parce que $Y_{K} \cap L_{u} = \varnothing$ --- $Y$ étant de
dimension $\leqslant n-2$ et $L_{u}$ générique de codimension $n-1$. Donc
$X_{u}$ est une courbe projective, d'ailleurs normale. Si $g$ en est le genre,
le corollaire au lemme de Bertini donne : \emph{$\pi_{1}(X)$ admet $2g$
générateurs topologiques}.\footnote{Le groupe fondamental d'une courbe
projective lisse de genre $g$ sur un corps algébriquement clos est
topologiquement engendré par $2g$ éléments, en toute caractéristique ---
en caractéristique $p$ c'est un théorème, obtenu par spécialisation depuis la
caractéristique zéro, et il n'y a pas de présentation finie du groupe entier :
les relations manquent. C'est exactement la différence entre l'énoncé a) et
l'énoncé b) de la page 40 (« de type fini » contre « de présentation finie »).}

\subsection*{Le passage à la limite sur les modèles}

Soit $K$ une extension de type fini de $k$. Les modèles propres et normaux de
$K$, ordonnés par domination, forment un ensemble ordonné filtrant croissant.

Soit $X$ un tel modèle et $U \subset X$ un ouvert. Deux faits, à distinguer
soigneusement :

\begin{enumerate}
  \item si $X$ est \emph{normal} et $U$ est un ouvert dense, l'homomorphisme
    $\pi_{1}(U) \to \pi_{1}(X)$ est \emph{surjectif} --- un revêtement étale
    connexe de $X$ reste connexe au-dessus de $U$ ;
  \item si $X$ est \emph{régulier} et $X - U$ de codimension $\geqslant 2$,
    l'homomorphisme est un \emph{isomorphisme} : c'est le \emph{théorème de
    pureté}.
\end{enumerate}

C'est cette distinction que le diagramme de la page 68 porte, l'isomorphisme du
bas étant annoté « th.\ de pureté ».\footnote{La phrase de la page, elle,
attribue la surjectivité à la pureté : « L'injection induite $U \to X$ étant le
complémentaire d'une partie fermée de codim.\ $\geqslant 1$, on voit que (par
pureté) on a un homo.\ \emph{surjectif} (si $X$ normal) ». En codimension 1 la
pureté ne dit rien, et la normalité suffit ; en codimension $\geqslant 2$ la
pureté donne l'isomorphisme, et c'est bien là que le diagramme la place. La
présente lecture sépare les deux énoncés.}

\begin{tikzcd}
\pi_1(X_1) \arrow[d, "f_{*}"'] & \pi_1(U_1) \arrow[l] & \\
\pi_1(X_0) & \pi_1(U_0) \arrow[l] \arrow[r, "\sim"] \arrow[ul, "g_{*}"] & \pi_1(U_1)
\end{tikzcd}

À la limite, on obtient un homomorphisme \emph{surjectif}
\[
  \pi_{1}(U, \Omega) \longrightarrow
  \varprojlim_{X \text{ modèle propre}} \pi_{1}(X, \Omega) ,
\]
$\Omega$ étant une extension algébriquement close de $K$. Le terme de droite est
ce que Grothendieck appelle le \emph{groupe fondamental birationnel} de $K/k$,
et qu'il note $\pi_{1}(K/k, \Omega)$.

\textbf{Théorème.} Si $K/k$ est une extension de type fini avec $k$
algébriquement clos, alors $\pi_{1}(K/k, \Omega)$ est topologiquement de type
fini.\footnote{Le mot \emph{birationnel} est de lui : « Le dernier est ce que
j'ai appelé le \emph{groupe fondamental birationnel} de $K/k$ ». Le passage à
la limite sur les modèles propres normaux d'un corps de fonctions est
aujourd'hui un objet standard, et la finitude est ce que la borne $2g$ ci-dessus
donne : le genre de la courbe générique ne dépend pas du modèle choisi, et
c'est le même entier qui majorera le rang de $\mathrm{Pic}$ à la page 76.}

\pagerange{70}{76}
\section*{« Pic birationnel » (pages 70 à 76)}

\subsection*{Le théorème}

\textbf{Théorème.} Soit $f : X \to Y$ un morphisme dominant \emph{birationnel}
de préschémas intègres normaux, propres sur $k$, avec $Y$ localement régulier.
Alors
\[
  f^{*} : \underline{\mathrm{Pic}}^{\tau}_{Y} \longrightarrow
  \underline{\mathrm{Pic}}^{\tau}_{X}
\]
est un \emph{isomorphisme}.

\emph{C'est un monomorphisme} parce que $f_{*}(\mathcal{O}_{X}) =
\mathcal{O}_{Y}$ universellement --- ce qui est le théorème principal de
Zariski, $f$ étant propre birationnel et $Y$ normal --- donc une immersion
fermée, s'agissant d'un homomorphisme de schémas en groupes propres.

\emph{Montrons qu'il est surjectif.} Soit $U$ l'ouvert de $Y$ où $f^{-1}$ est
défini ; son complémentaire est de codimension $\geqslant 2$. Comme $Y$ est
localement régulier, donc factoriel, la restriction
\[
  i^{*} : \mathrm{Pic}(Y) \xrightarrow{\ \sim\ } \mathrm{Pic}(U)
\]
est un isomorphisme.\footnote{La page écrit ici
« $\mathrm{Pic}(X) \xrightarrow{\ \sim\ } \mathrm{Pic}(U)$ », alors que sa
propre raison --- « comme $Y$ est factoriel » --- et son propre diagramme,
qui porte $i^{*} : \mathrm{Pic}(Y) \to \mathrm{Pic}(U)$, demandent l'un et
l'autre $\mathrm{Pic}(Y)$. C'est un lapsus, et l'énoncé écrit ici est celui que
l'argument utilise ; sans lui la factorialité de $Y$ ne servirait à rien.}
Soit $g : U \to X$ la section de $f$ sur $U$ et $i : U \to Y$ l'inclusion, de
sorte que $i = f \circ g$ et $i^{*} = g^{*} \circ f^{*}$ :

\begin{tikzcd}
& X \arrow[d, "f"] \\
U \arrow[ur, "g"] \arrow[r, "i"'] & Y
\end{tikzcd}

\begin{tikzcd}
& \mathrm{Pic}(X) \arrow[dl, "g^{*}"] \\
\mathrm{Pic}(U) & \mathrm{Pic}(Y) \arrow[u, "f^{*}"'] \arrow[l, "i^{*}"]
\end{tikzcd}

Puisque $i^{*}$ est un isomorphisme, $f^{*}$ est injectif et admet
$(i^{*})^{-1} g^{*}$ pour rétraction : $\mathrm{Pic}(Y)$ est un
\emph{facteur direct} de $\mathrm{Pic}(X)$, d'image l'ensemble des $\xi$ tels
que $\xi = f^{*}(i^{*})^{-1}g^{*}(\xi)$.

Le supplémentaire est engendré par les diviseurs $\sum n_{i} X_{i}$ portés par
les composantes exceptionnelles de $f$. Il reste donc à voir qu'une telle
combinaison $\tau$-équivalente à zéro est nulle --- « c'est là un lemme bien
connu », et c'est en effet la non-dégénérescence de la forme d'intersection sur
les diviseurs exceptionnels. Ainsi $f^{*}$ est une immersion \emph{surjective},
et il ne reste plus qu'à voir qu'elle induit un isomorphisme sur les espaces
tangents à l'origine, c'est-à-dire le

\textbf{Corollaire.} Sous les conditions du théorème,
\[
  f^{*} : H^{1}(Y, \mathcal{O}_{Y}) \longrightarrow H^{1}(X, \mathcal{O}_{X})
\]
est un isomorphisme.

\subsection*{Le corollaire en caractéristique $p$}

L'injectivité est acquise ; c'est la surjectivité qu'il faut. En caractéristique
$p > 0$, il suffit de prouver que sont bijectifs les deux homomorphismes
\[
  H^{1}(Y, \mathcal{O}_{Y})^{F} \longrightarrow H^{1}(X, \mathcal{O}_{X})^{F} ,
  \qquad
  {}_{F}H^{1}(Y, \mathcal{O}_{Y}) \longrightarrow
  {}_{F}H^{1}(X, \mathcal{O}_{X}) ,
\]
où $F$ est le Frobenius : les invariants d'une part, la partie tuée par $F$
d'autre part.

Le premier s'écrit, par la suite exacte d'Artin--Schreier
$0 \to \mathbb{Z}/p \to \mathcal{O} \xrightarrow{F-1} \mathcal{O} \to 0$,
\[
  H^{1}(Y, \mathbb{Z}/p\mathbb{Z}) \longrightarrow
  H^{1}(X, \mathbb{Z}/p\mathbb{Z}) ,
\]
et c'est un isomorphisme parce que $\pi_{1}(X) \to \pi_{1}(Y)$ en est un, par
le théorème de pureté sur $Y$ régulier.\footnote{L'identification demande
$X$ et $Y$ propres et connexes sur un corps algébriquement clos : la suite
d'Artin--Schreier donne
$0 \to H^{0}(\mathcal{O})/(F-1) \to H^{1}(\mathbb{Z}/p) \to
H^{1}(\mathcal{O})^{F} \to 0$, et le terme de gauche s'annule puisque
$H^{0}(\mathcal{O}) = k$ et que $F-1$ est surjectif sur un corps
algébriquement clos. La page écrit l'identification sans l'hypothèse ; elle est
supplée ici.}

Le second s'écrit, $X$ et $Y$ étant réduits,
\[
  H^{0}\bigl(Y, \mathcal{O}_{Y}/\mathcal{O}_{Y}^{p^{\nu}}\bigr) \longrightarrow
  H^{0}\bigl(X, \mathcal{O}_{X}/\mathcal{O}_{X}^{p^{\nu}}\bigr) ,
\]
ou encore, $k$ étant algébriquement clos donc parfait,
\[
  f^{*} : H^{0}\bigl(Y, \mathcal{O}_{Y}^{p^{-\nu}}/\mathcal{O}_{Y}\bigr)
  \longrightarrow
  H^{0}\bigl(X, \mathcal{O}_{X}^{p^{-\nu}}/\mathcal{O}_{X}\bigr) .
\]
Il est bijectif grâce à la \emph{normalité} de $X$ et de $Y$ : avec les
notations ci-dessus, on a le même triangle, où $g^{*}$ est injectif parce que
$\mathcal{O}_{X}^{p^{-\nu}}/\mathcal{O}_{X}$ est sans torsion ; donc $f^{*}$ et
$(i^{*})^{-1}g^{*}$ sont réciproques l'un de l'autre.

\begin{tikzcd}
& T(X) \arrow[dl, "g^{*}"] \\
T(U) & T(Y) \arrow[u, "f^{*}"'] \arrow[l, "i^{*}"]
\end{tikzcd}

C'est le seul endroit du carton où le découpage par le Frobenius soit écrit
pour lui-même, et il vaut d'être retenu comme méthode : \emph{la partie de
$H^{1}(\mathcal{O})$ fixée par le Frobenius est du groupe fondamental, la
partie tuée par lui est de la normalité}.

\subsection*{Le corollaire en caractéristique zéro}

On peut supposer que $k$ est le corps des fonctions d'un schéma intègre $S$ de
type fini sur $\mathbb{Z}$, et que $f : X \to Y$ provient d'un morphisme de
$S$-préschémas $\bar{f} : \bar{X} \to \bar{Y}$, intègres, $\bar{Y}$ lisse sur
$S$, $\bar{X}$ normal sur $S$, à fibres connexes. On peut alors supposer $X$ et
$Y$ plats sur $S$, et l'on a l'injectivité de
$R^{1}q_{*}(\mathcal{O}_{\bar{Y}}) \to R^{1}p_{*}(\mathcal{O}_{\bar{X}})$ ; la
constance de la dimension sur $S$, jointe au cas de caractéristique $> 0$
appliqué aux fibres, donne le résultat.\footnote{Une part importante des pages
73 et 74 est illisible et la transcription le dit. La structure de l'argument
--- étalement sur un schéma de type fini sur $\mathbb{Z}$, puis descente du cas
$p > 0$ aux fibres --- est ce que les phrases portantes autorisent à
restituer ; le détail du contrôle de la dimension ne l'est pas, et il n'est pas
inventé ici.}

\subsection*{Les deux corollaires et la question}

\textbf{Corollaire 1.} Invariance birationnelle de
$\underline{\mathrm{Pic}}^{\tau}_{X/k}$ pour $X$ propre et lisse sur $k$.

\textbf{Corollaire 2.} Soit $K$ une extension séparable primaire de type fini
de $k$ admettant un modèle propre et \emph{lisse} sur $k$. Alors
$\varinjlim_{X} \underline{\mathrm{Pic}}^{\tau}_{X/k}$, $X$ parcourant les
modèles propres et normaux de $K$ sur $k$, existe --- le système inductif étant
\emph{essentiellement constant}.

\textbf{Question.} Le corollaire 2 reste-t-il valable en supposant seulement
que $K$ admette un modèle \emph{normal} ? Cela équivaudrait à trois énoncés de
bornitude, sur tous les modèles propres normaux :
\begin{enumerate}
  \item $\dim \underline{\mathrm{Pic}}_{X/k}$ reste majoré ;
  \item l'ordre de la torsion de $\mathrm{Pic}(X/k)$ reste majoré ;
  \item $\dim H^{1}(X, \mathcal{O}_{X})$ reste majoré.
\end{enumerate}
Avec, en marge, l'observation que (i) résulte de (iii).

\subsection*{La borne uniforme, et la question finale}

\textbf{Théorème.} Soit $K$ une extension de type fini du corps $k$. Il existe
un entier $n(K/k)$ tel que, pour tout modèle $X$ de $K$ propre et normal sur
$k$,
\[
  \dim \underline{\mathrm{Pic}}_{X/k} \;\leqslant\; n(K/k) .
\]

\emph{Démonstration.} On se ramène au cas où $k$ est algébriquement clos, en
posant $\bar{K} = K \otimes_{k} \bar{k}$, extension de type fini de $\bar{k}$.
On utilise alors
\[
  \dim \underline{\mathrm{Pic}}_{X/k} \;=\; \tfrac{1}{2}\,
  \operatorname{rang}_{\mathbb{Z}_{\ell}}
  \bigl(\text{partie $\ell$-primaire de } \pi_{1}(X/k) \text{ rendu
  abélien}\bigr) ,
\]
et l'on conclut par le théorème de majoration uniforme du $\pi_{1}$ pour les
modèles normaux --- lequel résulte de la partie facile de Lefschetz, prise sur
les groupes fondamentaux rendus abéliens.\footnote{L'égalité est celle du
module de Tate : pour $X$ propre et normal sur $k$ algébriquement clos, la
partie libre du complété $\ell$-adique de $\pi_{1}(X)^{\mathrm{ab}}$ est
$T_{\ell}\bigl(\underline{\mathrm{Pic}}^{0}_{X/k}\bigr)$, de rang
$2\dim\underline{\mathrm{Pic}}^{0}$ ; et $\underline{\mathrm{Pic}}^{\tau}$ et
$\underline{\mathrm{Pic}}^{0}$ ont même dimension, leur quotient étant fini.
La page énonce l'égalité sans ces précisions ; c'est bien le rang de la partie
libre qui intervient, non celui du groupe entier.}

De façon plus précise : prenant un modèle $X$ projectif propre et normal, on
prend dans $X$ une « courbe générique $C$ » --- celle du lemme de Bertini
ci-dessus --- et pour $n(K)$ le \emph{genre} de cette courbe générique. C'est
donc, à un facteur près, la même borne que celle du groupe fondamental
birationnel : les deux moitiés de la chemise donnent le même entier.

\begin{tikzcd}
C' \arrow[r, hook] \arrow[d] & X' \arrow[d] \\
C_{k(u)} \arrow[r, hook] & X \arrow[r, hook] & \mathbf{P}^r
\end{tikzcd}

La chemise s'achève sur une question, et sur trois points d'interrogation en
marge : \emph{y a-t-il une démonstration du fait que cela donne une majoration,
n'utilisant pas la théorie du groupe fondamental et le fait que
$\underline{\mathrm{Pic}}^{\tau}_{X/k}$ est propre sur $k$ si $X$ est normale ?}

\pagerange{77}{80}
\section*{« Revêtements formels étales et théorème de Mumford » : Murre revient (pages 77 à 80)}

Un feuillet de sa main porte le titre, et sous lui vient une lettre de Murre
datée du 30 avril 1969 --- deux ans après l'échange sur la ramification
modérée. Elle accompagne la \emph{version préliminaire} des notes que Murre a
rédigées sur les revêtements formels modérément ramifiés et sur un résultat de
« plumbing », et elle demande un avis avant d'aller plus loin : « To me the
entire thing seems somewhat indigestible ». Le chapitre sur l'analogue du
théorème de Mumford reste à écrire.

Suivent deux pages de commentaires, chapitre par chapitre, dont trois méritent
d'être relevées parce qu'elles disent où la théorie fait mal.

\begin{itemize}
  \item Chapitre I : les morphismes étales de schémas formels sont définis
    comme systèmes inductifs de morphismes étales de schémas ordinaires, et le
    lemme crucial est que les voisinages étales formels sont localement des
    complétés de voisinages étales ordinaires.
  \item Chapitre III : « This is a nuisance ! » Murre n'a pas suivi le nouveau
    dispositif --- celui des données de ramification --- et il en donne la
    raison : « one gets the feeling that the notion of tamely ramified covering,
    although natural, is very unmanageable if one has to use \emph{the functor}
    instead of the geometric realization itself ». C'est la réponse de
    l'utilisateur au manifeste de la page 31, et elle mérite d'être lue à côté
    de lui : ce que Grothendieck tient pour indispensable à la compréhension,
    Murre le tient pour impraticable à la rédaction.
  \item Chapitre III encore : c'est en 3.6 que la réalisation géométrique
    \emph{détermine} le revêtement modérément ramifié, et Murre note que la
    normalité, plutôt que la régularité, suffit d'ordinaire --- avec cette
    difficulté que la normalité n'est peut-être pas stable par localisation
    étale formelle.
\end{itemize}

Le reste des commentaires est technique : quels lemmes servent où, quelle
hypothèse ($D_{i} \cap D_{j} = \varnothing$) est imposée par quel manque, et
l'aveu, pour le chapitre IX, que certains points ont été utilisés sans être
certains « because it seems to me that you used them implicitely and without
mentioning in your notes ». Les notes de lecture que Grothendieck a prises sur
ce texte sont aux pages 81, 83 et 85.

\pagerange{81}{85}
\section*{Notes de lecture, et l'appel d'un champ (pages 81, 83 et 85)}

Trois feuillets, en anglais, notent chapitre par chapitre ce que Grothendieck
veut voir changer dans le texte de Murre --- celui que la lettre du 30 avril
1969, page 78, accompagnait. Ce sont des notes de travail, faites de renvois et
de phrases coupées, et il n'y a pas lieu de les paraphraser toutes. Six items,
qui disent chacun quelque chose.

\begin{itemize}
  \item \emph{Se restreindre aux morphismes adiques.} L'ambiguïté sur ce qu'est
    un point d'un schéma formel est jugée « quite general » --- et le chapitre
    IV se voit prescrire d'\emph{introduire les points géométriques} de
    $\mathcal{X}$, plutôt que de raisonner sur
    $\operatorname{Spec}(\mathcal{O}_{\mathcal{X},s})$.
  \item \emph{Le chapitre III ne sert qu'à un énoncé.} « whole ch.\ III seems
    used only to get 3.5.8 ». Et il faudrait y ajouter deux choses : une
    interprétation topologique de $\pi_{1}^{D}$, et une section sur le cas où
    les diviseurs \emph{ont des équations} --- « This should make the situation
    much clearer ! ». C'est le cas particulier a) des pages 19 et 20, et c'est en effet
    celui où la donnée kummérienne se décrit sans détour.
  \item \emph{Le champ.} « Appl.\ descent for tamely ramified coverings ? At
    least one would like to have a \emph{champ} on
    $\mathrm{Et}(S)$ ! »\footnote{Le mot est en français dans ces notes
    anglaises, comme il l'est dans le tapuscrit de 1967 (page 22 : « qui est
    même un champ »). C'est exactement l'objet que les pages 190 à 200 de ce
    même carton définiront sous le nom de \emph{domaine de ramification} : un
    champ sur $\mathrm{Et}(S)$, muni d'une réalisation géométrique. Le carton
    ne date ni ces notes ni ces pages, et cette lecture n'en tire aucune
    chronologie ; le vœu et sa réalisation sont l'un et l'autre dans le même
    carton, à cent pages de distance.}
  \item \emph{La normalité est un artifice.} « Normality-assumption unnatural
    (formal normality, as formal schemes\dots) ». C'est le même reproche qu'il
    faisait à Murre en 1967 (page 8 : « I am convinced that these assumptions
    are anyhow artificial »), et que Murre lui retourne dans ses commentaires
    (page 79 : la normalité n'est peut-être pas stable par localisation étale).
  \item \emph{Où est le lemme d'Abhyankar ?} La question est posée deux fois,
    page 81 et en marge de la page 83 : « yoga Abhyankar lemma to be used ».
    C'est la reconnaissance que le résultat qui, dans le cas des croisements
    normaux, transforme la structure en propriété, n'a pas trouvé sa place dans
    la rédaction formelle.
  \item \emph{Le dernier item, page 85, est le plus concret}, et c'est un
    diagnostic de rédaction : ce que le chapitre X démontre longuement revient
    à montrer qu'un revêtement de $S$ est modéré relativement à $D_{1}, D_{2}$
    si et seulement si ses restrictions à $S_{1}$, $S_{2}$ le sont ---
    « \emph{should be trivial} with the conditions considered ! ».
\end{itemize}

\subsection*{Deux feuillets de son propre tapuscrit : l'extension de Raynaud}

Les pages 82 et 84 sont deux feuillets d'un tapuscrit français à lui, dans le
désordre --- ses pages 74 et 69 ---, portant des corrections de sa main. Ils
n'ont pas de rapport direct avec ce qui précède, mais l'un d'eux baptise un
objet, et le baptême est fait à la main sur la page.

La page 84 démontre que, dans le cas favorable où le schéma formel abélien
$\hat{B}$ est algébrisable et où $T_{o}$ est isotrivial, on trouve un schéma en
groupes $A^{\sharp}$ sur $S$, extension canonique
\[
  0 \longrightarrow T \longrightarrow A^{\sharp} \longrightarrow B
  \longrightarrow 0
\]
du schéma abélien $B$ par le tore $T$ qui relève $T_{o}$, caractérisée à
isomorphisme unique près par le fait de donner un système cohérent
d'isomorphismes d'extensions $A^{\sharp}_{n} \simeq A_{n}$, c'est-à-dire un
isomorphisme de schémas formels en groupes
$(A^{\sharp})^{\wedge} \simeq \hat{A}$ respectant les structures d'extension.
Elle dépend fonctoriellement de $A$ et sa formation commute au changement de
base.

La phrase centrale a été refaite à la main, et le carton conserve les deux
états : le tapuscrit annonçait « une extension canonique du schéma abélien $B$
\dots\ que nous appellerons l'\emph{extension de Raynaud} associée à $A$ », et
l'insertion manuscrite substitue « un schéma en groupes $A^{\sharp}$ sur $S$
appelé le \emph{groupe de Raynaud} associé au schéma en groupes $A$ sur
$S_{o}$ ». Le nom se déplace de l'extension à l'objet, sur la
page.\footnote{La page 82 --- sa page 74, donc postérieure dans le texte ---
porte le lemme qui sert : les extensions de $\underline{\mathbb{Z}}$ par
$G^{o\sharp}$ correspondent aux torseurs sous
$H = \underline{\mathrm{Hom}}(\underline{\mathbb{Z}}, G^{o\sharp})$, et,
$H$ étant fini localement libre sur $S$, la catégorie des torseurs sous $H$
équivaut par complétion formelle à celle des systèmes cohérents de torseurs
sous les $H_{n}$. La transcription signale que le caractère lu
$\underline{\mathbb{Z}}$ est douteux ; la lecture qui suit n'en dépend pas.}

\pagerange{86}{86}
\section*{Images inverses et voisinages infinitésimaux (page 86)}

Un feuillet d'une écriture très fine, sur un papier que traverse le tapuscrit
du verso ; la plus grande partie n'en est pas lisible et cette lecture n'y
supplée pas. Ce qui s'en dégage est une question, non un énoncé : l'image
inverse d'un module quasi-cohérent est quasi-cohérente, celle d'un module plat
est plate, et il faut vérifier que cela respecte la notion voulue ; puis, si
$X' \to X \times_{S} S'$ est non ramifié, il en va de même du morphisme induit
sur les espaces tangents, et l'on peut passer aux voisinages infinitésimaux de
tous ordres.

\begin{tikzcd}
X \times_{S} S' \arrow[r] \arrow[d] & X \arrow[d] \\
S' \arrow[r] & S
\end{tikzcd}

Le cube au crayon dont ce carré n'est qu'une face n'est pas reproduit : la
transcription en donne la raison, qui est la bonne --- le dessiner reviendrait
à affirmer ce qu'on n'a pas lu.

\pagerange{87}{89}
\section*{Le $\pi_{1}$ modéré d'un schéma formel le long de sa fibre spéciale (pages 87 et 89)}

Voici le calcul. Grothendieck pagine cette suite 1 à 6.

\subsection*{La situation}

Soit $\mathcal{X}$ un schéma formel régulier connexe, $X_{o} = \mathcal{X}_{o}$
un diviseur régulier --- la fibre spéciale ---, et
\[
  D \;=\; \bigcup_{i=0}^{n} D_{i}
\]
un diviseur à croisements normaux, les $D_{i}$ réguliers irréductibles, avec
$D_{o} = X_{o}$. On pose $D' = \sum_{i \neq 0} D_{i}$, on note $D'_{o}$ le
diviseur induit sur $X_{o}$, et $\mathcal{U} = \mathcal{X} - D_{o}$.

Il faut avoir en tête le dessin que la marge de la page porte : $D_{o} = X_{o}$
horizontal, coupé transversalement par $D_{1}$ et $D_{2}$, et le point de base
$\bar{\xi}_{o}$ sur $X_{o}$ en dehors des croisements.

\subsection*{La suite exacte}

Elle est double, et la seconde ligne dévisse le noyau de la première :
\[
  e \longrightarrow I_{o} \longrightarrow
  \pi_{1}^{t}(\mathcal{X}/D, \bar{\xi}_{1}) \longrightarrow
  \pi_{1}^{t}(X_{o}/D'_{o}, \bar{\xi}_{o}) \longrightarrow e ,
\]
\[
  \mu^{\infty}_{X_{o},\,\mathrm{mod}}(\bar{\xi}_{o}) \longrightarrow
  I_{o} \longrightarrow 0 ,
\]
où $\mu^{\infty}_{X_{o},\,\mathrm{mod}} = \varprojlim_{n} \mu_{n}$, la limite
étant prise sur les $n$ \emph{inversibles dans $\mathcal{O}_{X_{o}}$} --- donc,
si $p$ est la caractéristique de $X_{o}$, un groupe isomorphe à
$\prod_{\ell \neq p} \mathbb{Z}_{\ell}$ après choix d'une
orientation.\footnote{Le manuscrit écrit l'indice \emph{tame} en anglais ; il
est rendu ici \emph{mod}, pour \emph{modéré}, conformément à l'usage du reste
du carton. La restriction aux $n$ inversibles sur $X_{o}$ est portée en marge
de la page 87 et elle est essentielle : c'est elle qui fait de
$\mu^{\infty}$ un objet \emph{modéré}, et c'est elle qui, à l'exemple II
ci-dessous, l'annulera sur $\operatorname{Spec}\mathbb{Z}$.}

Le terme $\pi_{1}^{t}(X_{o}/D'_{o})$ s'identifie à $\pi_{1}^{t}(\mathcal{U})$
et à $\pi_{1}^{t}(\mathcal{X}/D')$ : c'est le théorème de spécialisation, écrit
à la volée sous la suite. Autrement dit : \emph{ôter la fibre spéciale ne change
pas le groupe fondamental modéré}, et tout le contenu du calcul est dans
l'extension par $I_{o}$, qui est la monodromie autour de cette fibre.

L'extension est \emph{centrale} et donc déterminée par une classe
\[
  \alpha \;\in\; H^{2}\bigl(\pi_{1}^{t}(\mathcal{U}),\, I_{o}\bigr) ,
\]
une fois choisie une orientation de $I_{o}$ via celles de
$\pi_{1}^{t}(\mathcal{U})$.

\subsection*{Lemme 1 : le revêtement universel modéré n'a pas de $H^{1}$}

\textbf{Lemme 1.} Soit $\tilde{\mathcal{U}}^{t}$ le revêtement universel de
$\mathcal{U}$ défini par $\pi_{1}^{t}$. Alors
$H^{1}(\tilde{\mathcal{U}}^{t}, \mathcal{J}) = 0$ pour $\mathcal{J}$ constant,
premier aux caractéristiques considérées.

C'est ce qui rend la suite spectrale de Hochschild--Serre utilisable en bas
degré : elle donne alors
\[
  0 \to H^{2}\bigl(\pi_{1}^{t}(\mathcal{U}), I\bigr)
  \xrightarrow{\ \rho\ } H^{2}(\mathcal{U}, I) \to
  H^{2}(\tilde{\mathcal{U}}^{t}, I)^{\pi_{1}^{t}} \to
  H^{3}\bigl(\pi_{1}^{t}, I\bigr) ,
\]
de sorte que $\alpha$ est connu dès qu'on connaît $\rho(\alpha)$ et qu'on sait
que son image plus loin est nulle. Et, toujours par le lemme 1 --- c'est
l'argument de Hurewicz ---,
\[
  H^{2}(\tilde{\mathcal{U}}^{t}, I)^{\pi_{1}^{t}} \;=\;
  \operatorname{Hom}\bigl(\pi_{2}(\tilde{\mathcal{U}}^{t}), I\bigr)^{\pi_{1}^{t}}
  \;=\; \operatorname{Hom}\bigl(\pi_{2}^{t}(\mathcal{U}), I\bigr)^{\pi_{1}^{t}} .
\]
C'est pourquoi $\pi_{2}$ apparaît ici : il mesure exactement ce qui empêche la
classe $\alpha$ d'être lue sur $H^{2}$ de $\mathcal{U}$.\footnote{Le
$\pi_{2}$ dont il s'agit est celui que la page 104 définit :
pour $X$ simplement connexe, $M \rightsquigarrow H^{2}(X,M)$ est exact donc
pro-représentable, et $\pi_{2}(X) $ est l'objet qui le pro-représente ; dans le
cas général on prend le revêtement universel. C'est le $\pi_{2}$ de l'homotopie
étale, systématisé par Artin et Mazur en 1969 ; le carton l'utilise ici comme
un objet acquis.}

\subsection*{Lemme 2 : la classe est celle de l'auto-intersection}

\textbf{Lemme 2.} Soit $i : X_{o} \to \mathcal{X}$ l'inclusion et soit $\beta$
la classe de l'auto-intersection de $X_{o}$ dans $\mathcal{X}$, c'est-à-dire la
classe du fibré normal
\[
  \mathcal{L}_{X_{o}/\mathcal{X}} \;\simeq\; i^{*}\bigl(\mathcal{O}(X_{o})\bigr)
\]
dans $\mathrm{Pic}(X_{o}) = H^{1}(X_{o}, \mathbb{G}_{m})$. Soit
$\sigma : H^{1}(X_{o}, \mathbb{G}_{m}) \to H^{2}(\mathcal{U}, I)$
l'homomorphisme canonique --- image de Kummer suivie de restriction. Alors
\[
  \rho(\alpha) \;=\; \sigma(\beta) .
\]

\emph{C'est le théorème.} Il dit que la monodromie modérée autour de la fibre
spéciale est déterminée par le fibré normal de cette fibre --- c'est-à-dire par
son nombre d'auto-intersection, si $X_{o}$ est une courbe dans une
surface.\footnote{La page 89 écrit « $\beta \in H^{1}(X_{o}, \mu)$ », la page
95 « $\alpha \in H^{2}(X_{o}, \mu^{\infty}_{X_{o},\mathrm{tame}})$ » pour la
même quantité. Les deux écritures sont correctes et différentes : la classe
d'un fibré en droites est dans $H^{1}(\mathbb{G}_{m}) = \mathrm{Pic}$, sa
classe de Chern dans $H^{2}(\mu_{n})$, et la suite de Kummer
$1 \to \mu_{n} \to \mathbb{G}_{m} \to \mathbb{G}_{m} \to 1$ envoie la première
sur la seconde. La lecture ci-dessus fait la distinction que les deux pages ne
font pas.}

Grothendieck ajoute la réserve qu'il faut : cela détermine l'extension
\emph{sous réserve} de contrôler le passage de $I$ à $\mu$, c'est-à-dire de
savoir que $\mu_{o} \to I_{o}$ ne perd rien sur l'image de
$\pi_{2}^{t}(\mathcal{U})$. Il pose $I'_{o} = \mu_{o}/\operatorname{Im}
\pi_{2}^{t}(\mathcal{U})$ et observe que l'image de $\beta$ tombe alors dans
\[
  \ker\Bigl(H^{2}(\mathcal{U}, I'_{o}) \to
  H^{2}(\tilde{\mathcal{U}}^{t}, I'_{o})\Bigr) \;\simeq\;
  H^{2}\bigl(\pi_{1}^{t}, I'_{o}\bigr) ,
\]
donc définit bien une extension.

\pagerange{91}{93}
\section*{Le cas où la classe s'annule : Picard formel (pages 91 et 93)}

Les pages 91 et 93 sont écrites d'une plume très fine sur un papier fortement
traversé par le tapuscrit du verso, et une bonne part n'en est pas lisible.
L'articulation est néanmoins claire, et le calcul du milieu de la page 93 est
complet.

Il s'agit du cas où la classe $\beta$ est dans le noyau de
$H^{2}(\mathcal{X}, \mu) \to H^{2}(\mathcal{U}, \mathcal{J}_{o})$. Supposons
$\mathcal{J}_{o} \simeq \mu_{n}$ avec $n$ inversible sur $X_{o}$. L'hypothèse
dit alors que l'image de
$\beta_{o} \in \mathrm{Pic}(X_{o}) = H^{1}(X_{o}, \mathbb{G}_{m})$ dans
$H^{2}(X_{o}, \mu_{n})$ est nulle, c'est-à-dire que
\[
  \beta_{o} \;=\; n\gamma_{o} \;+\; \sum_{i} \lambda_{i}\,\beta_{o}^{i} ,
\]
où les $\beta^{i}$ sont les classes des autres composantes du diviseur.

Il faut relever cette relation de $X_{o}$ à $\mathcal{X}$, et c'est ici que le
schéma formel travaille. Soit $X_{m}$ le $m$-ième voisinage infinitésimal de
$X_{o}$. Alors :

\begin{quote}
le noyau de $\mathrm{Pic}(X_{m}) \to \mathrm{Pic}(X_{o})$ est \emph{uniquement
divisible par $n$}.
\end{quote}

La raison, que la page ne donne pas, est immédiate et vaut d'être écrite : ce
noyau est filtré par les $H^{1}(X_{o}, \mathcal{I}^{k}/\mathcal{I}^{k+1})$, qui
sont des groupes de cohomologie de $\mathcal{O}_{X_{o}}$-modules ; et la
multiplication par $n$ est bijective sur un $\mathcal{O}_{X_{o}}$-module
puisque $n$ est inversible dans $\mathcal{O}_{X_{o}}$.\footnote{C'est
exactement l'endroit où l'hypothèse de \emph{modération} paie : la restriction
« $n$ inversible sur $X_{o}$ » de la page 87, qui semblait n'être qu'une
définition, est ce qui rend ce noyau uniquement divisible et permet le
relèvement. Sans elle, tout le raisonnement s'arrête ici. La page l'assène
(« on voit que le noyau \dots\ est toujours uniquement divisible par $n$ »)
sans la relier à sa définition.}

Il en va donc de même du noyau de
\[
  \mathrm{Pic}(\mathcal{X}) \;=\; \varprojlim_{m} \mathrm{Pic}(X_{m})
  \longrightarrow \mathrm{Pic}(X_{o}) ,
\]
d'où la surjectivité de ${}_{n}\mathrm{Pic}(\mathcal{X})$ sur la $n$-torsion du
quotient, et l'existence d'un $\gamma \in \mathrm{Pic}(\mathcal{X})$ relevant
$\gamma_{o}$, avec
\[
  \beta \;=\; n\gamma + \sum_{i} a_{i}\beta_{i}
  \qquad \text{dans } \mathrm{Pic}(\mathcal{X}) .
\]

\pagerange{95}{97}
\section*{Le théorème, et deux exemples (pages 95 et 97)}

\subsection*{L'énoncé}

Une fois obtenue une relation $n\gamma = \sum_{i} c_{i}\beta_{i}$ sur
$\mathcal{X}$, on applique le \emph{théorème de Kummer} pour $\mathcal{X}/D$ :
il fournit un revêtement de $\mathcal{X}$, modéré le long de $D_{o}$, galoisien
de groupe $\mathcal{J}_{o} \simeq \mu_{n}$, trivial au-dessus de $X_{o}$. On
obtient ainsi la suite exacte complète
\[
  \pi_{2}^{t}(X_{o}/D'_{o}, \bar{\xi}_{o}) \to
  \mu^{\infty}_{X_{o},\,\mathrm{mod}}(\bar{\xi}_{o}) \to
  \pi_{1}^{t}(\mathcal{X}/D, \bar{\xi}_{1}) \to
  \pi_{1}^{t}(X_{o}/D'_{o}, \bar{\xi}_{o}) \to 0 ,
\]
et la conclusion :

\begin{quote}
$\pi_{1}^{t}(\mathcal{X}/D, \bar{\xi}_{1})$, comme extension de
$\pi_{1}^{t}(X_{o}/D'_{o}, \bar{\xi}_{o})$, est connu \emph{entièrement} en
termes de $X_{o}/D'_{o}$ et de la classe
$\alpha \in H^{2}\bigl(X_{o}, \mu^{\infty}_{X_{o},\,\mathrm{mod}}\bigr)$,
qui est l'auto-intersection de $X_{o}$ dans $\mathcal{X}$.
\end{quote}

\subsection*{Exemple I : une courbe sur un corps algébriquement clos}

Soit $X_{o}$ une courbe sur $k$ algébriquement clos. Alors
$\pi_{2}^{t}(X_{o} - D'_{o}) = 0$ \emph{sauf} si $X_{o} \simeq
\mathbb{P}^{1}_{k}$ et $D'_{o} = \varnothing$, auquel cas
$\pi_{2}^{t}(X_{o}) = \mu^{\infty}_{X_{o},\,\mathrm{mod}}$. Deux cas, donc, et
ils sont opposés.

\begin{enumerate}
  \item[1)] Si $X_{o} - D'_{o} \neq \mathbb{P}^{1}_{k}$ : alors
    $\pi_{1}^{t}(\mathcal{X}/D, \bar{\xi}_{1})$ est une extension
    \emph{triviale} de $\pi_{1}^{t}(X_{o} - D_{o}, \bar{\xi}_{o})$ par
    $\mu^{\infty}_{\mathrm{mod}}(\bar{\xi}_{o})$. La monodromie et le groupe
    fondamental de la fibre se séparent, et il n'y a rien à calculer.
  \item[2)] Si $X_{o} - D'_{o} = \mathbb{P}^{1}_{k}$, c'est-à-dire
    $X_{o} = \mathbb{P}^{1}_{k}$ et $D'_{o} = \varnothing$ : alors
    \[
      \pi_{1}^{t}(\mathcal{X}/D, \bar{\xi}_{1}) \;\simeq\; \mu_{n}(k) ,
    \]
    où $n$ est la partie première à $p = \operatorname{car} k$ du degré de
    l'auto-intersection $X_{o} \cdot X_{o}$ dans $\mathcal{X}$.
\end{enumerate}

Le second cas est celui qu'il faut retenir, et c'est le plus concret du carton.
Une droite projective d'auto-intersection $-d$ dans une surface fibrée a, pour
groupe fondamental modéré de son voisinage formel épointé, le groupe des
racines $n$-ièmes de l'unité, $n$ étant la partie de $d$ première à la
caractéristique. La monodromie est cyclique d'ordre $d$ : c'est ce que l'on
attend d'une singularité quotient cyclique, et le calcul le rend.

\subsection*{Exemple II : les entiers}

\begin{enumerate}
  \item[a)] $X_{o} = \operatorname{Spec}\mathbb{Z}$. Alors
    $\mu^{\infty}_{\mathrm{mod},\,X_{o}} = 0$, donc
    \[
      \pi_{1}^{t}(\mathcal{X}/D, \bar{\xi}_{1}) \;\simeq\;
      \pi_{1}^{t}(X_{o}/D'_{o}, \bar{\xi}_{o}) :
    \]
    il n'y a pas de monodromie du tout.\footnote{La page l'affirme sans dire
    pourquoi, et la raison est la définition même : $\mu^{\infty}_{\mathrm{mod}}$
    est la limite des $\mu_{n}$ pour $n$ \emph{inversible dans}
    $\mathcal{O}_{X_{o}}$ ; sur $\operatorname{Spec}\mathbb{Z}$, aucun entier
    $n > 1$ n'est inversible, et la limite est prise sur la seule valeur
    $n = 1$. C'est l'exemple qui montre le mieux ce que « modéré » interdit.}
    Si $D'_{o} = \varnothing$, on trouve $0$ --- c'est le théorème de
    Minkowski : $\mathbb{Q}$ n'a pas d'extension non ramifiée. Si
    $D'_{o} \neq \varnothing$, la partie abélienne se calcule par la théorie du
    corps de classes et vaut
    \[
      \prod_{\ell \in D'_{o}} (\mathbb{Z}/\ell\mathbb{Z})^{*} ,
    \]
    les extensions abéliennes de $\mathbb{Q}$ modérément ramifiées en les seuls
    $\ell \in D'_{o}$ étant les sous-corps de
    $\mathbb{Q}(\zeta_{\prod \ell})$. Il ajoute que le groupe entier doit être
    fini par dévissage, et pose une question qu'il laisse ouverte sur les
    corps de nombres modérément ramifiés sur $\mathbb{Q}$.
  \item[b)] $X_{o} = \operatorname{Spec}\mathbb{Z} - S$, avec
    $S \neq \varnothing$. Alors
    \[
      \mu^{\infty}_{\mathrm{mod},\,X_{o}} \;=\;
      \prod_{\ell \in S} \mu^{\infty}(\ell)_{X_{o}} ,
    \]
    puisque les $n$ inversibles sur $\mathbb{Z}[1/S]$ sont exactement ceux dont
    les facteurs premiers sont dans $S$. Et la page se termine sur une
    question, à laquelle rien ne répond : « Mais $\pi_{2}^{t}(X_{o}/D'_{o})$
    ??? ».
\end{enumerate}

\pagerange{99}{100}
\section*{Abhyankar semi-global, et la même classe une deuxième fois (pages 99 et 100)}

Deux feuillets non paginés ouvrent une question nouvelle, et le second
réénonce, sous forme de programme, le théorème qu'on vient de démontrer.

\subsection*{Une version semi-globale du résultat d'Abhyankar}

Soit $X$ un schéma localement noethérien connexe et $D = \sum D_{i}$ un
diviseur à croisements normaux à composantes régulières. Supposons
$Z = \bigcap_{i} D_{i} \neq \varnothing$ et le couple $(X, Z)$
\emph{hensélien}. Posons $\mathcal{U} = X - D$ et soit $\xi$ un point
géométrique de $\mathcal{U}$. Alors :

\begin{enumerate}
  \item[a)] $\mathcal{U}$ est connexe, et
    $\pi_{1}(\mathcal{U})_{\mathrm{mod}} \to \pi_{1}(X)$ est surjectif ;
  \item[b)] pour chaque $D_{i}$, le groupe d'inertie relatif à $D_{i}$ est un
    sous-groupe distingué \emph{commutatif} $I_{i}$ de
    $\pi_{1}(\mathcal{U})_{\mathrm{mod}}$, et l'on a un épimorphisme canonique
    \[
      \prod_{\ell} T_{\ell}(\Omega^{*}) \longrightarrow I_{i} ,
    \]
    le produit étant pris sur les $\ell$ distincts de la caractéristique en
    $D_{i}$ ;
  \item[c)] les $I_{i}$ commutent entre eux, et leur \emph{produit} est le
    noyau de
    $\pi_{1}(\mathcal{U}, \xi)_{\mathrm{mod}} \to \pi_{1}(X, \xi)$.
\end{enumerate}

L'énoncé a) est immédiat. Pour b), l'invariance des $D_{i}$ dans un revêtement
modérément ramifié $X'$ de $X$ est de nature étale locale et $X'$ est connexe
donc irréductible, ce qui identifie $I_{i}$ ; le fait qu'il soit quotient du
produit des modules de Tate est le calcul local classique. Pour c), la
commutation des $I_{i}$ est locale en un point du
croisement.\footnote{Cette description --- le groupe fondamental modéré du
complémentaire d'un diviseur à croisements normaux dans un couple hensélien
est une extension de $\pi_{1}(X)$ par un produit de modules de Tate, un par
composante --- est le calcul qui, en géométrie logarithmique, donne le groupe
fondamental \emph{log}-étale d'un point logarithmique. Le carton en fait ici
une « version semi-globale du résultat d'Abhyankar » ; c'est le même énoncé.}

\subsection*{Le programme kummérien : la classe, une deuxième fois}

La page 100 reprend, pour un $X$ quelconque, ce que le lemme 2 démontrait pour
un schéma formel.

Soient $X$ un schéma, $(D_{i})_{1 \leqslant i \leqslant r}$ des diviseurs de
Cartier $\geqslant 0$, $(n_{i})$ des entiers $\geqslant 1$,
$\mu_{\underline{n}} = \prod_{i} \mu_{n_{i}}$ sur $X$, $\Gamma$ un schéma en
groupes, et une extension de groupes sur $X$
\[
  1 \longrightarrow \mu_{\underline{n}} \longrightarrow G \longrightarrow
  \Gamma \longrightarrow 1 .
\]
Donnons-nous $X' \to \tilde{X}_{1} \to X$ avec $G$ opérant sur $\tilde{X}_{1}$,
$\mu_{\underline{n}}$ y opérant trivialement, de sorte que $\tilde{X}_{1}$
devienne un \emph{torseur} sous $\Gamma$ ; et $X'$ kummérien au-dessus de
$\tilde{X}_{1}$ relativement aux $D_{i}$, c'est-à-dire
\[
  X' \;=\; \operatorname{Spec}
  \mathcal{O}_{\tilde{X}}[T_{1}, \dots, T_{r}]\,\big/\,
  \bigl(T_{1}^{n_{1}} - a_{1},\ \dots,\ T_{r}^{n_{r}} - a_{r}\bigr) .
\]
L'extension $G$ de $\Gamma$ par $\mu_{\underline{n}}$ définit une classe
$\alpha \in H^{2}(B_{\Gamma}, \mu_{\underline{n}})$, et le torseur
$\tilde{X}$ fournit un homomorphisme
\[
  \varphi : H^{2}(B_{\Gamma}, \mu_{\underline{n}}) \longrightarrow
  H^{2}(X, \mu_{\underline{n}}) \;\simeq\;
  \prod_{i} H^{2}(X, \mu_{n_{i}}) .
\]

\begin{quote}
\textbf{Ce qu'on veut prouver.} Si $\varphi(\alpha) = (\xi_{i})_{1 \leqslant i
\leqslant r}$, alors
\[
  \xi_{i} \;=\; c\bigl(\mathcal{O}_{X}(D_{i})\bigr) \pmod{n_{i}} .
\]
\end{quote}

C'est le lemme 2 de la page 89, écrit sans schéma formel : \emph{la classe de
l'extension d'inertie est la classe de Chern du diviseur}. La page 102 en donne une rédaction en corollaire ; et la page 200, tout à la fin du
carton, l'écrit une quatrième fois, pour la gerbe des racines
$\underline{n}$-ièmes des $D_{i}$ : « L'existence d'une section dépend de la
nullité d'un élément de $H^{2}(X, \mu_{\underline{n}})$, qui n'est autre,
comme on devine, que la classe de cohomologie des $D_{i}$ ». Quatre énoncés,
cent dix pages, une seule identité.

La page se clôt sur une ligne qui indique où cela devait aller : « Variantes
\dots\ formelle, ou \dots\ loc.\ annelés ».

\pagerange{101}{102}
\section*{Deux pages isolées : un symbole local, et le théorème kummérien redit (pages 101 et 102)}

\subsection*{Le symbole local}

La page 101 est le carbone d'un tapuscrit français, sans lien avec le reste du
carton. Elle achève la démonstration d'une réciprocité pour les symboles locaux
sur une variété abélienne $A$ : pour deux cycles $\underline{a}$,
$\underline{b}$ et un diviseur $X$, on veut
\[
  j_{o}(X_{\underline{a}}, \underline{b}) \;=\;
  j_{o}(X^{-}_{\underline{b}}, \underline{a}) .
\]
Le cas où $\underline{b}$ appartient au noyau d'Albanese se traite par
linéarité et par l'invariance du symbole sous les translations, ce qui ramène
au cas déjà traité ; le cas général se ramène, en modifiant $\underline{b}$, à
celui où les supports sont étrangers, et l'on conclut par l'invariance du
symbole $i$ sous l'application $\zeta(x) = a + b - x$.

C'est de la théorie des symboles locaux au sens de Serre, et le carton n'y
revient pas.

\subsection*{Le théorème kummérien, troisième énoncé}

La page 102 --- numérotée 3 par lui, dans un cercle --- reprend en corollaire le
programme de la page 100. Elle est écrite très vite et raturée d'un bout à
l'autre ; la transcription en laisse une large part illisible et cette lecture
n'invente pas ce qui manque. Ce qui se lit est l'articulation, et c'est elle
qui compte.

Soit $X$ un schéma normal localement intègre, $(D_{i})$ des diviseurs de
Cartier, et $X'$ un revêtement dont les groupes d'inertie relatifs aux $D_{i}$
sont tels que $X'$ soit modérément ramifié le long des $D_{i}$. On a alors une
extension
\[
  1 \longrightarrow I \longrightarrow G \longrightarrow \Gamma
  \longrightarrow 1 ,
\]
avec $\tilde{X} = X'/I$ revêtement principal de $X$ de groupe $\Gamma$, d'où
une classe $\alpha \in H^{2}(\Gamma, I)$, puis une classe
\[
  \alpha \in H^{2}\bigl(X,\ \tilde{X} \times^{\Gamma} I\bigr) .
\]
On dispose d'un isomorphisme canonique
$\tilde{D} \times^{\Gamma} I \simeq \prod_{i} \mu_{n_{i},\,D}$ sur $D$,
compatible aux opérations de $\Gamma$, et la conclusion visée est la même que
celle des pages 89, 100 et 200 : la classe est celle du fibré normal,
\[
  c\bigl(\mathcal{O}_{X}(D)\bigr) \in H^{2}(X, \mu_{\underline{n}}) ,
  \qquad
  c\bigl(T_{D/X}\bigr) \in H^{2}(D, \mu_{\underline{n}}) .
\]

\pagerange{103}{104}
\section*{Le théorème de monodromie locale (pages 103 et 104)}

\subsection*{La relation, et d'où elle vient}

La page 103 est barrée de deux longs traits diagonaux, et c'est la seule page de cette suite qui se lise proprement de bout en bout.

Le contexte est celui d'un corps local de corps résiduel fini à $q$ éléments.
Le quotient modéré de son groupe de Galois est engendré par deux éléments : un
relèvement $f$ du Frobenius, et un générateur $g$ de l'inertie modérée --- qui,
après choix d'isomorphismes $\mathbb{Z}_{\ell}(1) \simeq \mathbb{Z}_{\ell}$
pour tout $\ell$, engendre topologiquement
$\prod_{\ell \neq p} \mathbb{Z}_{\ell}$. Ils satisfont la relation
\[
  (6.4) \qquad f\,g\,f^{-1} \;=\; g^{q} ,
\]
qui exprime que le Frobenius multiplie le tour par $q$. Le groupe est le
produit semi-direct de $\mathbb{Z}$ par $\prod_{\ell\neq p}\mathbb{Z}_{\ell}$,
$\mathbb{Z}$ opérant par multiplication par $q$.

\subsection*{Le lemme}

\textbf{Lemme.} Soient $F$, $G$ deux \emph{automorphismes} d'un espace
vectoriel $E$ de dimension finie sur un corps, et $q > 1$ un entier tels que
\[
  F G F^{-1} \;=\; G^{q} .
\]
Alors $G$ est \emph{quasi-unipotent} : ses valeurs propres sont toutes des
racines de l'unité ; autrement dit, il existe $N \geqslant 1$ tel que $G^{N}$
soit unipotent.\footnote{La page écrit « deux endomorphismes », tout en
utilisant $F^{-1}$ dans la relation et en plaçant l'ensemble $\Phi$ des valeurs
propres dans le groupe multiplicatif $\bar{R}^{*}$. Il faut $F$ inversible pour
que la relation ait un sens, et $G$ inversible pour la conclusion : si $0$
était valeur propre de $G$, il le serait de $G^{q}$ et l'argument ne
l'excluerait pas, alors que $0$ n'est pas racine de l'unité. Les deux
hypothèses sont supplées ici.}

\emph{Démonstration.} L'ensemble $\Phi$ des valeurs propres de $G$ dans une
clôture algébrique du corps de base est aussi celui de $FGF^{-1}$, donc celui
de $G^{q}$ : c'est donc une partie finie de $\bar{R}^{*}$ stable par
$\lambda \mapsto \lambda^{q}$. Comme $\Phi$ est fini, pour tout
$\lambda \in \Phi$ il existe deux entiers $N' \neq N''$ avec
$\lambda^{q^{N'}} = \lambda^{q^{N''}}$, d'où $\lambda^{q^{N'} - q^{N''}} = 1$ :
$\lambda$ est une racine de l'unité.

C'est le \emph{théorème de monodromie locale}, dans sa forme purement
algébrique : la relation de commutation entre Frobenius et inertie interdit à
la monodromie d'être quelconque. Il n'y a pas d'hypothèse de géométrie dans le
lemme --- seulement la relation $(6.4)$ et la finitude de la
dimension.\footnote{C'est le lemme qui figure dans SGA 7, exposé I, comme
première étape du théorème de monodromie ; il y est attribué à Grothendieck et
la démonstration est celle-ci. Que la page du carton soit barrée de deux
diagonales n'est pas un jugement sur l'énoncé : c'est la marque, fréquente dans
ce carton, d'un texte reporté ailleurs.}

\subsection*{$\pi_{2}$}

La page 104 définit, en cinq lignes, l'objet dont les pages 87 à 97 se servaient sans le définir.

\begin{enumerate}
  \item[a)] Si $X$ est simplement connexe, le foncteur
    $M \rightsquigarrow H^{2}(X, M)$, sur les groupes abéliens finis, est
    exact --- $H^{1}(X, M)$ étant nul --- donc pro-représentable, et l'on pose
    \[
      H^{2}(X, M) \;=\; \operatorname{Hom}_{\mathrm{cont}}
      \bigl(\pi_{2}(X),\, M\bigr) .
    \]
  \item[b)] Dans le cas général, $\pi_{2}(X, \xi) = \pi_{2}\bigl(\tilde{X}(\xi)
    \bigr)$, où $\tilde{X}(\xi)$ est le revêtement universel pointé défini par
    $\xi$.
\end{enumerate}

C'est le $\pi_{2}$ de l'homotopie étale, et le procédé --- définir un groupe
d'homotopie comme le pro-objet qui pro-représente un foncteur de cohomologie
--- est exactement celui que la lettre à Serre de 1959 emploie pour $\pi_{1}$
(page 126) et que les pages 164 à 185 systématiseront. Le bas de la page porte
un croquis en perspective de trois diviseurs se coupant à croisements normaux ;
il n'est pas repris, n'étant pas un diagramme.

\pagerange{105}{110}
\section*{Un feuillet d'EGA IV, et de l'algèbre commutative (pages 105 à 110)}

\subsection*{« Attention à l'ambiguïté dans \emph{lisse} ! »}

La page 105 est le feuillet IV-878 du tapuscrit d'EGA IV, portant ses
corrections. Il court de la proposition (17.4.10) --- un morphisme localement
de présentation finie est non ramifié si et seulement si ses fibres le sont ---
jusqu'au théorème (17.5.1), qui donne les quatre caractérisations équivalentes
de la lissité en un point : lissité, platitude plus lissité de la fibre,
régularité du morphisme, et lissité formelle de
$\mathcal{O}_{X,x}$ comme $\mathcal{O}_{Y,y}$-algèbre pour les topologies
discrètes.

Deux corrections sont substantielles : « localement \emph{noethérien} » est
biffé et remplacé par « localement \emph{de type fini} » dans l'argument sur la
fibre, et les deux équivalences sont étiquetées en marge. La troisième est un
avertissement, et il vaut pour tout le carton : « Attention à l'ambiguïté dans
« lisse » ! » --- le mot ayant servi, dans les premières rédactions, tantôt
pour \emph{lisse} et tantôt pour \emph{simple}.\footnote{L'ambiguïté n'est pas
théorique : la page 74 du carton, dans « Pic birationnel », écrit
« $X$ propre et \emph{simple} sur $k$ » là où le corollaire 1 demande
\emph{lisse}, et le corollaire 2 de la même page demande un « modèle propre et
\emph{simple} ». C'est l'un des rares endroits où l'avertissement de la marge
d'EGA porte sur le carton lui-même.}

\subsection*{Le $\pi_{1}$ local le long de deux branches}

La page 106 dessine deux branches $X_{T'}$ et $X_{T''}$ se coupant en un point
$p$, et un petit disque $Z = \operatorname{Spec}
\hat{\mathcal{O}}_{X,a}$ autour du croisement. Les deux traces
$Z' = X_{T'} \cap Z$ et $Z'' = X_{T''} \cap Z$ correspondent aux corps
résiduels $k(p')$ et $k(p'')$, qui sont les corps locaux en $a$ des deux
branches. Le diagramme de groupes fondamentaux qui en résulte,

\begin{tikzcd}[column sep=small, row sep=small, nodes={font=\scriptsize}]
 & \hat{\mathbb{Z}} = \pi_{1}(Z') \arrow[dl] \arrow[dr] & \pi_{1}(Z'') = \hat{\mathbb{Z}}
 \arrow[dl] \arrow[d] \\
\pi_{1}(X_{T'}) \arrow[dr] & \pi_{1}(X_{T''}) \arrow[d] & \pi_{1}(Z) =
\hat{\mathbb{Z}}^{2} \arrow[dl] \\
 & \pi_{1}(X_{T}) &
\end{tikzcd}

est celui qu'un théorème de descente devrait clore, et la page le dit sans le
faire : « On doit prouver un th.\ de descente de nature essentiellement locale,
plus déterminer certains hom.\ canoniques ». C'est la situation locale des
pages 41 et 87, vue une troisième fois.

\subsection*{Polynômes unitaires et idéaux étrangers}

Les pages 107, 109 et 110 sont de l'algèbre commutative élémentaire, mise en
place pour un usage qui n'est pas dit. Le point est le suivant. Pour deux
idéaux $Q$, $Q'$ d'un anneau commutatif $A$, on a l'équivalence
\[
  Q + Q' = A \;\Longleftrightarrow\; Q \cap Q' = Q\,Q' ,
\]
et sous cette condition $A/QQ' \to A/Q \times A/Q'$ est un isomorphisme. D'où :
les décompositions de $B/R$ en produit de deux anneaux correspondent aux
couples d'idéaux $J$, $J'$ avec $JJ' = R$ et $J + J' = B$ ; et si $B = A[T]$ et
$R = (P)$ avec $P$ unitaire, ces décompositions sont exactement les
factorisations $P = QQ'$ en polynômes unitaires. L'unicité permet de se ramener
au cas local, et l'existence se voit alors sur le complété.

C'est le mécanisme de factorisation dont on se sert pour décomposer une algèbre
finie étale en facteurs connexes ; le carton s'en servira à la page 152, où la
proposition 3.2 caractérise les morphismes finis étales de rang $n$.

\pagerange{112}{115}
\section*{« Théorème de Bertini et variantes » (pages 112 à 115)}

Un tapuscrit corrigé de sa main. Il énonce ce que la page 66 --- « $\pi_{1}$ birationnel » --- utilisait sans le démontrer, puis va plus loin.

Dans tout ce qui suit $k$ est un corps et $\mathbf{P}^{r} = \mathbf{P}^{r}_{k}$.

\subsection*{Les deux énoncés de base}

\textbf{Théorème de Bertini 1.} Soit $f : X \to \mathbf{P}^{r}_{k}$ un
morphisme, $X$ un schéma algébrique géométriquement irréductible, et
$n = \dim \overline{f(X)}$. Soit $L_{t}$ le sous-espace linéaire
\emph{générique} de codimension $s \leqslant n-1$. Alors
$Y_{t} = f_{t}^{-1}(L_{t})$, dans $X_{t} = X \otimes_{k} k(t)$, est
géométriquement irréductible.

\textbf{Énoncé local 2.} Soit $f : X \to \mathbf{P}^{r}$ avec $X$ algébrique.
Soit $Z \subset X$ l'ensemble des points où $f$ est ramifié \emph{ou} où $X$
n'est pas géométriquement unibranche --- respectivement pas normal, pas lisse.
Si $\dim \overline{f(Z)} < s$, alors $Y_{t} = f_{t}^{-1}(L_{t})$ est
géométriquement unibranche --- respectivement normal, lisse.

Le second est le mécanisme de tout ce qui suit, et il faut voir sa forme : la
propriété qu'on veut sur la section générique s'obtient en \emph{écartant} par
la dimension le lieu où elle échoue sur $X$.

\subsection*{Les corollaires}

\textbf{Corollaire 3.} Sous les conditions de 1 et 2, avec $k$ algébriquement
clos, le morphisme canonique $\pi_{1}(Y_{\bar{t}}) \to \pi_{1}(X)$ est
\emph{surjectif}.

\textbf{Remarque 4.} Utilisant les théorèmes de changement de base pour le
$\pi_{1}$, on conclut que
\[
  (\ast) \qquad \pi_{1}(Y_{\bar{t}}) \longrightarrow \pi_{1}(X_{\bar{t}})
\]
est surjectif \emph{si $X$ est propre}, et que
\[
  (\ast\ast) \qquad \pi_{1}^{!}(Y_{\bar{t}}) \longrightarrow
  \pi_{1}^{!}(X_{\bar{t}})
\]
est surjectif sans hypothèse de propreté, où $!$ désigne le plus grand quotient
profini premier à l'exposant caractéristique de $k$.\footnote{Le tapuscrit
écrit, pour $(\ast\ast)$, « $\pi_{1}^{!}(X_{\bar{t}}) \to
\pi_{1}^{!}(X_{\bar{t}})$ », les deux membres identiques. Le membre de gauche
doit être $Y_{\bar{t}}$, comme en $(\ast)$ ; c'est une faute de frappe et elle
est corrigée ici.} On \emph{n'affirme pas} que $(\ast)$ soit toujours surjectif
sans propreté : c'est exactement la différence que $(\ast\ast)$ mesure, et elle
est du côté sauvage.

\textbf{Corollaire 4 (de 1).} Si $X$ est de plus propre, alors pour
\emph{toute} sous-variété linéaire $L_{t}$ de codimension $s$ ---
générique ou non --- $Y_{t} = f_{t}^{-1}(L_{t})$ est géométriquement connexe.
C'est le théorème de connexité.

\textbf{Corollaire 5.} Sous les hypothèses de 3, avec $X$ propre,
$\pi_{1}(Y_{\bar{t}}) \to \pi_{1}(X)$ est surjectif, l'identification
$\pi_{1}(X) \simeq \pi_{1}(X_{\bar{t}})$ venant du théorème de changement de
base. En marge : « cf.\ SGA 1 X 2.11 ».

\subsection*{Le théorème 6 : la version modérée, pour $t$ non générique}

C'est le résultat propre du tapuscrit, et il répond au besoin qu'énonce la page
113 : un énoncé de surjectivité pour un $t$ \emph{non} générique, sans
hypothèse de propreté, et pour des quotients plus gros que ceux de
$(\ast\ast)$.

\textbf{Théorème 6.} Supposons $k$ algébriquement clos, $X$ propre,
$f : X \to \mathbf{P}^{r}$, $L_{t}$ un sous-espace linéaire de codimension $s$
défini sur $k$, et $D \subset X$ un sous-schéma fermé, tels que :

\begin{enumerate}
  \item[a)] $Y_{t} = f^{-1}(L_{t})$ est non vide ;
  \item[b)] pour tout $x \in Y_{t}$ : $X$ est lisse en $x$, $Y_{t}$ est lisse
    en $x$ et de codimension $s$ dans $X$ en $x$, et $f$ est net en $x$ ;
  \item[c)] pour tout $x \in D_{t} = D \cap Y_{t}$ : ou bien
    $\operatorname{codim}_{x}(D_{t}, Y_{t}) \geqslant 2$, ou bien $D$ est un
    diviseur sur $X$ en $x$ qui induit sur $Y_{t}$ en $x$ un diviseur à
    croisements normaux.
\end{enumerate}

Supposons $U = X - D$ irréductible et posons $U_{t} = U \cap Y_{t}$. Soit
$\pi^{t}_{1}(U)$ le quotient de $\pi_{1}(U)$ classifiant les revêtements
modérément ramifiés le long des composantes de codimension 1 de $D$ qui
rencontrent $Y_{t}$, et $\pi^{t}_{1}(U_{t})$ le quotient analogue pour
$D_{t}$. Alors
\[
  \pi^{t}_{1}(U_{t}) \longrightarrow \pi^{t}_{1}(U)
\]
est \emph{surjectif}.

\emph{Démonstration.} Pour $L_{t'}$ générique la conclusion est le corollaire
3 : un revêtement étale connexe $V$ de $U$ reste connexe au-dessus de
$U_{t'}$. Il faut la faire descendre à $t$. Sur un voisinage ouvert $T$ de
$\bar{t}$ dans la grassmannienne, la famille des $Y_{t'}$ est un morphisme
propre et lisse $Y \to T$, et celle des $D_{t'}$ un sous-schéma fermé $\Delta$
de $Y$, réunion d'un $\Delta_{1}$ de codimension $\geqslant 2$ sur chaque
fibre --- qui ne compte pas, par le théorème de pureté relative --- et d'un
diviseur $\Delta_{2}$ à croisements normaux relatifs. Le revêtement $V$
définit un revêtement $V'$ de $U' = Y - \Delta$ ; par pureté, $V'$ se prolonge
en un revêtement étale $V''$ de $U'' = Y - \Delta_{2}$, modérément ramifié sur
les fibres relativement à $\Delta_{2}$. On conclut par le \emph{théorème de
spécialisation pour le groupe fondamental modéré}.\footnote{Le tapuscrit
renvoie à « SGA 1 XIII (exposé de Mme Raynaud) », au futur : « qui
figure(ra) ». C'est le second renvoi du carton à ce théorème, après celui de la
page 22 ; les deux le citent comme l'outil qui manque, et la ramification
modérée au sens large de la page 21 est ce qui le rend applicable.}

\subsection*{Remarque 7 : où il bute}

Avec les notations de SGA 7 XVIII, le théorème 6 s'applique à tout pinceau de
Lefschetz $D$ et donne la surjectivité de
\[
  \pi^{t}_{1}(D - \check{D}) \longrightarrow
  \pi^{t}_{1}(\mathbf{P} - \check{X}) ,
\]
d'où le théorème de conjugaison des cycles évanescents --- donc
l'irréductibilité de la cohomologie évanescente --- \emph{pourvu qu'on ne soit
pas dans le cas $\operatorname{car} k = 2$ avec $\dim X$ impaire}.

Dans ce cas-là, dit-il, « je ne sais pas prouver la conjugaison sous
$\pi_{1}(D - \check{D})$, sauf pour un pinceau \emph{générique}, puisqu'alors
la représentation envisagée est sauvage ; j'ignore si elle reste vraie pour un
pinceau \emph{général} ». Il ajoute que la conjugaison des sous-groupes de
monodromie locale \emph{dans} $\operatorname{Aut}(H^{n-1}(Y))$, qui est très
élémentaire, reste vraie et suffit pour le théorème caractérisant le cas où la
condition (A) n'est pas vérifiée.

C'est l'un des deux endroits du carton où Grothendieck écrive noir sur blanc
qu'il ne sait pas ; l'autre est la page 179.

\pagerange{116}{120}
\section*{« Comportement fonctoriel des groupes d'inertie » (pages 116 à 120)}

Un second tapuscrit corrigé, en cinq numéros dont le cinquième déborde sur la
page 121. Il traite ce que la lettre de 1967 signalait comme la cause de
l'effondrement du théorème 6.4 de Murre.

\subsection*{1. Ce qu'est un groupe d'inertie, et de quoi il dépend}

Soient $X$ un schéma, $U$ un ouvert, $G$ un groupe profini, $U'$ un revêtement
principal de $U$ de groupe $G$, $x$ un point géométrique de $X$, $\bar{X}$ le
localisé strict de $X$ en $x$, $\bar{U}$ l'image inverse de $U$ et
$\bar{U}' = U' \times_{U} \bar{U}$. Donnons-nous de plus un $\bar{X}$-schéma
$\bar{X}'$ à opérateurs $G$ prolongeant $\bar{U}'$.

Le cas d'intérêt est celui où $x \in X - U$ est adhérent à $U$, où $\bar{X}'$
est entier sur $\bar{X}$, et où
$\mathcal{O}_{\bar{X}} = \bar{p}_{*}(\mathcal{O}_{\bar{X}'})^{G}$ --- de sorte
que $G$ opère transitivement sur la fibre $\bar{X}'_{\bar{x}}$. Si $X$ est
normal en $x$ et $\bar{U}$ non vide, on prend pour $\bar{X}'$ le
\emph{normalisé} de $\bar{X}$ relativement à $\bar{U}'$, comme limite projective
des normalisés dans les quotients finis.

Pour $\xi \in \bar{X}'_{\bar{x}}$, le stabilisateur
\[
  G_{\xi} \;=\; \{\, g \in G \;:\; g\cdot\xi = \xi \,\}
\]
est le \emph{groupe d'inertie relatif à $\xi$}. Quand $G$ opère transitivement
sur la fibre, les $G_{\xi}$ forment une classe de sous-groupes conjugués, qu'on
appelle par abus \emph{les} sous-groupes d'inertie relatifs à $x$.

Il faut lire l'avertissement qui suit, parce qu'il est le sujet de tout ce
tapuscrit : \emph{cette classe n'est pas déterminée par la seule connaissance
de $x$}, elle dépend du choix du prolongement $\bar{X}'$ de $\bar{U}'$. La
convention --- $X$ normal en $x$, $x$ spécialisation d'un point de $U$,
prolongement par normalisation --- est ce qui la fixe, et elle est sous-entendue
partout ensuite.

\subsection*{2. Fonctorialité}

Soit une seconde situation $(Y, V, H, V')$, un morphisme $f : X \to Y$ avec
$f(U) \subset V$, un morphisme $f' : U' \to V'$ au-dessus de lui, et un
homomorphisme $\varphi : G \to H$ rendant $f'$ équivariant. Le cas courant est
celui où $U'$ et $V'$ sont des quotients des revêtements universels de $U$ et
$V$ pointés en $a$ et $b = f(a)$, où $G$ et $H$ sont des quotients de
$\pi_{1}(U,a)$ et $\pi_{1}(V,b)$, et où $\varphi$ est induit par
$f_{*}$ ; alors $f'$ et $\varphi$ sont déterminés sans ambiguïté par $f$, et
leur existence signifie seulement que $f_{*}$ passe aux quotients.

Étant donné un prolongement $\bar{f}' : \bar{X}' \to \bar{Y}'$, on a alors la
relation évidente
\[
  (2.1) \qquad \varphi(G_{\xi}) \;\subset\; H_{\bar{f}'(\xi)} :
\]
un morphisme envoie inertie dans inertie. Et si les deux actions sont
transitives sur les fibres, $\varphi$ applique tout groupe d'inertie de $G$
relatif à $x$ dans un groupe d'inertie de $H$ relatif à $y = f(x)$.

Reste l'existence de $\bar{f}'$. En pratique $\bar{U}'$ est schématiquement
dense dans $\bar{X}'$ et $\bar{Y}'$ est séparé sur $\bar{Y}$, ce qui rend
$\bar{f}'$ unique s'il existe ; et il existe --- en l'interprétant comme
$\bar{X}' \to \bar{Y}' \times_{\bar{Y}} \bar{X}'$ --- lorsque $\bar{Y}'$ est
entier sur $\bar{Y}$, par EGA II 6.1.14. La commutation aux actions de $G$ est
alors automatique.

\subsection*{3. Le calcul local}

Supposons $\mathcal{O}_{X,x}$ régulier et $\bar{U} = \bar{X} - \bar{D}$ avec
$\bar{D}$ un diviseur à croisements normaux dans le schéma strictement local
régulier $\bar{X}$, et supposons la ramification de $\bar{U}'$ \emph{modérée}
aux points maximaux de $\bar{D}$. Le calcul local classique donne alors : pour
tout $\xi$, le groupe d'inertie $G_{\xi}$ est canoniquement un \emph{quotient}
du groupe
\[
  (3.1) \qquad I \;=\; \prod_{\ell \neq p}
  \underline{Z}_{\ell}(1)\bigl(k(\bar{x})\bigr)^{C}
  \;=\; \prod_{\ell \neq p} T_{\ell}\bigl(k(\bar{x})\bigr)^{C} ,
\]
où $C$ est l'ensemble des composantes irréductibles de $\bar{D}$ et $p$ la
caractéristique résiduelle en $x$. \emph{Un facteur par branche du diviseur} :
c'est la structure entière de l'inertie modérée.

Le même calcul montre que le facteur correspondant à une composante
$\bar{D}_{c}$ coïncide avec le groupe d'inertie en n'importe quel point de
$\bar{X}'$ au-dessus du point maximal de $\bar{D}_{c}$ qui générise $\xi$ ---
et de tels points existent toujours. Donc \emph{la classe de conjugaison de ce
facteur ne dépend que de la composante irréductible de $D = X - U$}. C'est
l'énoncé qui manquait à Murre, et c'est celui que la page 7 annonçait.

\subsection*{4. La matrice des multiplicités}

Reprenons la situation de 2 avec, des deux côtés, la structure de 3 :
$\bar{V} = \bar{Y} - \operatorname{supp}\bar{E}$, $\bar{E}$ à croisements
normaux, $C'$ l'ensemble de ses composantes. Pour tout $c' \in C'$, l'image
inverse de $\bar{E}_{c'}$ est un diviseur sur $\bar{X}$ de support contenu dans
celui de $\bar{D}$, donc
\[
  \bar{f}^{*}(\bar{E}_{c'}) \;=\; \sum_{c \in C} n_{c,c'}\,\bar{D}_{c} ,
\]
les $n_{c,c'}$ étant des entiers $\geqslant 0$ non tous nuls pour $c'$ fixé.
Un calcul d'images inverses de revêtements kummériens donne alors la
commutativité de

\begin{tikzcd}[column sep=small, row sep=small, nodes={font=\scriptsize}]
G_{\xi} \arrow[r, "\varphi"] & H_{y} \\
\prod_{\ell} \underline{Z}_{\ell}(1)(\bar{x})^{C} \arrow[u] \arrow[r, "N"] &
\prod_{\ell} \underline{Z}_{\ell}(1)(\bar{y})^{C'} \arrow[u]
\end{tikzcd}

les flèches \emph{verticales} étant les épimorphismes canoniques de (3.1) et
$N$ l'homomorphisme de matrice $(n_{c,c'})$, compte tenu que
$k(\bar{y}) \to k(\bar{x})$ induit
$\underline{Z}_{\ell}(1)(\bar{y}) \xrightarrow{\ \sim\ }
\underline{Z}_{\ell}(1)(\bar{x})$.\footnote{Le tapuscrit écrit d'abord
« horizontales », corrigé de sa main en « verticales » : ce sont bien les
projections $I \to G_{\xi}$ et $I' \to H_{y}$ qui sont les épimorphismes
canoniques, la flèche horizontale du bas étant $N$ et celle du haut $\varphi$.}

\textbf{Critère.} Pour que $N$ soit un isomorphisme, il faut et il suffit que
$\operatorname{card} C = \operatorname{card} C'$ et que l'entier ordinaire
$\det N$ soit égal à $\pm\, p^{r}$, $p$ étant l'exposant caractéristique de
$k(x)$ et de $k(y)$. Dans ce cas les sous-groupes d'inertie de $H$ en $y$ sont
exactement les conjugués des images de ceux de $G$ en $x$.

Le critère est joli et il est exact : $N$ doit être inversible sur
$\prod_{\ell \neq p}\mathbb{Z}_{\ell}$, ce qui revient à demander que $\det N$
soit inversible dans cet anneau, c'est-à-dire premier à tout $\ell \neq p$,
c'est-à-dire une puissance de $p$ au signe près. La caractéristique n'entre
donc qu'à un seul endroit, et sous cette forme.

\subsection*{5. Le cas d'un diviseur image inverse}

Le numéro 5 particularise à $V = Y - \operatorname{supp} E$ avec $E$ un
diviseur $\geqslant 0$ sur $Y$ et $U = f^{-1}(V)$, sous deux conditions : que
les anneaux locaux soient noethériens aux points concernés, et que
$D = f^{*}(E)$ soit défini et régulier. Le tapuscrit s'interrompt au milieu de
cette seconde condition ; il reprend à la page 121 avec la condition de transversalité et la conclusion sur les classes de conjugaison
d'homomorphismes $\underline{Z}_{\ell}(1) \to G$.

\pagerange{121}{122}
\section*{Une classe de conjugaison par branche (pages 121 et 122)}

Le tapuscrit commencé à la page 116 s'achève. Sous les conditions du numéro 5 --- $D =
f^{*}(E)$ régulier, $\operatorname{supp} E$ irréductible et régulier aux points
de $f(\operatorname{supp} D)$, avec la condition de transversalité --- les
images par $\varphi : G \to H$ des sous-groupes d'inertie de $G$ en un point
$x$ de $\operatorname{supp} D$ sont des sous-groupes d'inertie de $H$ relatifs
au point maximal de $\operatorname{supp} E$.

Il en résulte le corollaire pratique : pour toute représentation linéaire de
$G$ dans un module $M$ qui se factorise par $H$, \emph{les images dans
$\operatorname{Aut}(M)$ des sous-groupes d'inertie de $G$ relatifs aux divers
points de $\operatorname{supp} D$ sont conjuguées}. C'est ce qu'on veut savoir
quand on étudie une famille : la monodromie locale est la même, à conjugaison
près, tout le long du diviseur.

\subsection*{L'énoncé le plus précis}

On peut demander mieux. Pour tout $\xi$ au-dessus d'un point géométrique
$\bar{x}$ de $\operatorname{supp} D$, considérons l'homomorphisme
\[
  \underline{Z}_{\ell}(1)(\bar{x}) \longrightarrow G
\]
dont l'image est le sous-groupe d'inertie correspondant. Pour $\xi$ variable et
$\bar{x}$ fixé, il varie par composition avec un automorphisme intérieur de
$G$ : on a donc, pour $\bar{x}$ fixé, une \emph{classe de conjugaison}
d'homomorphismes.

Pour $\bar{x}$ variable dans une composante $D_{i}$ de l'ensemble des points
réguliers de $D$ de caractéristique $\neq \ell$, on passe d'un tel homomorphisme
à un autre en composant avec l'un des isomorphismes
$\underline{Z}_{\ell}(1)(\bar{x}) \to \underline{Z}_{\ell}(1)(\bar{x}')$
associés à une classe de chemins de $\bar{x}$ à $\bar{x}'$. Si cet isomorphisme
ne dépend pas du choix de la classe de chemins --- c'est-à-dire si le faisceau
$\ell$-adique de Tate $\underline{Z}_{\ell}(1)$ est \emph{trivial} sur
$D_{i}$, ce qui se voit déjà sur sa restriction au point générique --- alors on
identifie tous les $\underline{Z}_{\ell}(1)(\bar{x})$ à un même groupe
$\underline{Z}_{\ell}(1)(D_{i})$, et l'on obtient une classe de conjugaison
d'homomorphismes
\[
  \underline{Z}_{\ell}(1)(D_{i}) \longrightarrow G
\]
bien déterminée par la seule composante $D_{i}$.

Si de plus $\underline{Z}_{\ell}(1)$ est trivial sur l'ouvert $X_{\ell}$ des
points de $X$ de caractéristique $\neq \ell$ et que celui-ci est connexe, on
identifie tout à un unique $\underline{Z}_{\ell}(1)$, et l'on obtient, pour
chaque $i$, une classe de conjugaison d'homomorphismes
$\underline{Z}_{\ell}(1) \to G$. \emph{Ces classes peuvent être distinctes pour
des $i$ distincts} --- et c'est bien pour cela qu'on les a construites. La
formation en est fonctorielle relativement à un morphisme $f : X \to Y$ du type
étudié aux numéros 2 et 4, « ce que le lecteur explicitera lui-même s'il en a
besoin ».

Une note en marge de la page 121, à ne pas laisser passer : « Remords au
n\textsuperscript{o} 3, 4 ». Elle porte sur les deux numéros où la classe
d'inertie est construite, et le carton n'en dit pas davantage.

\pagerange{124}{126}
\section*{Lettre à Serre, 18 octobre 1959 : le groupe fondamental infinitésimal (pages 124 à 126)}

\subsection*{Le contexte, en trois lignes}

La lettre s'ouvre sur Tate et les courbes elliptiques, dont Grothendieck dit
n'avoir « pas du tout compris pourquoi ses résultats suggèreraient l'existence
d'une telle définition » --- il s'agit d'une définition globale des variétés
analytiques sur un corps valué complet --- et reste sceptique. Puis, sans
transition :

\begin{quote}
J'ai réfléchi un petit peu à la partie « infinitésimale » du groupe
fondamental, juste assez pour me convaincre que ça existe et est raisonnable.
\end{quote}

\subsection*{Le dispositif}

On se donne un schéma connexe $S$ (par exemple un schéma algébrique sur un
corps $k$), une catégorie auxiliaire $\mathcal{C}$ (par exemple les schémas
algébriques finis sur $k$) où les produits fibrés existent, et une catégorie
$\mathcal{G}$ dont les objets sont des groupes de $\mathcal{C}$ (par exemple
les groupes algébriques finis sur $k$), soumise à :

\begin{enumerate}
  \item[(i)] $\mathcal{G}$ est stable par produits, et par noyaux de couples de
    morphismes --- le sous-groupe maximum où deux morphismes coïncident, qui
    existe puisqu'il s'exprime par un produit fibré ;
  \item[(ii)] les images existent, sont des quotients, et sont dans
    $\mathcal{G}$ ;
  \item[(iii)] toute suite décroissante de sous-groupes appartenant à
    $\mathcal{G}$ d'un objet de $\mathcal{G}$ est stationnaire.
\end{enumerate}

Et un foncteur covariant $F$ de $\mathcal{G}$ dans les schémas en groupes sur
$S$ (par exemple $G \mapsto S \times_{k} G$), soumis à :

\begin{enumerate}
  \item[(iv)] $F$ est \emph{exact} : il commute aux produits, aux noyaux de
    couples, aux images ;
  \item[(v)] tout $H$ de la forme $F(G)$ est « spécial », c'est-à-dire admet
    une suite exacte de schémas en groupes finis et plats sur $S$
    \[
      e \longrightarrow H_{\mathrm{inf}} \longrightarrow H \longrightarrow
      H_{\mathrm{s\acute{e}p}} \longrightarrow e ,
    \]
    où $H_{\mathrm{inf}}$ est purement infinitésimal --- la projection sur $S$
    est géométriquement injective --- et $H_{\mathrm{s\acute{e}p}}$ non ramifié
    sur $S$ ;
  \item[(vi)] $S$ est réduit.
\end{enumerate}

Il ajoute, et c'est important pour lire ce qui suit : « Les conditions (v) et
(vi) ont une sale gueule, et sont essentiellement provisoires. » Elles ne
servent qu'à un lemme.

\textbf{Lemme.} Soient $H$ comme en (v), $S$ comme en (vi), $P$ et $Q$ deux
fibrés principaux homogènes sous $H$, et $u$, $v$ deux isomorphismes de $P$
dans $Q$ transformant un point marqué donné en un autre donné. Alors $u = v$.

\emph{Démonstration.} En tordant $H$, on se ramène au cas où $P$ est trivial ;
alors il suffit d'observer qu'une section de $Q$ est un isomorphisme de $S$ sur
une composante connexe de $Q_{\mathrm{r\acute{e}d}}$.

Le lemme est ce qui fait tenir toute la construction : \emph{un fibré
principal pointé n'a pas d'automorphisme pointé}. Sans cela, la classification
n'est pas un ensemble mais un groupoïde, et le foncteur ci-dessous ne serait
pas à valeurs dans les ensembles.

\subsection*{Le foncteur, et ce qu'il vérifie}

Soit $a$ un point marqué de $S$ --- une extension algébriquement close du corps
résiduel d'un $s \in S$. Pour $G \in \mathcal{G}$, on note
\[
  Z(S, a\,;\, G) \;=\; Z(G)
\]
l'ensemble des classes d'isomorphie de fibrés principaux homogènes sur $S$, de
groupe $F(G)$, munis d'un point marqué au-dessus de $a$. C'est un foncteur
covariant de $\mathcal{G}$ dans les ensembles, et il vérifie :

\begin{enumerate}
  \item[1)] il commute aux produits, puisque $F$ y commute ;
  \item[2)] il commute aux noyaux de couples : si
    $G'' \to G \rightrightarrows G'$ est exacte dans $\mathcal{G}$, alors
    $Z(G'') \to Z(G) \rightrightarrows Z(G')$ l'est, c'est-à-dire que $Z(G'')$
    s'identifie à l'ensemble des éléments de $Z(G)$ ayant même image par les
    deux flèches. Cela résulte de l'exactitude de $F$ et du lemme.
\end{enumerate}

\subsection*{La pro-représentabilité}

De ces deux propriétés il résulte \emph{formellement} qu'il existe un système
projectif filtrant $(G_{i})_{i \in I}$ dans $\mathcal{G}$, à morphismes de
transition épimorphiques, essentiellement unique, avec un isomorphisme
fonctoriel
\[
  Z(G) \;=\; \varinjlim_{i} \operatorname{Hom}_{\mathcal{G}}(G_{i}, G) .
\]

C'est ce système projectif --- pris modulo l'équivalence qui, intuitivement,
revient à passer à la limite projective des $G_{i}$ --- qui \emph{est} le
groupe fondamental. On le note $\pi_{1}^{\mathcal{G}}(S, a)$ ; et dans le cas
d'un corps de base $k$ avec $\mathcal{G}$ la catégorie des groupes algébriques
finis sur $k$, on écrira $\pi_{1}(S/k, a)$ : « c'est le groupe proalgébrique
(mais avec partie infinitésimale) fondamental du $k$-schéma $S$ ».

Il faut voir la forme du geste, parce qu'il revient dix fois dans le carton :
\emph{on ne construit pas le groupe, on construit le foncteur qu'il devrait
représenter, et on montre que ce foncteur est pro-représentable.} Les pages 164
à 185 feront la même chose pour des groupes structuraux infinis, la page 104
pour $\pi_{2}$, et la page 208 axiomatisera la catégorie où c'est
possible.\footnote{L'existence du système projectif est démontrée à la page 127
par les \emph{couples minimaux}, et c'est le seul endroit où l'argument est
écrit. Un couple $(G, z)$, $z \in Z(G)$, est dit \emph{minimal} si $z$ ne
provient d'aucun $z'$ pour un sous-groupe $G' \subsetneq G$ ; il est
\emph{majoré} par $(G', z')$ s'il existe $u : G' \to G$ avec $z = Z(u)(z')$.
La condition (iii) assure que tout couple est majoré par un couple minimal, et
la propriété 2) que le $u$ est alors \emph{unique}. Les couples minimaux
forment donc un système projectif, filtrant parce que $(G,z)$ et $(G',z')$ sont
tous deux majorés par $(G\times G', (z,z'))$. C'est exactement l'argument que
les feuillets 136 à 140 refont à la main, et deux fois.}

\subsection*{Le dévissage sur un corps parfait}

Lorsque $k$ est parfait, tout groupe algébrique fini se décompose en partie
infinitésimale et partie réduite, et le groupe fondamental est le produit
semi-direct de sa \emph{partie séparable} --- un groupe profini ordinaire sur
lequel $\operatorname{Gal}(\bar{k}/k)$ opère --- par sa \emph{partie
infinitésimale}, laquelle ne dépend d'ailleurs plus du point base.

Deux remarques qu'il faut lire ensemble. La première : la partie séparable est
« plus grande » que le groupe fondamental ordinaire, parce qu'elle classifie
les revêtements principaux dont le groupe structural est un groupe fini
\emph{et séparable sur $k$} --- c'est-à-dire un groupe fini ordinaire sur
lequel on fait opérer $\operatorname{Gal}(\bar{k}/k)$ de façon non
nécessairement triviale. La seconde : « J'avoue que si $k$ n'est pas
algébriquement clos, le groupe fondamental ci-dessus n'est sans doute pas le
bon » ; il faudrait prendre celui de $S \times_{k} \bar{k}$, muni de sa donnée
de descente, donc défini en réalité sur $k$.\footnote{Un paragraphe entier est
biffé à cet endroit, sur les points rationnels du groupe fondamental de
$\operatorname{Spec}(k)$ et le centre de $\operatorname{Gal}(\bar{k}/k)$. La
transcription le conserve, comme elle doit ; il n'est pas repris ici, l'auteur
l'ayant retiré.}

\pagerange{127}{129}
\section*{Ce que la lettre conjecture, et ce qu'elle calcule (pages 127 à 129)}

\subsection*{La suite exacte d'homotopie, à trouver}

« Je ne sais encore comment devrait être formulée la suite exacte d'homotopie ;
pour bien faire, l'inclusion dans le groupe fondamental d'une partie
infinitésimale devrait donner une théorie satisfaisante du comportement du
groupe fondamental par spécialisation. »

L'espoir est le suivant. Soit $f : X \to Y$ propre et séparable à fibres
absolument connexes --- c'est-à-dire $f_{*}(\mathcal{O}_{X}) = \mathcal{O}_{Y}$
--- muni d'une section qui détermine des points-base sur les fibres. Alors les
groupes fondamentaux \emph{totaux} des fibres devraient former sur $Y$ un
pro-schéma en groupes filtrant : il existerait un système projectif
$\pi_{1}(X/Y, s) = (G_{i})_{i \in I}$ de schémas en groupes « spéciaux » sur
$Y$, à morphismes de transition épimorphiques, dont les groupes fondamentaux
des fibres se déduiraient par simple spécialisation.

Et la conjecture proprement dite : $\pi_{1}(X/Y)$, $\pi_{1}(X)$ et l'image
inverse par $f$ de $\pi_{1}(Y)$ devraient être reliés par une suite exacte.
Il ajoute --- et c'est la seule fois du carton où il rêve à voix haute ---
« Il y a de la joie en perspective pour les groupes d'homotopie supérieurs
\dots ».

Une conséquence immédiate est notée : pour des schémas complets $X$, $Y$ sur
$k$ algébriquement clos,
\[
  \pi_{1}(X \times Y/k) \;=\; \pi_{1}(X/k) \times \pi_{1}(Y/k) .
\]
C'est l'énoncé qui figure à la page 62 --- mais pour le quotient
premier à la caractéristique, non pour le groupe total. La différence est
exactement la partie infinitésimale, et c'est elle qui rend l'énoncé de 1959
plus fort que celui de 1967.

\subsection*{Le calcul sur une variété abélienne}

C'est le morceau de la lettre.

Utilisant les raisonnements de Serre, on conclut que si $X$ est une variété
abélienne sur $k$, alors $\pi_{1}(X/k)$ est \emph{abélien}, et qu'un revêtement
principal minimal $X'$ de $X$ est une variété abélienne isogène à $X$. On en
déduit
\[
  \pi_{1}(X/k) \;=\; \varprojlim_{n} {}_{n}X ,
\]
où ${}_{n}X$ est le noyau de la multiplication par $n$ --- \emph{avec sa partie
infinitésimale, bien sûr} --- l'homomorphisme de ${}_{mn}X$ dans ${}_{n}X$
étant induit par la multiplication par $m$.

Et, par la dualité de Cartier : ce groupe fondamental est dual du groupe
ind-algébrique
\[
  \varinjlim_{n} {}_{n}X^{*} ,
\]
où $X^{*}$ est la variété duale.

Prenant la $p$-composante, « on trouve ce qui, pour le nombre premier $p$, doit
jouer le rôle du module de Weil » --- c'est-à-dire l'analogue en
caractéristique $p$ du module de Tate. Il ajoute qu'il est hors de doute qu'on
peut associer fonctoriellement à un groupe proalgébrique infinitésimal abélien
un module sur l'anneau de Witt, et que ce module se décompose en les trois
morceaux correspondant aux trois types principaux de groupes algébriques
abéliens ; ce qui donnerait au théorème de Serre « une formulation plus
naturelle », avec une démonstration uniforme sans distinction entre
$\ell \neq p$ et $\ell = p$, et montrerait que « ta somme abracadabrante se
comporte bien quand on fait varier $X$ dans une famille ».

\subsection*{Ce que c'est, aujourd'hui}

L'objet construit dans cette lettre --- le pro-schéma en groupes finis qui
classifie les torseurs pointés sous les schémas en groupes finis, avec sa
partie infinitésimale --- est ce qu'on appelle aujourd'hui le \emph{schéma en
groupes fondamental}. La construction publiée, par la voie tannakienne des
fibrés essentiellement finis, est due à Madhav Nori (thèse de 1976, article de
1982), qui démontre en particulier que pour une variété abélienne
$\pi^{N}(X) = \varprojlim_{n} {}_{n}X$, et la dualité de Cartier avec
$\varinjlim {}_{n}X^{*}$.\footnote{Les deux constructions ne sont pas la même :
celle de Nori est tannakienne et suppose $X$ propre réduit connexe avec un
point rationnel ; celle de la lettre est par les torseurs et repose sur les
conditions (v) et (vi), dont Grothendieck écrit lui-même qu'elles « ont une
sale gueule, et sont essentiellement provisoires ». La lettre est une esquisse
avec des hypothèses avouées comme provisoires, non une théorie. Ce que la
présente note constate est que l'objet, la méthode de pro-représentabilité, le
dévissage semi-direct sur un corps parfait, le calcul sur une variété abélienne
et la dualité de Cartier sont dans une lettre du 18 octobre 1959 restée
inédite. Elle ne constate rien de plus : aucune filiation n'est ici affirmée,
et la question de savoir ce que la littérature publiée établit à ce sujet n'est
pas de la compétence de cette édition.}

\subsection*{Le calendrier}

La lettre se termine ainsi, et il faut la lire jusqu'au bout :

\begin{quote}
J'espère arriver dans l'année prochaine à une théorie satisfaisante du groupe
fondamental, et achever la rédaction des Chapitres IV, V, VI, VII (ce dernier
étant le groupe fondamental), en même temps que des catégories. Dans deux ans
résidus, dualité, intersections, Chern, Riemann-Roch. Dans trois ans
cohomologie de Weil, et un peu d'homotopie si Dieu veut. Et entre-temps, je ne
sais quand, le grand théorème d'existence avec Picard etc, un peu de courbes
algébriques, les schémas abéliens. Sauf difficultés imprévues ou enlisement, le
multiplodoque devrait être fini d'ici 3 ans, ou 4 au maximum. On pourra
commencer à faire de la géométrie algébrique !
\end{quote}

\pagerange{130}{134}
\section*{Les « feuilles anciennes » : représentabilité et recollement (pages 130 à 134)}

Onze feuillets manuscrits que la chemise appelle « \emph{Vrai} groupe
fondamental [feuilles anciennes] » reprennent la même théorie de sa main. Ils
sont d'une écriture pâle et lourdement raturée, et la transcription y laisse
beaucoup d'illisible ; ce qui suit est ce qui s'énonce.

\subsection*{Le foncteur des fibrés pointés}

Soit $f : X \to S$ propre et plat, muni d'une section $\sigma$, avec
$f_{*}(\mathcal{O}_{X}) = \mathcal{O}_{S}$ \emph{universellement}. Soit $G$ un
groupe fini et étale sur $S$ --- le cas important étant $G = S \times \Gamma$
pour $\Gamma$ un groupe fini ordinaire. On note $F_{G}(T)$ l'ensemble des classes d'isomorphie de fibrés principaux
homogènes sur $X_{T}$, de groupe $G_{X_{T}}$, munis d'une section le long de
$\sigma_{T}$ --- c'est-à-dire d'un isomorphisme
$\sigma_{T}^{*}(X') \simeq G_{T}$.

\textbf{Question.} Le foncteur $T \mapsto F_{G}(T)$ est-il représentable ? Si
oui, le schéma qui le représente est \emph{étale} sur $S$. On note que
$F_{G}$ est compatible aux limites et à la descente fidèlement plate
quasi-compacte.

C'est le même foncteur que le $Z(S,a;G)$ de la lettre, rendu relatif : la
lettre le prenait sur un $S$ fixe et cherchait à le pro-représenter ; ici on le
prend en famille et l'on cherche à le représenter. La différence n'est pas de
degré : la seconde question est celle de la \emph{constructibilité en famille}
que la page 44 appelait de ses vœux.

\subsection*{Un critère de représentabilité}

\textbf{Critère.} Soit $S$ un préschéma localement noethérien et
$F : (\mathrm{Sch})_{/S} \to (\mathrm{Ens})$ un foncteur vérifiant :

\begin{enumerate}
  \item[(i)] $F$ est compatible aux limites inductives d'anneaux ;
  \item[(ii)] $F$ est compatible à la descente fidèlement plate quasi-compacte
    et à la localisation ;
  \item[(iii)] si $A$ est un anneau local noethérien \emph{complet} sur $S$ de
    corps résiduel $k$, alors $F(A) \simeq F(k)$ ;
  \item[(iv)] $F$ est séparé ;
  \item[(v)] pour tout corps $k$ sur $S$, $F(k)$ est fini, et
    $F(k) \to F(k')$ est bijectif pour $k'$ une extension algébrique séparable
    de $k$.\footnote{Le qualificatif de $k'$ porte sur la page une mention
    biffée que la transcription ne lit pas ; « alg.\ sép.\ » et « alg.\ clos »
    y voisinent. La condition est donnée ici telle que la page la laisse lire,
    et elle n'est pas complétée par conjecture.}
\end{enumerate}

Alors $F$ est représentable par un schéma \emph{étale, séparé et quasi-fini}
sur $S$.

C'est la condition (iii) qui porte le poids : elle dit que le foncteur ne voit
que la fibre spéciale d'un anneau local complet, ce qui est la version
foncteur du fait qu'un revêtement étale d'un local complet est déterminé par sa
fibre spéciale. La démonstration est esquissée puis s'arrête au milieu d'une
phrase.

\subsection*{Le lemme de recollement}

\textbf{Lemme.} Soit $S$ localement noethérien, et pour tout $s \in S$ un
préschéma $X(s)$ quasi-fini étale séparé sur
$\operatorname{Spec}(\mathcal{O}_{s})$, avec, pour $s \in \overline{\{t\}}$,
des morphismes $\varphi_{st} : X(t) \to X(s)$ au-dessus de
$\operatorname{Spec}(\mathcal{O}_{t}) \to \operatorname{Spec}(\mathcal{O}_{s})$
vérifiant $\varphi_{st}\varphi_{tu} = \varphi_{su}$. Alors il existe un
préschéma $X$ quasi-fini séparé étale sur $S$ et des isomorphismes
\[
  \varphi_{s} : X \times_{S} \operatorname{Spec}(\mathcal{O}_{s})
  \xrightarrow{\ \sim\ } X(s)
\]
compatibles, et $X$ est unique à isomorphisme unique près.

\begin{tikzcd}
X(s) \arrow[d] & X(t) \arrow[l, "\varphi_{st}"'] \arrow[d] \\
\operatorname{Spec}(\mathcal{O}_{s}) & \operatorname{Spec}(\mathcal{O}_{t}) \arrow[l]
\end{tikzcd}

L'unicité résulte d'un énoncé plus général de pleine fidélité sur les
préschémas de type fini sur $S$. Sur l'existence, Grothendieck écrit : « Je ne
suis pas [sûr] de mon raisonnement » --- et la marge propose la stratégie qui
convient : montrer que pour tout $s$, $X(s)$ se réalise sur un voisinage ouvert
$U$ de $s$, puis recoller.

\pagerange{136}{140}
\section*{Groupes admissibles, couples minimaux, et la proposition 4 (pages 136 à 140)}

\subsection*{Les groupes admissibles}

Soit $S$ un schéma \emph{réduit connexe}. Appelons \emph{admissible} un schéma
en groupes $G$, fini et plat sur $S$, admettant une suite exacte
\[
  e \longrightarrow U \longrightarrow G \longrightarrow \Gamma \longrightarrow
  e
\]
avec $U \to S$ plat, fini et géométriquement injectif, et $\Gamma \to S$ plat,
fini et séparable. Alors le sous-groupe $U$ est \emph{déterminé de façon
unique} --- avec, en marge, la question : « pour cette question, l'hypothèse
« $S$ réduit » est-elle nécessaire ? ».

C'est la condition (v) de la lettre, débarrassée du foncteur $F$ : un groupe
est admissible s'il se dévisse en une partie infinitésimale et une partie
étale.

\textbf{Proposition 1.} Soit $S' \to S$ un morphisme et $G$ un $S$-groupe fini.
Si $G$ est admissible, $G \times_{S} S'$ l'est --- c'est trivial, la suite
exacte se changeant de base. Réciproquement, si $G \times_{S} S'$ est admissible
et si $S' \to S$ est fidèlement plat, alors $G$ est admissible, par descente et
unicité de $U$.\footnote{La page énonce les deux directions dans l'ordre
inverse : « Si $G \times_{S} S'$ est admissible, $G$ aussi ; la réciproque est
vraie si $S' \to S$ est fidèlement plat », avec « Direct est trivial » en
dessous. C'est le changement de base qui est trivial et la descente qui demande
la platitude fidèle ; les deux directions sont remises ici dans l'ordre que
l'argument impose. Ce qui est écrit à la suite --- « Réciproque résulte du
théorème de \dots\ et unicité » --- confirme la lecture : c'est bien l'unicité
de $U$ qui sert à descendre.}

\textbf{Corollaire.} Si $G$ est admissible et $P$ un fibré principal sous $G$,
alors $\operatorname{Aut}(P)$, comme $S$-groupe, est admissible.

\textbf{Proposition 2.} Si $(P, \xi)$ est un fibré principal \emph{pointé} sous
un $S$-groupe admissible $G$, alors $\operatorname{Aut}(P, \xi)$ est réduit à
l'identité.

\textbf{Proposition 3.} Deux sections d'un préschéma réduit connexe $P$ sur $S$
qui coïncident en un point sont identiques. Plus précisément : si $s$ est une
section de $P/S$, il existe un sous-préschéma fermé $Q$ de $P$ tel que
$Q_{\mathrm{r\acute{e}d}} \to S$ soit un isomorphisme, d'isomorphisme
réciproque induit par $s$.

Les propositions 2 et 3 sont le lemme de la lettre, découpé en deux : la
seconde est le fait géométrique, la première la conséquence dont on se sert.

\subsection*{Minimalité}

Un couple $(P, \xi)$, relatif à un $G$ admissible, est dit \emph{minimal} s'il
n'existe pas de sous-groupe admissible $G' \subset G$ et de $(P', \xi')$
relatif à $G'$ tels que
\[
  (P, \xi) \;\simeq\; (P', \xi') \times_{G'} G .
\]
Autrement dit : $(P,\xi)$ n'est pas obtenu par \emph{extension du groupe
structural} depuis un sous-groupe strict.\footnote{La page écrit
$(P',\xi') \times_{G} G'$, avec les deux groupes échangés. L'extension du
groupe structural d'un $G'$-fibré à un $G$-fibré est le produit contracté
$P' \times_{G'} G$ : c'est le groupe \emph{du fibré} qui va en indice. La
transcription signale la lecture et la rapproche du $(P,\xi) \times_{G}
(G' \times G')$ des pages 138 et 140, où l'indice est bien le groupe du fibré
--- ce qui confirme que la page 137 est un lapsus et non une autre convention.}

On pose alors
\[
  \mathfrak{Z}(S, a\,;\, G) \;=\; \mathfrak{Z}(G) \;=\;
  \bigl\{\text{classes d'isomorphie de $G$-fibrés principaux pointés au-dessus
  de } a\bigr\} ,
\]
foncteur covariant en le $S$-groupe admissible $G$. \emph{C'est le $Z(S,a;G)$
de la lettre}, écrit d'une autre lettre, et c'est aussi le
$\pi^{1}(S,\xi\,;\,G)$ des pages 164 et suivantes : trois notations, un seul
foncteur, et c'est la lecture du carton entier qui permet de le dire.

\subsection*{La proposition 4, démontrée deux fois}

\textbf{Proposition 4.} Soit $\alpha \in \mathfrak{Z}(G)$ \emph{minimal}. Alors
\[
  u \longmapsto u_{*}(\alpha) , \qquad
  \operatorname{Hom}^{\mathrm{gr}}_{S}(G, G') \longrightarrow \mathfrak{Z}(G')
\]
est \emph{injective}.

\emph{Démonstration.} Soient $u, v : G \to G'$ avec $u_{*}(\alpha) =
v_{*}(\alpha)$. Soit $(P, \xi)$ de classe $\alpha$ et $(P', \xi')$ de classe
$u_{*}(\alpha) = v_{*}(\alpha)$. Le couple $(u,v) : G \to G' \times G'$ donne
\[
  (Q, \eta) \;=\; (P, \xi) \times_{G} (G' \times G')
  \;=\; \bigl[(P,\xi) \times_{u} G'\bigr] \times
        \bigl[(P,\xi) \times_{v} G'\bigr] ,
\]
et par hypothèse il existe un unique homomorphisme $(P', \xi') \to (Q, \eta)$
compatible avec la diagonale $G' \to G' \times G'$.

\begin{tikzcd}
G \arrow[r, "{(u,v)}"] & G' \times G' \\
H \arrow[u] \arrow[r] & G' \arrow[u]
\end{tikzcd}
\qquad
\begin{tikzcd}
(P, \xi) \arrow[r] & (Q, \eta) \\
(R, \lambda) \arrow[u] \arrow[r] & (P', \xi') \arrow[u]
\end{tikzcd}

Soit $(R, \lambda)$ le produit fibré de $(P, \xi)$ et $(P', \xi')$ au-dessus de
$(Q, \eta)$, et $H$ le noyau du couple $G \rightrightarrows G'$ ---
c'est-à-dire le sous-groupe où $u$ et $v$ coïncident. Alors $(R, \lambda)$ est
un fibré principal \emph{de groupe $H$}, et la minimalité de $(P, \xi)$ force
$H = G$, donc $u = v$.

C'est l'énoncé sur lequel toute la pro-représentabilité repose : c'est lui qui
rend \emph{unique} le morphisme entre deux couples minimaux, donc qui fait des
couples minimaux un système projectif et non un groupoïde. La lettre l'utilise
page 127 sans le démontrer ; ces feuillets le démontrent, et deux fois --- la
page 139 porte une première tentative qui s'arrête sur la phrase « Soit
$(R,\lambda)$ le produit fibré de $(P,\xi)$ et $(P',\xi')$ sur $(Q,\eta)$ », et
la page 140 la reprend depuis le début et la conduit à
terme.\footnote{L'abandon de la page 139 est conservé dans la transcription
pour cette raison : il marque l'endroit où l'argument a été refondé. La reprise
de la page 140 change le dispositif --- elle introduit d'emblée la catégorie
$\mathcal{C}$ de $S$-groupes finis et plats stable par produits, noyaux de
couples et images, c'est-à-dire les conditions (i) et (ii) de la lettre --- et
c'est ce cadre-là qui permet de conclure.}

La fin de la démonstration --- l'extension de base $(S_{1}, a_{1}) \to (S, a)$
et l'identification de $P_{1}$ à $G'_{1}$ --- est aux pages 141 et 142, et c'est de là que sort le pro-objet
$\pi_{1}^{\mathcal{C}}(S, \xi)$.

\pagerange{141}{142}
\section*{Le pro-objet $\pi_{1}^{\mathcal{C}}(S,\xi)$ (pages 141 et 142)}

\subsection*{La fin de la proposition 4}

La démonstration commencée page 140 s'achève par extension de la base. On se
donne $(S_{1}, a_{1}) \to (S, a)$ ; la situation devient celle de $G_{1}$,
$G'_{1}$, $(P_{1}, \xi_{1})$, $(Q_{1}, \eta_{1})$, et le fibré $P_{1}$
s'identifie alors à $G'_{1}$ muni de son marquage $\varepsilon'_{1}$, tandis que
$Q_{1}$ s'identifie à $G'_{1} \times_{S_{1}} G'_{1}$ muni de
$(\varepsilon'_{1}, \varepsilon'_{1})$. Il suffit alors de voir que
$(P'_{1}, \xi'_{1}) \to (Q_{1}, \eta_{1})$ est l'inclusion diagonale
$G'_{1} \to G'_{1} \times_{S_{1}} G'_{1}$.

Notant $w$ l'isomorphisme $(P, \xi) \times_{u} G' \to (P, \xi) \times_{v} G'$
et $w_{1}$ son extension à $S_{1}$, l'assertion revient à dire que $w_{1}$
transporte la section naturelle de $P_{1} \times_{u_{1}} G'_{1}$ sur celle de
$P_{1} \times_{v_{1}} G'_{1}$ ; et cela résulte de l'unicité de
$P \to P \times_{u} G'$ compatible au marquage.

\begin{tikzcd}[column sep=large]
P \arrow[r, "\mathrm{can}"] \arrow[rr, bend right=30, "\mathrm{can}"'] & P \times_u G' \arrow[r, "w"] & P \times_v G'
\end{tikzcd}

\subsection*{Le pro-objet}

Restent deux propriétés formelles. D'abord la multiplicativité,
\[
  \mathfrak{Z}(G \times G') \;=\; \mathfrak{Z}(G) \times \mathfrak{Z}(G') ,
\]
et ensuite la condition minimale sur les sous-groupes appartenant à
$\mathcal{C}$ d'un objet de $\mathcal{C}$.\footnote{La phrase de la page ---
« Les s-gpes $G \in \mathcal{C}$ d'un $G \in \mathcal{C}$ satisfaisant la
condition minimale » --- fait servir la même lettre deux fois, pour le
sous-groupe et pour l'objet. La transcription le signale et le conserve ; ici
les deux sont séparés, et il s'agit bien de la condition (iii) de la lettre à
Serre : la chaîne descendante des sous-groupes est stationnaire.}

\emph{Ces propriétés formelles suffisent} pour construire un pro-objet de
$\mathcal{C}$, au moyen d'un système projectif d'épimorphismes, que l'on note
\[
  \pi_{1}^{\mathcal{C}}(S, \xi) .
\]

C'est le même objet que la lettre de 1959 construit page 126, obtenu ici sans
passer par les couples minimaux, mais par les seules propriétés du foncteur.
Deux exemples sont amorcés et laissés inachevés : $\mathcal{C}$ une classe de
groupes définis par un groupe abstrait --- « on retrouve le groupe
fondamental » --- et $\mathcal{C}$ provenant d'un corps
$k$.\footnote{Les deux pages 141 et 142 sont au crayon très pâle et sont les
plus difficiles du carton ; la transcription y laisse beaucoup d'illisible et
cette lecture ne le comble pas. Ce qui est restitué ici sont les énoncés
portants, non les phrases qui les joignent.}

\pagerange{144}{146}
\section*{Données de descente : la relation de cocycle (pages 144 à 146)}

Deux feuillets écrivent la condition de cocycle pour une donnée de descente le
long d'un morphisme, dans le langage des groupoïdes fondamentaux.

On dispose d'un $S' \to S$ et de ses puissances fibrées, avec les projections
$p_{i}^{\alpha}$ de $\pi''_{\alpha}$ dans $\pi'$ et $p_{ij}^{\alpha\beta}$ de
$\pi'''_{\beta}$ dans $\pi''_{\alpha}$, indexées par les trois couples
$(1,2)$, $(2,3)$, $(1,3)$. Les relations entre elles font intervenir des
automorphismes intérieurs, et le carré qui les résume est celui-ci :

\begin{tikzcd}[column sep=large]
E \arrow[r, "A_{\alpha'}"] \arrow[d, "a'_{\beta}"'] & E \arrow[r, "a''_{\beta}", "\sim"'] & E \arrow[d, "A_{\alpha''}"] \\
E \arrow[r, "A_{\alpha'''}"'] & E \arrow[r, "a'''_{\beta}"', "\sim"] & E
\end{tikzcd}

d'où la relation, que la page encadre :
\[
  \boxed{\; A_{\alpha''}\, a''_{\beta}\, A_{\alpha'} \;=\;
  a'''_{\beta}\, A_{\alpha'''}\, a'_{\beta} \;}
\]

C'est la condition de cocycle, écrite comme la commutation d'un rectangle. La
page 146 la reprend en termes de \emph{chemins} --- $g_{s''} = g_{s'''}g_{s'}$,
d'où $g = 1$ pour des chemins triviaux --- puis abandonne : elle est barrée
presque entièrement. Ce qui reste au bas de la feuille est le diagramme
simplicial $S \leftarrow S' \leftarrow S'' \leftarrow S'''$ avec les
$R_{i}$ au-dessus, et trois mots isolés.

Cette relation reviendra sous une forme réglée à la page 180, où les données de
descente sur un faisceau constant sont décrites comme des applications
$f : K_{1} \to G$ vérifiant $f \circ p_{1} = (f \circ p_{0})(f \circ p_{2})$.
C'est la même identité, et c'est là seulement qu'elle est écrite proprement.

\pagerange{148}{150}
\section*{« Compléments aux notes IHES » : le $\pi_{1}$ d'un produit fibré (pages 148 à 150)}

\subsection*{L'énoncé}

Soient $X$ et $Y$ deux schémas sur $S$, avec $X$ propre, séparable et surjectif
sur $S$, et $Y$ non vide. Soit $c$ un point géométrique de $X \times_{S} Y$,
d'images $a$, $b$, $d$ dans $X$, $Y$, $S$. On dispose de l'homomorphisme
canonique
\[
  \pi_{1}(X \times_{S} Y,\, c) \longrightarrow
  \pi_{1}(X, a) \times_{\pi_{1}(S, d)} \pi_{1}(Y, b) .
\]

\textbf{Il est toujours surjectif.} Il est \emph{bijectif} dès que la condition
suivante est vérifiée : il existe un $S'$ non vide, étale sur $S$, tel que
$Y' = Y \times_{S} S'$ admette une section sur $S'$ --- c'est-à-dire, ce qui
revient au même, tel qu'il existe un $S$-morphisme $S' \to Y$ --- et tel que
$\pi_{1}(X') \to \pi_{1}(X)$ soit injectif.

Le cas particulier qui sert : lorsque $S = \operatorname{Spec} k$, on obtient
\[
  \pi_{1}(X \times_{k} Y) \;\simeq\; \pi_{1}(X) \times_{\pi_{1}(k)}
  \pi_{1}(Y)
\]
pour $X$ propre universellement connexe sur $k$ et $Y$ simplement connexe sur
$k$, sans plus.

Il faut mesurer ce que la condition coûte. Elle demande un point rationnel de
$Y$, au moins après une extension étale de la base --- ce qui n'est pas
automatique --- et l'injectivité de $\pi_{1}(X') \to \pi_{1}(X)$. Sans elle,
seule la surjectivité subsiste, et la différence est réelle : c'est elle qui
sépare le théorème de Künneth pour $\pi_{1}$ de la simple existence d'un
homomorphisme.

\subsection*{La démonstration}

Elle est page 150, et c'est une échelle de suites exactes. Les deux noyaux
\[
  \operatorname{Ker}\bigl[\pi_{1}(X) \to \pi_{1}(S)\bigr]
  \qquad\text{et}\qquad
  \operatorname{Ker}\bigl[\pi_{1}(X \times_{S} Y) \to \pi_{1}(Y)\bigr]
\]
s'identifient tous deux à l'image de $\pi_{1}(\bar{X}_{s})$ --- le groupe
fondamental de la fibre géométrique --- et la comparaison des deux lignes donne
le résultat.

\begin{tikzcd}[column sep=small, row sep=small, nodes={font=\scriptsize}]
& \operatorname{Im}(\pi_1(\bar{X}_s) \to \pi_1(X)) \arrow[d, no head] & & & \\
e \arrow[r] & \operatorname{Ker}[\pi_1(X) \to \pi_1(S)] \arrow[r, "\sim"] & \pi_1(X) \times_{\pi_1(S)} \pi_1(Y) \arrow[r] & \pi_1(Y) \arrow[r] & e \\
e \arrow[r] & \operatorname{Ker}[\pi_1(X \times_S Y) \to \pi_1(Y)] \arrow[u] \arrow[d, no head] \arrow[r] & \pi_1(X \times_S Y) \arrow[u] \arrow[r] & \pi_1(Y) \arrow[u] \arrow[r] & e \\
& \operatorname{Im}(\pi_1(\bar{X}_s) \to \pi_1(X \times_S Y)) & & &
\end{tikzcd}

C'est la propreté de $X$ sur $S$ qui fait marcher les deux identifications :
elle donne la suite exacte d'homotopie, dont les deux lignes sont deux
instances.

La page 149, entre les deux, ne porte qu'un croquis : $X \to S$ avec $S'$ à
côté. C'est le changement de base dont les deux autres pages discutent, et
c'est tout ce que la feuille contient.

\pagerange{151}{151}
\section*{Le topos des $\pi$-ensembles (page 151)}

Une page entière barrée de trois traits diagonaux, et elle mérite d'être lue
malgré la biffure : elle vérifie que la catégorie des ensembles finis à
opérateurs est un topos et que sa cohomologie est celle du groupe.

Soit $\pi$ un groupe et $\mathcal{E}^{\pi}$ la catégorie des ensembles finis où
$\pi$ opère à gauche, l'objet final étant le point $e$, l'objet initial
$\varnothing$. Le $\pi$-ensemble $\pi_{s} = \pi$ (translations à droite) a pour
groupe d'automorphismes $\pi$ lui-même, par $x \mapsto xg^{-1}$.

Un faisceau d'ensembles sur $\mathcal{E}^{\pi}$ est alors la même chose qu'un
ensemble fini $E$ où $\pi$ opère, via
\[
  E \;=\; F(\pi_{s}) , \qquad
  F(U) \;=\; \operatorname{Hom}_{\pi}(U, E) .
\]
Plus précisément, si $U = \coprod_{i} \pi/\pi_{i}$ est la décomposition en
orbites, alors la catégorie fibrée $\mathcal{F}^{\pi}$ vérifie
$\mathcal{F}_{U} \simeq \prod_{i} E^{\pi_{i}}$ ; et
\[
  H^{0}(F) = \Gamma(F) = E^{\pi} , \qquad
  H^{1}(F) \simeq H^{1}(\pi, E) , \qquad
  H^{i}(F) \simeq H^{i}(\pi, E) ,
\]
la dernière ligne valant pour $F$ et $E$ commutatifs.

C'est le \emph{topos classifiant} $B\pi$, et l'identification de sa cohomologie
à celle du groupe. La page 209 s'en servira sans le redémontrer : « prenant
pour $\mathcal{T}$ le topos $(B\pi_{i})$ ».

\pagerange{152}{156}
\section*{Un morphisme fini étale de rang $n$ est un $P/H$ (pages 152 à 156)}

\subsection*{La proposition 3.2}

\textbf{Proposition.} Soit $f : X \to Y$ un morphisme fini. Les conditions
suivantes sont équivalentes.

\begin{enumerate}
  \item[(i)] $f$ est étale, et $X$ est de rang localement égal à $n$ sur $Y$.
  \item[(ii)] Il existe $Y_{1} \to Y$ fidèlement plat tel que
    $X \times_{Y} Y_{1}$ soit $Y_{1}$-isomorphe à $Y_{1} \times I_{n}$, où
    $I_{n}$ est un ensemble à $n$ éléments.
  \item[(ii bis)] La même chose avec $Y_{1}$ un revêtement galoisien, de groupe
    le groupe symétrique $\mathfrak{S}_{n}$.
  \item[(iii)] $X$ est isomorphe à un $P/H$, où $P$ est un revêtement étale
    galoisien de $Y$ de groupe $G$ et $H$ un sous-groupe de $G$.
  \item[(iii bis)] La même chose avec $G = \mathfrak{S}_{n}$ et $H$ le
    stabilisateur d'un point de $I_{n}$.
\end{enumerate}

Entre (ii) et (ii bis), et entre (iii) et (iii bis), la page écrit
« kif -- kif » : c'est la même chose.

\subsection*{La démonstration de Serre}

L'implication (i) $\Rightarrow$ (iii bis) est créditée sur la page : « la
démonstration qui suit [est] due à J.-P.\ Serre ». Elle vaut d'être écrite en
entier, car elle est le mécanisme de toute la théorie.

Soit $X^{n} = X \times_{Y} \cdots \times_{Y} X$ le produit fibré de $n$
copies, sur lequel $\mathfrak{S}_{n}$ opère par permutation des facteurs. Soit
$P \subset X^{n}$ le complémentaire du lieu où deux coordonnées coïncident ---
c'est-à-dire l'ensemble des $n$-uplets de points \emph{deux à deux distincts}
au-dessus d'un même point de $Y$. C'est une partie \emph{ouverte et fermée} de
$X^{n}$, parce que le lieu de coïncidence en est une (la diagonale d'un
morphisme étale est ouverte et fermée) ; elle est stable sous
$\mathfrak{S}_{n}$, qui y opère \emph{sans point fixe} par construction. Enfin,
$X$ étant de rang $n$ sur $Y$, la fibre géométrique de $P$ au-dessus d'un point
de $Y$ est exactement l'ensemble des $n!$ bijections de $I_{n}$ sur la fibre de
$X$.

Donc $P \to Y$ est un revêtement étale galoisien de groupe
$\mathfrak{S}_{n}$ --- la page en fait la conclusion d'une proposition 2.5
qu'elle cite sans l'énoncer et qui n'est pas dans ce carton.\footnote{« En
vertu de \emph{prop.\ 2.5}, $P \to Y$ [est] un revêtement étale galoisien de
$Y$ de groupe $G_n$ » (page 153). Le renvoi est à un énoncé antérieur des
mêmes notes, absent de la cote ; il est donné ici pour que le lecteur sache
que cette étape est citée et non démontrée sur place.} Et la dernière
projection $X^{n} \to X$ est invariante sous
le stabilisateur $H = \mathfrak{S}_{n-1}$ du dernier indice, d'où
$X \simeq P/H$.\footnote{C'est la construction dite du \emph{revêtement des
repères} : $P$ est le schéma des isomorphismes $I_{n} \to X$ au-dessus de $Y$,
et le fait qu'il soit galoisien de groupe $\mathfrak{S}_{n}$ est ce qui permet
de ramener tout revêtement fini étale à un revêtement galoisien. C'est l'énoncé
qui rend la théorie de Galois abstraite non vide : sans lui, l'existence d'un
revêtement galoisien majorant un revêtement donné serait à démontrer cas par
cas.}

Les implications restantes se ferment : (iii bis) $\Rightarrow$ (iii) est
triviale ; (iii) $\Rightarrow$ (i) résulte d'un lemme antérieur ; (iii)
$\Rightarrow$ (ii bis) en prenant $Y_{1} = P$ et en utilisant
$(X \times G)/H \simeq X \times G/H$ ; et (ii bis) $\Rightarrow$ (i) est
immédiat.

\subsection*{La proposition 3.1 : l'existence de $X/G$}

\textbf{Proposition.} Soit $X$ étale sur $Y$ et quasi-compact, et $G$ un groupe
fini opérant par $Y$-automorphismes. Alors $X/G$ existe, $X \to X/G$ est fini,
et $X/G \to Y$ est \emph{étale}.

\emph{Démonstration.} Soit $I$ l'ensemble des composantes connexes de $X$, sur
lequel $G$ opère ; en regroupant les composantes par orbite, $X$ est la somme
des $X_{(\alpha)} = \bigcup_{i \in \alpha} X_{i}$, stables sous $G$, et
$X/G = \bigsqcup_{\alpha} X_{(\alpha)}/G$. On se ramène donc au cas où $G$
opère transitivement sur les composantes connexes. Soit $X_{0}$ l'une d'elles
et $G_{0}$ son stabilisateur : alors $X \simeq X_{0} \times_{G_{0}} G$ et
$X/G \simeq X_{0}/G_{0}$.

On peut alors supposer que $G_{0}$ opère fidèlement ; mais il opère aussi
\emph{sans inertie}, et il reste à voir que $X/G$ est étale sur $Y$. Faisant
une extension fidèlement plate de la base, on se ramène au cas où l'extension
résiduelle $k(x)/k(y)$ est triviale, et l'on conclut sur les complétés :
\[
  \hat{\mathcal{O}}_{y} \longrightarrow \hat{\mathcal{O}}_{x'}
  \xrightarrow{\ \sim\ } \hat{\mathcal{O}}_{x} .
\]

\begin{tikzcd}[column sep=small, row sep=small]
X \arrow[dr] & \\
& X' = X/G \arrow[d] \\
& Y \ni y
\end{tikzcd}
\qquad
\begin{tikzcd}[row sep=small]
k(x) \arrow[d, no head] \\
k(x') \arrow[d, no head] \\
k(y)
\end{tikzcd}

La tour des corps résiduels est dessinée ici dans l'ordre des inclusions,
$k(y) \subset k(x') \subset k(x)$.\footnote{La page la dessine dans l'ordre
inverse --- $k(x')$ en haut, puis $k(x)$, puis $k(y)$ --- ce que la
transcription signale. C'est un lapsus de tracé, et non une convention : la
suite de l'argument, qui passe de $\hat{\mathcal{O}}_{y}$ à
$\hat{\mathcal{O}}_{x'}$ puis à $\hat{\mathcal{O}}_{x}$, impose l'ordre écrit
ici.}

\textbf{Corollaire.} Si $X$ est étale sur $Y$ en $x$ et si $g : X \to X'$ est
un $Y$-morphisme étale en $x$, alors $X'$ est étale sur $Y$ en
$g(x)$.\footnote{Les pages 155 et 156, au crayon sur une transparence de
verso, sont les plus difficiles de cette suite après la 141 ; la transcription y laisse
en particulier la condition (i) de la proposition 3.1 largement illisible. Ce
qui est énoncé ci-dessus est ce que les phrases portantes autorisent, et
l'hypothèse « quasi-affine » que la page ajoute en cours de démonstration est
celle sous laquelle le quotient existe.}

\pagerange{157}{158}
\section*{« Catégorie galoisienne » : Gal 1 à Gal 6 (pages 157 et 158)}

Ce sont les deux pages les plus citables de cette suite. Elles écrivent la liste
d'axiomes, et l'on y voit la liste se faire : un axiome est écrit puis biffé en
entier, les numéros sont réattribués, et le mot « indécomposable » est remplacé
interlinéairement par « connexe » à chacune de ses occurrences.

\subsection*{Les données}

\begin{enumerate}
  \item[(i)] Une catégorie $\mathcal{C}$.
  \item[(ii)] Un foncteur $F : \mathcal{C} \to (\text{ensembles finis})$.
\end{enumerate}

\subsection*{Les axiomes}

\begin{itemize}
  \item[\textbf{Gal 1}] Dans $\mathcal{C}$, les sommes et les produits finis
    existent --- d'où l'existence de l'objet initial $0$ et de l'objet final
    $1$.

  \item[\textbf{Gal 4}] $F$ commute aux sommes et aux produits.
\end{itemize}

On dit alors qu'un objet $X$ est \emph{connexe} si $X = X' \sqcup X''$ implique
$X' \simeq 0$ ou $X'' \simeq 0$.

\begin{itemize}
  \item[\textbf{Gal 2}] Tout $X$ s'écrit $X \simeq \coprod_{i \in I} X_{i}$
    avec les $X_{i}$ connexes non nuls, et toute décomposition
    $X \simeq X' \sqcup X''$ provient d'une partition $I = I' \sqcup I''$. Les
    $X_{i}$, déterminés de façon essentiellement unique avec leurs flèches vers
    $X$, sont les \emph{composantes connexes} de $X$. On a $X = 0$ si et
    seulement si $I = \varnothing$.

  \item[\textbf{Gal 3}] Soit $f : X \to Y$ avec $Y$ connexe non nul. Alors les
    composantes $X_{i}$ de $X$ telles que $f$ induise un \emph{isomorphisme}
    $f_{i} : X_{i} \xrightarrow{\ \sim\ } Y$ sont en correspondance
    biunivoque avec les \emph{sections} $g : Y \to X$ de $f$, la section
    associée à $X_{i}$ étant le composé
    $Y \xrightarrow{f_{i}^{-1}} X_{i} \hookrightarrow X$. Il suffit de le
    vérifier pour $X = Z \times Y$ et $f = p_{2}$.

  \item[\textbf{Gal 5}] Si $f : X \to Y$ induit une bijection
    $F(X) \xrightarrow{\ \sim\ } F(Y)$, c'est un isomorphisme. Cela implique en
    particulier que $F(X) = \varnothing$ entraîne $X = 0$. Il suffit de le
    vérifier quand $X$ est une composante connexe d'un produit $Z \times Y$ et
    $f$ induit par $p_{2}$.

  \item[\textbf{Gal 6}] Soit $X$ connexe non nul et \emph{galoisien},
    c'est-à-dire tel que $G = \operatorname{Aut}(X)$ opère
    \emph{transitivement} dans $F(X)$. Alors pour tout sous-groupe $H$ de $G$,
    le quotient $X/H$ existe dans $\mathcal{C}$ et
    \[
      F(X/H) \;\simeq\; F(X)/H .
    \]
\end{itemize}

\subsection*{Ce qu'on voit dans la fabrication}

Un axiome antérieur --- « $F(X) = \varnothing \Rightarrow X$ est une unité à
gauche » --- est biffé en entier, et pour cause : il est \emph{conséquence} de
Gal 5, comme la page le note aussitôt entre crochets. C'est un axiome retiré
parce que redondant, et la même chose se reproduira page 211 avec l'axiome
\textbf{G5}, dont il écrira qu'il le soupçonne redondant sans le
démontrer.\footnote{Cette liste n'est pas celle de SGA 1, exposé V, § 4, dont
la numérotation et le contenu diffèrent --- là, G3 est la factorisation d'un
morphisme en un épimorphisme suivi d'un isomorphisme sur un facteur direct, et
G4 à G6 portent sur l'exactitude et la conservativité de $F$. C'est un état
antérieur, où la décomposition en composantes connexes est encore un axiome
(Gal 2) plutôt qu'une conséquence. La page 209 du même carton donnera un
troisième état, sous le nom de G0 à G5, et pour une classe d'objets plus large.}

Le remplacement de « indécomposable » par « connexe » n'est pas une correction
de style : c'est le moment où le vocabulaire de la géométrie passe dans une
axiomatique qui n'en a pas besoin, et c'est le mot que les pages 190 à 216
emploieront sans plus jamais l'expliquer.

\subsection*{La seconde version : des points au lieu d'un foncteur}

La page 158 recommence, avec un dispositif différent : une catégorie
$\mathcal{C}$, une sous-catégorie \emph{pleine} $\mathcal{P}$ dont les objets
s'appellent les \emph{points}, et une classe de morphismes dits \emph{non
ramifiés}. Les neuf conditions, numérotées $1^{\mathrm{o}}$ à
$9^{\mathrm{o}}$ dans un ordre qui n'est pas celui de leur écriture, disent :
que les sommes finies et les produits fibrés existent ; que les morphismes non
ramifiés sont stables par composition et changement de base ; qu'il y a des
composantes connexes ; que les points sont connexes, de sorte que
$\operatorname{Hom}(\xi, X \sqcup Y) = \operatorname{Hom}(\xi, X) \sqcup
\operatorname{Hom}(\xi, Y)$ ; qu'une section d'un morphisme non ramifié sur un
objet connexe est un isomorphisme sur une composante connexe ; et qu'un
morphisme non ramifié dont les fibres $\operatorname{Hom}(\xi, X) \to
\operatorname{Hom}(\xi, Y)$ sont réduites à un élément est un isomorphisme.
Enfin, la neuvième donne l'existence de $X/H$ comme en Gal 6.

Sous ces conditions on définit $\pi_{1}(X, x)$ pour $x \in
\operatorname{Hom}(\xi, X)$, et --- c'est la dernière ligne de la page ---
« c'est un foncteur \emph{covariant} en $(X, x)$ » :

\begin{tikzcd}
\xi \arrow[r, "x"] \arrow[d] & X \arrow[d] \\
\xi' \arrow[r, "x'"'] & X'
\end{tikzcd}

Le foncteur fibre a été remplacé par la famille des points ; c'est le passage
du groupe fondamental au \emph{groupoïde} fondamental, et c'est la structure
que les catégories multigaloisiennes de la page 208 axiomatiseront.

\pagerange{159}{160}
\section*{« Groupe fondamental dans les catégories » (pages 159 et 160)}

Deux dernières pages poussent l'abstraction d'un cran, et c'est le cadre le
plus large du carton.

Soit $\mathcal{C}$ une catégorie munie d'une famille $(D)$ de morphismes
satisfaisant :

\begin{enumerate}
  \item[(i)] $(D)$ contient les identités et est stable par composition ;
  \item[(ii)] $(D)$ est stable par changement de base ;
  \item[(iii)] les morphismes de $(D)$ sont des morphismes de descente
    \emph{universels stricts} ;
  \item[(iv)] les morphismes qui « descendent » par un $g \in (D)$ le font pour
    un $f \in (D)$ --- mais \emph{on ne suppose pas} qu'on puisse faire
    descendre tout morphisme par un $f \in (D)$ ;
  \item[(v)] la projection $X \times I \to X$ est dans $(D)$.
\end{enumerate}

\emph{Exemple.} $\mathcal{C}$ la catégorie des préschémas, $(D)$ les morphismes
\emph{plats et propres} --- ou plats et quasi-propres.

On développe alors la théorie des fibrés principaux dans ce cadre : les groupes
structuraux sont les $X \times G$ pour $G$ un groupe abstrait fini, et les
fibrés principaux sur $X$ de groupe $G$ sont ceux dont la projection
$P \to X$ appartient à $(D)$.

Il faut voir ce que la clause (iv) refuse. Elle dit qu'on ne demande pas à
$(D)$ d'être une classe de recouvrements au sens d'une topologie : tout objet
n'est pas descendable. C'est la même retenue que la page 190 exprimera en
posant l'axiome \textbf{R3} --- une famille de sous-objets dont les
réalisations recouvrent est \emph{de descente effective} --- au lieu de
supposer que tout est recouvert. Le cadre le plus abstrait du carton est aussi
celui qui promet le moins, et c'est là son intérêt.

Les trois dernières lignes de la page 160 amorcent les notions d'espaces à
opérateurs admissibles et d'homomorphismes admissibles de groupes structuraux,
et s'arrêtent.

\pagerange{161}{162}
\section*{La normalité n'est pas une commodité (pages 161 et 162)}

\subsection*{L'énoncé}

Soit $f : X \to Y$ un morphisme propre dominant, $Y$ \emph{normal} connexe et
localement noethérien. Alors l'image de
\[
  \pi_{1}(X) \longrightarrow \pi_{1}(Y)
\]
est d'indice fini.\footnote{La page écrit « a un conoyau fini ». Pour un
morphisme de groupes profinis l'image n'est pas nécessairement distinguée et il
n'y a pas de conoyau ; ce que l'énoncé veut dire, et ce que la démonstration
donne, est que l'image est d'\emph{indice fini}. C'est la formulation retenue
ici.} La démonstration se ramène au cas $Y = \operatorname{Spec}(K)$, en
remplaçant $K$ par une extension séparable finie convenable et $X$ par une
composante connexe de $X'$.

\subsection*{Le contre-exemple}

C'est la moitié lisible de la page, et c'en est le propos :

\begin{quote}
Si $Y$ n'est pas \emph{normal}, le théorème est faux.
\end{quote}

Prenons pour $Y$ une courbe rationnelle sur $k$ algébriquement clos ayant un
point double ordinaire --- une cubique nodale, par exemple --- et pour $X$ sa
normalisée. Alors $X \simeq \mathbb{P}^{1}_{k}$, donc
\[
  \pi_{1}(X) \;=\; 0 ,
\]
tandis que
\[
  \pi_{1}(Y) \;\simeq\; \hat{\mathbb{Z}} .
\]
L'image est triviale dans un groupe infini : elle n'est pas d'indice fini.

Le groupe $\pi_{1}(Y)$ se voit à l'œil. Prenons $n$ copies
$C_{1}, \dots, C_{n}$ de $\mathbb{P}^{1}$ et recollons $\infty_{i}$ à
$0_{i+1}$ cycliquement ; le collier ainsi obtenu s'envoie sur $Y$, chaque
$C_{i}$ par la normalisation, et le morphisme est fini étale de degré $n$ ---
au point double, les deux branches s'envoient isomorphiquement sur les deux
branches. C'est un revêtement galoisien de groupe $\mathbb{Z}/n$, et il existe
\emph{pour tout $n$}. Le groupe fondamental est donc $\hat{\mathbb{Z}}$ tout
entier.\footnote{La page porte « $\pi_{1}(Y) = \hat{\mathbb{Z}}$
privé de sa $p$-composante », ces derniers mots étant une lecture que la
transcription signale comme douteuse. Elle n'est en tout cas pas nécessaire :
les revêtements en collier existent en degré $p$ comme en tout autre degré,
parce qu'ils viennent du graphe dual de la courbe et non de sa géométrie ---
c'est le $H^{1}$ d'un cercle, qui ne connaît pas la caractéristique. La
restriction au premier à $p$ est celle qui s'impose pour les revêtements
\emph{d'une courbe lisse}, où le $p$-groupe fondamental est sauvage ; elle ne
s'impose pas ici. Le contre-exemple est d'ailleurs plus fort sans elle.}
Enfin, l'hypothèse $X \simeq \mathbb{P}^{1}$ --- c'est-à-dire $Y$
\emph{rationnelle} --- est ce qui donne $\pi_{1}(X) = 0$ ; la page ne l'écrit
pas et elle est nécessaire.

Grothendieck ajoute une observation qui explique où le raisonnement casserait :
dans ce cas, on ne peut pas former $\underline{\mathrm{Hom}}_{Y}(X, Y \sqcup
Y)$. C'est le schéma des morphismes du lemme 1 de la page 47, et son absence
est ce qui interdit de descendre l'isomorphisme générique. La normalité de $Y$
était ce qui le faisait exister.

La page 162 ne porte qu'un carré encadré, $f_{*}(X_{1})$ au-dessus de
$S \leftarrow T$, et rien d'autre.

\pagerange{164}{169}
\section*{Le progroupe fondamental pour les revêtements infinis (pages 164 à 169)}

Le feuillet qui ouvre la suite porte de sa main, avec une rature :
« \emph{Le Groupe} \; Progroupe fondamental pour les revêtements infinis
principaux ». La rature est où le nom a changé : ce n'est pas un groupe, c'est
un pro-groupe.

\subsection*{Le foncteur}

Soit $S$ un préschéma connexe et $\xi$ un point géométrique. Pour tout groupe
$\pi$ --- \emph{pas nécessairement fini} --- on note $\pi^{1}(S, \xi\,;\, \pi)$ l'ensemble des classes d'isomorphie de
fibrés principaux homogènes sur $S$, de groupe structural $\pi$, pointés
au-dessus de $\xi$.

C'est le même foncteur que le $Z(S,a;G)$ de la lettre à Serre et que le
$\mathfrak{Z}(S,a;G)$ des feuilles anciennes, avec la restriction de finitude
levée ; et c'est cette levée qui fait toute la difficulté de ces pages.

\subsection*{$\alpha$) Le foncteur est exact à gauche}

\textbf{a) Il commute aux produits finis.} Trivial.

\textbf{b) Il transforme une injection en une injection.} Soit
$\pi \hookrightarrow \pi'$. La raison est le lemme d'unicité déjà rencontré :
entre deux $\pi'$-fibrés pointés il existe \emph{au plus un} isomorphisme
respectant les points marqués. Si $P'_{1}$ et $P'_{2}$ proviennent de
$\pi$-fibrés $P_{i} \subset P'_{i}$ avec $\xi_{i} = \xi'_{i}$, l'unique
isomorphisme $u' : P'_{1} \to P'_{2}$ se restreint donc en
$u : P_{1} \to P_{2}$.

Et la page 165 dit exactement d'où vient la réduction du groupe structural :
\emph{$P_{i}$ se déduit de $P'_{i}$ par une section de $P'_{i}/\pi$ pointée par
$\xi'_{i}$ ; et une telle section, si elle existe, est unique}. Réduire le
groupe structural, c'est trancher une section ; le point marqué la rend unique.

\textbf{c) Il préserve l'exactitude d'un couple.} Si
$\pi \to \pi' \rightrightarrows \pi''$ est exacte, il en va de même des images.
On considère $Q$, le sous-schéma des coïncidences des deux morphismes
$u_{1}, u_{2} : P' \to P''$ ; il est ouvert et fermé, contient $\xi'$, est
stable par $\pi$, et lui correspond donc un sous-préschéma ouvert et fermé
$S'$ de $P'/\pi$, fidèlement plat et quasi-compact sur
$S$.\footnote{Les pages 164 et 165 sont d'une écriture très rapide et la
transcription y laisse la plupart des phrases de liaison illisibles. Les trois
énoncés a), b), c) et les deux raisons portantes --- l'unicité de l'isomorphisme
pointé, l'unicité de la section --- sont ce que la page donne ; le détail des
vérifications ne l'est pas et n'est pas reconstruit ici.}

\subsection*{d) La filtration décroissante}

C'est le point nouveau, et il demande une hypothèse : \emph{$S$ localement
noethérien, ou bien $\pi$ artinien} (condition de chaîne descendante).

Soit $c \in \pi^{1}(S,\xi;\pi)$ correspondant à $(P, \xi)$. Considérons deux
sous-groupes $\pi'$, $\pi''$ auxquels on peut restreindre le groupe structural.
Alors
\[
  P/(\pi' \cap \pi'') \longrightarrow P/\pi' \times P/\pi''
\]
est une \emph{immersion ouverte et fermée}, et le point $(y', y'')$ donné par
les deux réductions est dans son image. On peut donc restreindre le groupe
structural à $\pi' \cap \pi''$.

Il en résulte que \emph{les groupes auxquels on peut restreindre forment un
système filtrant décroissant} $(\pi_{i})$, avec des $P_{i} \subset P$
correspondants. C'est l'analogue, pour des groupes infinis, de la condition de
minimalité de la lettre de 1959 : là, la chaîne descendante était stationnaire
par hypothèse ; ici elle ne l'est pas, et il faut montrer que l'intersection
porte encore un fibré.

\subsection*{L'intersection porte un fibré}

C'est ce que démontrent les pages 167 et 169.

La question est locale sur $S$, et par étalement sur un $S^{(0)}$ de type fini
sur $\mathbb{Z}$ on peut supposer $S$ \emph{noethérien}. Soit alors $P^{0}$ la
composante connexe de $\xi$ dans $P$ : elle est ouverte et fermée, $P$ étant
localement noethérien donc localement connexe.

\textbf{Assertion.} $P^{0} \to S$ est \emph{surjectif}.

L'image est constructible et stable par générisation ; on se ramène au cas d'un
anneau de valuation complet, où l'on conclut.

Alors, puisque $P^{0} \subset \bigcap_{i} P^{(i)}$ et que chacun des $P^{(i)}$
est une réunion de composantes connexes de $P$, il en va de même de leur
intersection $P'$. Celle-ci est stable sous $\pi' = \bigcap_{i} \pi_{i}$, et
ses fibres géométriques sont non vides puisqu'elles contiennent celle de
$P^{0}$. Donc \emph{$P'$ est un fibré principal homogène sous $\pi'$}.

\subsection*{La conséquence}

\textbf{Conséquence.} Si $\pi$ varie dans une catégorie de groupes artiniens
stable par produits finis et par sous-groupes, le foncteur
\[
  \pi \longmapsto \pi^{1}(S, \xi\,;\, \pi)
\]
est \emph{strictement pro-représentable}.

\subsection*{L'application : ce que le groupe élargi classe}

La page 168 dit à quoi cela sert, et c'est l'énoncé qui donne son sens à ces pages.

Une relation étale est dite \emph{à groupe structural $\pi$} si le schéma des
isomorphismes $\underline{\mathrm{Isom}}_{S}(E_{S}, T)$ est un revêtement
principal de groupe $\pi_{S}$. Si $S$ est connexe muni d'un point géométrique
$\xi$, alors \emph{se donner un $T$ muni d'une telle structure revient à se
donner des opérations continues de $\tilde{\pi}_{1}(S, \xi)$ sur $F$},
c'est-à-dire un homomorphisme continu $\tilde{\pi}_{1}(S,\xi) \to \pi$.

\textbf{Exemple.} Les schémas en groupes commutatifs sans torsion sur $S$, de
type fini, correspondent aux $\tilde{\pi}_{1}(S,\xi)$-modules qui sont des
$\mathbb{Z}$-modules de type fini.

C'est exactement ce que le groupe fondamental ordinaire ne peut pas faire : un
$\mathbb{Z}$-module de type fini n'est pas un ensemble fini, et le groupe
profini ne classe que des ensembles finis. Le groupe élargi est la réponse à ce
manque.\footnote{Cet objet --- le groupe fondamental classifiant les objets
localement constants à fibres non nécessairement finies --- porte aujourd'hui
plusieurs noms selon la classe de coefficients qu'on autorise : \emph{groupe
fondamental élargi} dans SGA 3, et \emph{groupe fondamental pro-étale} dans le
travail de Bhatt et Scholze (2015), qui en donne la construction topologique
correcte pour un schéma quelconque. Le carton en donne ici la construction par
pro-représentabilité, et ne discute pas la topologie du pro-groupe obtenu.}

\pagerange{170}{171}
\section*{Les cinq classes de recouvrements (pages 170 et 171)}

Ces deux pages sont écrites au net et se lisent presque entièrement --- ce qui,
dans ce carton, mérite d'être signalé.

Soit $S$ un préschéma. Pour tout groupe \emph{discret} $G$ on pose
\[
  H^{1}_{\mathcal{C}}(S, G) \;=\; \varinjlim_{T/S} H^{1}(T/S, G) ,
  \qquad
  H^{1}(T/S, G) \;=\; H^{1}\bigl(\pi_{0}(K_{T/S}),\, G\bigr) ,
\]
la limite étant prise sur une catégorie $\mathcal{C}$ de $S$-préschémas. Les
cinq classes qui importent sont :

\begin{itemize}
  \item[$C_{1}$] les $S$-préschémas fidèlement plats quasi-compacts ;
  \item[$C_{2}$] les mêmes, de type fini ;
  \item[$C_{3}$] les mêmes, quasi-finis ;
  \item[$C_{4}$] les mêmes, finis principaux ;
  \item[$C_{5}$] les mêmes, finis étales principaux.
\end{itemize}

On a des applications injectives canoniques $H^{1}_{C_{i}}(S,G)
\hookrightarrow H^{1}_{C_{j}}(S,G)$ pour $C_{i} \subset C_{j}$.

\subsection*{Pourquoi cela n'a d'intérêt que pour $G$ infini}

Si $G$ est \emph{fini}, tous ces homomorphismes sont bijectifs, et
$H^{1}(S,G)$ est simplement l'ensemble des classes de revêtements étales
principaux de groupe $G$, soit
\[
  \operatorname{Hom}\bigl(\pi_{1}(S,a),\, G\bigr) \big/
  \operatorname{int}(G) .
\]
La raison est le théorème de descente : $G_{T}$ est alors fini, donc affine sur
$T$, et la descente fidèlement plate quasi-compacte des schémas affines est
effective. Toutes les classes donnent la même réponse, et l'appareil est vide.

D'où la phrase qui justifie ces pages : \emph{« Notre intérêt ici est le cas $G$
infini »}.

\subsection*{Les comparaisons}

\begin{itemize}
  \item En général, $H^{1}_{C_{3}} \simeq H^{1}_{C_{2}}$ : quasi-fini ou de
    type fini, cela ne change rien.
  \item Si $S$ est le spectre d'un anneau de Dedekind,
    $H^{1}_{C_{2}} \xrightarrow{\ \sim\ } H^{1}_{C_{1}}$, d'où
    $H^{1}_{C_{3}} \simeq H^{1}_{C_{2}} \simeq H^{1}_{C_{1}}$.
  \item Si $S$ est le spectre d'un anneau de valuation discrète complet à corps
    résiduel algébriquement clos, $H^{1}_{C_{i}} \simeq H^{1}(\hat{V}/V, -)$
    pour $i = 1, 2, 3$ : une seule extension suffit.
  \item Si $S$ est le spectre d'un anneau local \emph{complet} $V$ de corps
    résiduel $k$, alors $H^{1}_{C_{5}} \simeq H^{1}_{C_{4}} \simeq
    H^{1}_{C_{3}} \simeq H^{1}_{C_{2}}$, et de plus
    $H^{1}_{C_{i}}(V, -) \simeq H^{1}_{C_{i}}(k, -)$ pour $i = 2,3,4,5$.
\end{itemize}

La dernière ligne est la version « groupes infinis » du fait que le groupe
fondamental d'un anneau local complet est celui de son corps résiduel.

\pagerange{173}{173}
\section*{Un plan en sept sections (page 173)}

Le feuillet qui précède porte, de sa main, le seul mot « Plan ». Voici ce qu'il
annonce, pour deux théorèmes de changement de base :

\begin{itemize}
  \item[I] Premier théorème de changement de base pour $\pi_{1}$, ou théorème
    de comparaison.
  \item[II] Deuxième théorème de changement de base pour $\pi_{1}$.
\end{itemize}

\begin{enumerate}
  \item Morphismes $(-1)$-connexes et $0$-connexes.
  \item Morphismes \emph{locaux} $(-1)$-connexes et $0$-connexes.
  \item Caractérisation des morphismes locaux ou globaux $i$-connexes
    ($i = 0, 1$).
  \item Deuxième théorème du changement de base pour $\pi_{1}$ : énoncés et
    corollaires.
  \item Démonstration des théorèmes.
  \item Généralisation à l'aide de la résolution des singularités.
  \item Application : la dimension cohomologique de certains corps.
\end{enumerate}

Le vocabulaire des morphismes $i$-connexes --- un morphisme est $(-1)$-connexe
s'il est surjectif au sens des faisceaux, $0$-connexe s'il induit de plus une
bijection sur les $\pi_{0}$, $1$-connexe s'il induit un isomorphisme sur les
$\pi_{1}$ --- est celui de l'homotopie, transposé. Le plan n'est pas exécuté
dans ce carton.

\pagerange{175}{176}
\section*{$\pi_{2}$ contre la tour des revêtements (pages 175 et 176)}

Le feuillet qui ouvre la suite porte, de sa main, « $\Pi_{2}$ ». Ce sont deux
pages de calcul, sans phrases.

\subsection*{Le calcul}

On dispose de la tour d'un revêtement $X^{i} \to X$, de son revêtement
universel $\tilde{X}$, et des revêtements universels itérés
$\tilde{\tilde{X}}^{i}$ :

\begin{tikzcd}[column sep=small, row sep=small, nodes={font=\scriptsize}]
\widetilde{\widetilde{X}}{}^i \arrow[d] & \\
\widetilde{X}^i \arrow[d] & \widetilde{X} \arrow[d, no head] \\
X^i \arrow[r] & X
\end{tikzcd}

et l'on pose $\tilde{\pi}/H_{i} = \pi_{i}$. Le calcul est
\[
  \pi^{2}(X, G) \;\simeq\; H^{2}(\check{X}, G) \;\simeq\;
  \varinjlim_{i} H^{2}(X_{i}, G) \;\simeq\;
  \varinjlim_{i} H^{2}(\tilde{X}_{i}, G) ,
\]
la suite spectrale de Hochschild--Serre donnant
\[
  E_{2}^{p,q} \;=\; H^{p}\bigl(H_{i},\, H^{q}(\tilde{X}, G)\bigr)
  \;\Longrightarrow\; H^{*}(\tilde{X}_{i}, G) ,
\]
puis, en passant à la limite sur $i$, la suite exacte
\[
  \cdots \to \varinjlim_{i} H^{2}(H_{i}, G) \to \pi^{2}(X, G) \to
  \varinjlim_{i} \pi^{2}(\tilde{X}_{i}, G)^{H_{i}} \to
  \varinjlim_{i} H^{3}(H_{i}, G) .
\]
Grothendieck entoure les deux extrémités, qui sont ce qu'il voulait :
\[
  \operatorname{Ker}\bigl(\pi^{2}(X,G) \to \pi^{2}(X^{k},G)\bigr) \;\simeq\;
  \varinjlim_{i} H^{2}(H_{i}, G) ,
\]
\[
  \operatorname{Coker} \;\simeq\;
  \operatorname{Ker}\Bigl(\varinjlim_{i} H^{3}(H_{i}, G) \longrightarrow
  H^{3}(\tilde{X}, G)\Bigr) .
\]
Autrement dit : \emph{ce que le passage aux revêtements fait perdre et gagne en
$\pi_{2}$ est de la cohomologie des groupes de la tour}.\footnote{L'exposant du
premier terme est tracé comme un « I » romain sur la page, là où quatre lignes
plus bas le même objet est noté $\pi^{2}$ ; la transcription le signale et le
conserve. C'est $\pi^{2}$ qui est écrit ici.}

\subsection*{Le cas d'une fibration}

La page 176 fait la même chose pour une fibration $\bar{F} \to X \to Y$, avec
la suite exacte d'homotopie $\pi_{1}(\bar{F}) \to \pi_{1}(X) \to \pi_{1}(Y)
\to e$, et aboutit à l'annulation
\[
  H^{1}\bigl(Y'/Y,\ \operatorname{Hom}(\pi_{1}(\bar{F}),\, \mathfrak{z})\bigr)
  \;=\; 0 ,
\]
qui est la condition sous laquelle le passage de la fibre à l'espace total est
sans obstruction.\footnote{Le groupe noté $\mathfrak{z}$ ne porte pas de nom
sur la page ; c'est le centralisateur $G^{\pi_{1}(\bar{F})}$ qui apparaît deux
lignes plus haut. La page est faite de formules sans phrases liantes, et cette
lecture n'invente pas les phrases : elle donne le calcul et son point
d'arrivée.}

\pagerange{178}{180}
\section*{Descente : pro-représentabilité, et une question ouverte (pages 178 à 180)}

Le feuillet qui ouvre la suite porte, de sa main : « $\pi_{0}$, $\pi_{1}$ et
formalisme simplicial (Revêtements infinis) ».

\subsection*{La proposition}

Soit $S$ un préschéma \emph{connexe}, $S' \to S$ un morphisme
\emph{submersif} --- par exemple fidèlement plat quasi-compact, ou propre
surjectif --- et $S''$, $S'''$ le carré et le cube fibrés. On suppose que $S'$
est \emph{somme} de préschémas connexes, c'est-à-dire que $\pi_{0}(S')$ est
discret. Soit $\xi$ un point géométrique de $S$.

Alors, pour un groupe $G$ variable, le foncteur \emph{covariant}
\[
  G \longmapsto \pi^{1}(S'/S, \xi\,;\, G)
\]
--- l'ensemble des classes, à isomorphisme près, de \emph{données de descente}
sur $G_{S'}$ relatives à $S' \to S$, pointées au-dessus de $\xi$ --- est
strictement pro-représentable :
\[
  \pi^{1}(S'/S, \xi\,;\, G) \;=\; \varinjlim_{i}
  \operatorname{Hom}(\pi_{i},\, G) .
\]

Deux raffinements. Si $\pi_{0}(S'')$ est quasi-compact, les $\pi_{i}$ sont de
type fini. Et si $S''$ n'a qu'un nombre fini $n''$ de composantes connexes, le
foncteur est \emph{représentable} :
\[
  \pi^{1}(S'/S, \xi\,;\, G) \;\simeq\;
  \operatorname{Hom}\bigl(\pi_{1}(S'/S, \xi),\, G\bigr) ,
\]
avec $\pi_{1}(S'/S,\xi)$ un groupe \emph{discret} de type fini, engendré par
$n''$ générateurs.\footnote{La page écrit d'abord « contravariant », puis le
biffe. C'est bien un foncteur covariant : une donnée de descente sur $G_{S'}$
se pousse le long d'un homomorphisme $G \to G'$. La correction est de sa main
et elle est suivie ici.}

C'est le \emph{groupe fondamental de la descente} : un groupe discret de type
fini, engendré par autant d'éléments que $S''$ a de composantes connexes. La
finitude du nombre de générateurs vient donc de la combinatoire du
recouvrement, non de la géométrie --- exactement comme au contre-exemple de la
page 161.

\subsection*{La question ouverte}

Considérons de même le foncteur
$G \mapsto \pi^{1}_{\mathrm{eff}}(S'/S, \xi\,;\, G)$ des données de descente
\emph{effectives} pointées. Il est également pro-représentable, par un système
projectif de \emph{quotients} des $\pi_{i}$.

Et vient l'aveu :

\begin{quote}
J'ignore, si $S$, $S'$, $S''$ sont noethériens, si ce dernier foncteur est
représentable ; \emph{a fortiori} j'ignore si dans ce cas
$\pi^{1} = \pi^{1}_{\mathrm{eff}}$, c'est-à-dire si \emph{toute} donnée de
descente sur $G_{S'}$ est \emph{effective}.
\end{quote}

Il ajoute que c'est vrai sous une hypothèse supplémentaire, et note en marge :
« Pb intér.\ pour l'affirmative [\dots\ et $S' \to S$ fidèlement plat
quasi-compact] ».

Il vaut la peine de dire pourquoi la question est ouverte, parce que la raison
se lit en rapprochant cette page de la page 170.

Lorsque $G$ est \emph{fini}, le schéma en groupes constant $G_{S'}$ est fini
sur $S'$, donc \emph{affine} sur $S'$ ; et le théorème de descente fidèlement
plate de Grothendieck dit précisément que la descente est effective pour les
morphismes affines. La question est alors réglée, et c'est ce que la page 170
constatait en disant que les cinq classes donnent la même réponse pour $G$
fini.

Lorsque $G$ est \emph{infini}, $G_{S'}$ est la somme disjointe d'une famille
infinie de copies de $S'$ : elle n'est pas quasi-compacte sur $S'$, donc pas
affine, et le théorème de descente ne l'atteint plus. \emph{La question est
donc ouverte exactement dans le cas que la page 170 déclarait être le seul
intéressant.}

\subsection*{La traduction simpliciale}

La page 180 réduit tout cela à de la combinatoire. On dispose du diagramme
semi-simplicial des $X_{i}$, de leurs fibres $X_{i}^{\xi}$, et des ensembles
\[
  K_{i}^{\xi} \;=\; \pi_{0}(X_{i}^{\xi}) ,
\]
avec, en dessous, la version pointée $k_{i}$ et les inclusions $i_{j}$. Alors :

\begin{quote}
Les données de descente sur $E_{X_{0}}$ correspondent aux données de descente
sur $E_{K_{0}}$, c'est-à-dire aux applications ensemblistes
$f : K_{1} \to G$ vérifiant
\[
  f \circ p_{1} \;=\; (f \circ p_{0})\,(f \circ p_{2}) ,
\]
deux telles applications étant équivalentes si
$f'(q) = h(p_{1}q)\, f(q)\, h(p_{0}q)^{-1}$ ; et les données \emph{pointées}
sont celles qui vérifient de plus $f(i_{1}(q)) = e_{G}$ pour tout
$q \in k_{1}$, l'équivalence étant restreinte aux $h$ tels que
$h(i_{0}(q)) = 1$.
\end{quote}

C'est la relation de cocycle non abélienne en degré 1, et c'est la même que
celle que la page 144 encadrait : $A_{\alpha''} a''_{\beta} A_{\alpha'} =
a'''_{\beta} A_{\alpha'''} a'_{\beta}$. Elle est ici écrite proprement, sur
un ensemble semi-simplicial, ce qui rend visible sa nature : \emph{un
$1$-cocycle non abélien sur le nerf du recouvrement}. La
pro-représentabilité de la proposition ci-dessus n'est alors rien d'autre que
la présentation par générateurs de ce nerf --- d'où les $n''$ générateurs. Les
pages 181 à 185 poursuivent ce calcul, sur des foncteurs semi-cosimpliciaux
abstraits.

\pagerange{181}{185}
\section*{Foncteurs semi-cosimpliciaux, et le van Kampen (pages 181, 182, 184 et 185)}

Ces quatre pages continuent le calcul de la page 180 en le rendant abstrait :
au lieu du diagramme simplicial d'un recouvrement, un diagramme de foncteurs.
Elles sont d'un crayon très pâle sur un feuillet dont le verso transparaît, et
la transcription y laisse beaucoup d'illisible.

\subsection*{Quand un foncteur est pro-représentable}

Soit $F$ un endofoncteur de la catégorie des ensembles. On cherche quand
\[
  F(I) \;\simeq\; C(K, I) ,
\]
$C(K,I)$ désignant les applications continues d'un espace $K$ totalement
discontinu dans l'ensemble discret $I$ --- c'est-à-dire les fonctions
localement constantes à valeurs dans $I$. Trois conditions nécessaires :

\begin{enumerate}
  \item[a)] $F$ commute aux limites projectives finies, c'est-à-dire est exact
    à gauche ;
  \item[b)] $F\bigl(\bigcap_{\alpha} I_{\alpha}\bigr) =
    \bigcap_{\alpha} F(I_{\alpha})$ pour toute famille filtrante de
    sous-ensembles de $I$ ;
  \item[c)] $\varprojlim_{j} \pi_{0j} \to \pi_{0i}$ est \emph{surjectif} pour
    tout $i$.
\end{enumerate}

Et le point : \emph{ces conditions sont nécessaires et suffisantes}, et
équivalent à dire que $F$ est \emph{strictement pro-représentable} par un
système projectif strict d'ensembles $(\pi_{0i})$ ; on prend alors
$K = \varprojlim \pi_{0i}$ muni de la topologie limite des topologies discrètes.

La condition c) est celle qui distingue un système projectif \emph{strict} d'un
système quelconque : elle dit qu'aucune information n'est perdue en passant à
la limite. Sans elle, $K$ pourrait être vide sans qu'aucun $\pi_{0i}$ le soit.

Une remarque en note, avec sa référence : si $F$ est représenté de deux façons
par des $C(L, I)$ avec $L$ un groupe topologique, on a bien
$\pi_{0}(L) \to K$, mais si $L$ n'est \emph{pas compact} ce n'est ni un
isomorphisme, ni surjectif, ni injectif.\footnote{La référence portée en marge
est « Bourb.\ Topologie II, p.\ 232, prop.\ 6 ». C'est la mise en garde qui
justifie tout le soin apporté à la stricte pro-représentabilité : le groupe
topologique qui « représente » n'est déterminé que sous une hypothèse de
compacité, et c'est pourquoi le carton parle de \emph{pro}-objets et non de
groupes.}

\subsection*{Le noyau, et le foncteur semi-cosimplicial}

Soit $F^{*} : F^{0} \rightrightarrows F^{1}$ un couple de foncteurs, et
\[
  \pi^{0}(F^{*}) \;=\; \operatorname{Ker}\bigl(F^{0} \rightrightarrows
  F^{1}\bigr) .
\]
Si $F^{0}$ et $F^{1}$ sont pro-représentables par $K_{0}$ et $K_{1}$, les deux
flèches proviennent de $K_{0} \leftleftarrows K_{1}$, et $\pi^{0}(F^{*})$ l'est
aussi.

On considère alors un \emph{foncteur semi-cosimplicial}

\begin{tikzcd}[column sep=small, row sep=small]
F^0 \arrow[r, bend left=12] \arrow[r, bend right=12]
  & F^1 \arrow[r, bend left=20] \arrow[r] \arrow[r, bend right=20] & F^2
\end{tikzcd}

auquel on associe deux foncteurs $\pi^{0}(F^{*}; E)$ et $\pi^{1}(F^{*}; G)$,
pour $E$ un ensemble et $G$ un groupe variables ; et, si l'on se donne de plus
un \emph{pointage} $\xi : F^{*} \to f^{*}$,

\begin{tikzcd}[column sep=small, row sep=small, nodes={font=\scriptsize}]
F^0 \arrow[r, bend left=12] \arrow[r, bend right=12] \arrow[d, "\xi^0"]
  & F^1 \arrow[r, bend left=20] \arrow[r] \arrow[r, bend right=20]
    \arrow[d, "\xi^1"] & F^2 \arrow[d, "\xi^2"] \\
f^0 \arrow[r, bend left=12] \arrow[r, bend right=12]
  & f^1 \arrow[r, bend left=20] \arrow[r] \arrow[r, bend right=20] & f^2
\end{tikzcd}

les versions pointées $\pi^{0}(F^{*}, \xi; E)$ et $\pi^{1}(F^{*}, \xi; G)$.
Lorsque tout est pro-représentable, on se ramène par un procédé standard à un
\emph{ensemble semi-simplicial pointé}
\[
  K_{0} \leftleftarrows K_{1} \Lleftarrow K_{2} ,
  \qquad
  k_{0} \leftleftarrows k_{1} \Lleftarrow k_{2} ,
\]
et l'on retrouve exactement les cocycles de la page 180. C'est le $\pi^{0}$ et
le $\pi^{1}$ non abéliens d'un objet semi-cosimplicial.

\subsection*{Le van Kampen}

La page 185 se termine sur le programme, et il faut le lire :

\begin{quote}
Il y a lieu de généraliser ces constructions --- on a vu que des $K_{i}$, des
$\overline{K}_{i}$, on a des catégories galoisiennes à chaque point, pour
décider de descente : le \emph{Van Kampen} \dots
\end{quote}

C'est-à-dire : la descente d'un revêtement le long d'un recouvrement est un
calcul de $\pi^{0}$ et de $\pi^{1}$ sur le nerf de ce recouvrement, et le
théorème de van Kampen --- qui reconstruit le groupe fondamental d'un espace à
partir de ceux d'un recouvrement --- en est la forme classique. Le carton s'en
servait déjà librement à la page 60 ; il en donne ici la raison.\footnote{Le
cadre général qu'appellent ces pages --- pro-objets, ensembles semi-simpliciaux,
$\pi^{i}$ non abéliens d'un objet cosimplicial --- est celui de l'homotopie
étale d'Artin et Mazur, dont le mémoire paraît en 1969. Le carton n'est pas
daté sur ces pages et cette lecture n'en tire pas de chronologie.}

\pagerange{183}{187}
\section*{Trois feuillets d'un autre ordre (pages 183, 186 et 187)}

\subsection*{Un feuillet d'esquisses}

La page 183 est un feuillet en travers, couvert de tracés sans ordre de
lecture : le diagramme $\overline{K}_{0} \leftleftarrows \overline{K}_{1}
\Lleftarrow \overline{K}_{2}$ avec les mentions « fini » et « groupe
topologique discret » ; la formule d'équivalence des données de descente de la
page 180 ; l'anneau $k[[s,t]]$ ; un chemin d'arêtes de $x$ à $a$ dont les
flèches ne vont pas toutes dans le même sens ; et, encadré, un symbole
$\pi_{1}^{\#}(S,a)$ surmonté de « $S$ connexe ». Le croisillon est tracé deux
fois de la même façon, donc c'est une notation ; rien sur la feuille n'en fixe
le sens, et cette lecture ne le conjecture pas.

\subsection*{Une proposition sur les morphismes de type fini}

La page 186 est écrite au net, à l'encre bleue, et porte une proposition sans
rapport avec ce qui l'entoure.

\textbf{Proposition.} Soit $f : X \to Y$ un morphisme avec 1) $Y$ localement
noethérien ; 2) $X$ n'ayant qu'un nombre fini de composantes irréductibles ;
3) $f$ localement de type fini et \emph{séparé}. Alors $f$ est \emph{de type
fini}.

La démonstration se ramène à $Y$ affine noethérien et $X$ intègre, puis
construit un $X''$ fini au-dessus de $X$ et birationnel au-dessus d'un $X'$
intègre, de sorte que $X'' \to X$ soit fini et surjectif ; il suffit alors de
voir $X''$ de type fini sur $Y$, et $X'' \to X'$ étant birationnel entre
schémas intègres, on conclut par la normalisation.\footnote{L'hypothèse 1)
porte, entre « $Y$ loc.\ noeth. » et « de type fini », un mot que la
transcription ne lit pas, et les hypothèses 1) et 3) sont l'une et l'autre
biffées et réécrites sur la page. Telle qu'elle se laisse lire, la proposition
n'assure pas visiblement la quasi-compacité de $X$, qui est ce qu'il faut
démontrer ; c'est sans doute là que le mot manquant se trouve. La proposition
est donc enregistrée ici comme la page la porte, et non réparée : le seul
endroit d'où la quasi-compacité pourrait venir est l'hypothèse illisible.}

\subsection*{Un opérateur d'homotopie pour $d''$}

La page 187 est barrée de deux longues diagonales et se lit néanmoins d'un bout
à l'autre. Elle est d'analyse complexe --- le seul feuillet du carton qui le
soit --- et donne un opérateur d'homotopie explicite pour l'opérateur de
Cauchy--Riemann $d''$ en $n$ variables.

Sur un polydisque, on dispose pour chaque variable de deux opérateurs
intégraux : le noyau de Cauchy, qui projette sur les fonctions holomorphes en
$z_{j}$,
\[
  U_{j} f \;=\; \frac{1}{2\pi i} \int_{B_{j}}
    f(z_{1}, \ldots, t, \ldots, z_{n})\, \frac{dt}{t - z_{j}} ,
\]
et l'opérateur de Cauchy--Pompeiu, qui résout $\partial/\partial\bar{z}_{j}$,
\[
  T_{j} f \;=\; \frac{1}{2\pi i} \iint_{K_{j}}
    f(z_{1}, \ldots, t, \ldots, z_{n})\,
    \frac{dt \wedge \overline{dt}}{t - z_{j}} .
\]
On introduit ensuite, à partir des fonctions symétriques élémentaires
$p_{k} = \sum x_{1}x_{2}\cdots x_{k}$, la combinaison
\[
  \varphi_{n}(x_{1}, \ldots, x_{n}) \;=\; \frac{1}{n}\left(1 +
  \frac{p_{1}}{n-1} + \frac{p_{2}}{\binom{n-1}{2}} + \cdots + p_{n-1}\right) ,
\]
et l'on note $\varphi_{n}(U ; j_{1}, \ldots, j_{q})$ ce que devient
$\varphi_{n}(U_{1}, \ldots, U_{n})$ quand on remplace par zéro les
$U_{j_{1}}, \ldots, U_{j_{q}}$. L'opérateur d'homotopie, pour une forme
$\omega = a\, d\bar{z}_{j_{1}} \wedge \cdots \wedge d\bar{z}_{j_{q}}$ avec
$q \geqslant 1$, est alors
\[
  P\omega \;=\; \varphi_{n}(U ; j_{1}, \ldots, j_{q}) \ast
  \sum_{\gamma=1}^{q} (-1)^{\gamma-1} (T_{j_{\gamma}} a)\;
  d\bar{z}_{j_{1}} \wedge \cdots \wedge
  \widehat{d\bar{z}_{j_{\gamma}}} \wedge \cdots \wedge d\bar{z}_{j_{q}} ,
\]
et il vérifie
\[
  P\, d'' f \;=\; f - U_{1}U_{2}\cdots U_{n}\, f ,
\]
c'est-à-dire : \emph{$f$ moins sa partie holomorphe est un bord}. C'est le
lemme de Dolbeault--Grothendieck, avec son opérateur écrit
explicitement.\footnote{Le signe qui sépare $\varphi_{n}$ de la somme est un
astérisque appuyé que la page n'explique pas ; la transcription le signale et
cette lecture ne lui suppose pas de sens. La combinaison $\varphi_{n}$ est ce
qui rend l'opérateur indépendant de l'ordre choisi sur les indices.}

\pagerange{190}{191}
\section*{Domaines de ramification : les axiomes R1 à R4 (pages 190 et 191)}

Grothendieck pagine cette suite 1 à 10, puis 11 à 16 aux pages 201 à 206.

\subsection*{La définition}

Soit $S$ un schéma, $\mathrm{Et}(S)$ son site étale et $\mathcal{U}$ un
univers. Un \emph{domaine de ramification} sur $S$ est un couple $(R, r)$ formé
d'un \emph{champ en $\mathcal{U}$-catégories multigaloisiennes} $R$ sur
$\mathrm{Et}(S)$ et d'un morphisme \emph{continu}
\[
  r : R \longrightarrow \mathrm{Fl}_{S} ,
\]
appelé \emph{réalisation géométrique}, où $\mathrm{Fl}_{S}$ est la catégorie
fibrée sur $\mathrm{Et}(S)$ dont la fibre en $S'$ est $(\mathrm{Sch})_{/S'}$.
Un objet $X$ de $R_{S'}$ est un \emph{revêtement de type $R$ sur $S'$}, et
$r(X)$ est sa réalisation géométrique.

\begin{quote}
\textbf{Il faut s'arrêter ici.} Le carton emploie les mêmes mots pour deux
objets différents. La \emph{donnée} de ramification de 1967 --- notée
$\mathcal{R} = (\mathcal{C}, s)$ plus haut --- est un
catalogue auxiliaire : une catégorie fibrée quelconque, avec un foncteur vers
les schémas, dont on \emph{déduit} la catégorie des revêtements. Le
\emph{domaine} de ramification $(R, r)$ de cette page \emph{est} la catégorie
des revêtements, prise comme donnée, et c'est un champ. Les deux ne sont pas de
même nature, et le second est ce que la page 24 annonçait qu'on serait
« infailliblement conduit » à prendre.
\end{quote}

Les axiomes sont les suivants.

\begin{itemize}
\item[\textbf{R1}] $r$ commute aux \emph{sommes finies} et au \emph{passage au
  quotient} par une relation d'équivalence. Il n'est en général \emph{pas}
  exact à gauche.

\item[\textbf{R2}] Pour $S' \in \mathrm{Et}(S)$ et $X \in R_{S'}$,
  l'application $X' \mapsto r(X')$ est une \emph{bijection} de l'ensemble des
  sous-objets directs de $X$ sur l'ensemble des parties ouvertes et fermées de
  $r(X)$.

\item[\textbf{R3}] Pour une famille $(X_{i})$ --- pas nécessairement finie ---
  de sous-objets de $X$ telle que les $r(X_{i})$ recouvrent $r(X)$, la famille
  $X_{i} \to X$ est une famille de descente \emph{effective}.

\item[\textbf{R4}] Pour $X \to Y$ dominant dans $R_{S}$, on veut que les
  $r(X_{i})$ recouvrent $r(Y)$.\footnote{L'axiome R4 est le plus abîmé des
  quatre sur la page ; la transcription en laisse une partie illisible. Il est
  donné ici tel qu'il se lit, et non complété. Sa fonction se déduit néanmoins
  de l'usage qui en est fait : R1 et R3 servent à faire de la surjectivité une
  topologie, R2 à identifier sous-objets et parties ouvertes et fermées, et R4
  à assurer que la surjectivité se propage --- c'est ce que la page 191 note :
  « on a utilisé seulement R1 ; R3 signifie que les topologies sur $R_{S'}$
  sont localement finies ; R4 que \dots\ L'axiome sur le changement de base est
  trivial. »}
\end{itemize}

Les marges de la page portent trois desiderata qui ne sont pas des axiomes mais
disent ce qu'on voudrait : que les foncteurs de changement de base soient
\emph{exacts} ; que dans $\mathrm{Fl}_{S}$ le quotient soit déterminable ; et
que la famille des objets de $R_{S'}$ soit localement constante finie.

\subsection*{Proposition 1 : c'est une topologie}

Disons qu'une famille de flèches $X_{i} \to X$ de $R_{S'}$ est
\emph{couvrante} si la famille $r(X_{i}) \to r(X)$ est surjective ---
c'est-à-dire, grâce à R1 et « c'est Kif Kif », si les images $X_{i} \subset X$
sont telles que les $r(X_{i})$ recouvrent $r(X)$.

\textbf{Proposition 1.} Ces familles sont les familles couvrantes d'une
\emph{topologie} sur $R_{S'}$, et le foncteur de changement de base
$R_{S'} \to R_{S''}$ est \emph{continu}.

La vérification des quatre axiomes d'une topologie est page 191, et elle est
courte :

\begin{enumerate}
  \item[a)] une famille majorée par une famille couvrante est couvrante ---
    trivial ;
  \item[b)] une somme est couvrante --- trivial ;
  \item[c)] stabilité par changement de base : grâce à b) on se ramène au cas
    où les $X_{i} \to X$ sont des monomorphismes ; mais alors $r$ commute aux
    produits fibrés $X_{i} \times_{X} X'$, et l'on gagne ;
  \item[d)] nature locale : si des $X'_{\alpha} \to X$ couvrent $X$ et que pour
    chaque $\alpha$ les $X_{i} \times_{X} X'_{\alpha}$ couvrent
    $X'_{\alpha}$, alors les $X_{i}$ couvrent $X$ --- même réduction aux
    monomorphismes.
\end{enumerate}

Le mécanisme est le même dans c) et d), et il vaut d'être nommé : \emph{$r$
n'est pas exact à gauche en général, mais il l'est sur les monomorphismes},
et b) permet toujours de s'y ramener. C'est là tout le contenu de R1 et R2.

\textbf{Corollaire.} Prenant les topos associés aux $R_{S'}$, on trouve un
\emph{champ en topos} $\tilde{R}_{S'}$ au-dessus de la sous-catégorie des
objets localement constants finis de $R_{S'}$.

Un bloc rayé de plusieurs diagonales pose ensuite une question à laquelle la
page ne répond pas : pour $X \to Y$ avec $Y = \coprod_{n} Y_{n}$ et chaque
$X \times_{Y} Y_{n} \to Y_{n}$ de somme effective, la famille des
$Y_{n} \to Y$ est-elle couvrante ? Elle est laissée ouverte.

\pagerange{192}{194}
\section*{Objets subordonnés, et l'exemple 1 de 1967 retrouvé (pages 192 à 194)}

\subsection*{Vide et connexe}

\textbf{Proposition 2} (utilise R1 et R2). Pour $X \in R_{S'}$ :
\[
  X = \varnothing \;\Longleftrightarrow\; r(X) = \varnothing ,
  \qquad
  X \text{ connexe} \;\Longleftrightarrow\; r(X) \text{ connexe} .
\]

\subsection*{Subordination}

Soit $C$ une catégorie multigaloisienne et $P$ un objet. Un objet $X$ de $C$
est dit \emph{subordonné} à $P$ si $X_{P} \to P$ est \emph{constant} --- au
sens où $X_{P}$ est somme de copies de $P$.\footnote{La définition est abîmée
sur la page 192 et la transcription y laisse deux passages illisibles. Elle est
fixée ici par l'usage qu'en fait la page 194, où les objets subordonnés à $P$
sont mis en bijection avec les objets « constants par morceaux » sur $r(P)$ ;
c'est cette phrase-là qui décide, non le tracé de la page 192. La condition de
connexité que la page ajoute est celle sous laquelle « constant » et
« constant par morceaux » coïncident.}

Pour $P$ fixé, les objets subordonnés à $P$ forment une sous-catégorie
\emph{pleine} stable par limites projectives finies et par facteurs directs ---
car les objets constants par morceaux sur $P$ forment une telle sous-catégorie
de $C_{/P}$ --- donc une \emph{catégorie multigaloisienne pleine} $C^{(P)}$
de $C$.

Prenant $C = R_{S'}$ et les objets \emph{localement} subordonnés à $P$, on
obtient une sous-catégorie fibrée pleine multigaloisienne $R^{P}$ de $R$.

\subsection*{R2 équivaut à la fidélité de $r$}

Un feuillet inséré entre ses pages 3 et 4, et rayé de plusieurs diagonales,
porte une proposition qu'il vaut la peine de retenir :

\textbf{Proposition.} Sous R1, R3 et R4, l'axiome \textbf{R2 équivaut à ce que
$r$ soit \emph{fidèle}}.

\emph{Démonstration.} Supposons $r$ fidèle. Soient $X'$, $X''$ deux sous-objets
de $X$ avec $r(X') \supset r(X'')$ dans $r(X)$. Comme
$r(X' \cap X'') = r(X') \cap r(X'')$, les inclusions
$X' \cap X'' \subset X'$ et $X' \cap X'' \subset X''$ deviennent des
isomorphismes par $r$, ce qui ramène au cas d'une inclusion
$X' \subset X''$, soit $X'' = X' \sqcup X_{1}$. Alors
$r(X'') = r(X') \sqcup r(X_{1})$, et $r(X') = r(X'')$ équivaut à
$r(X_{1}) = \varnothing$, donc à $X_{1} = \varnothing$.

Cette proposition mérite d'être mise en regard des pages 12 à 20. En 1967, la
\emph{non-}pleine-fidélité de la réalisation géométrique était le fait
fondamental --- « il importera de bien distinguer entre un revêtement de type
$\mathcal{R}$ et sa réalisation géométrique » --- et la \emph{fidélité} était
une propriété qu'on vérifiait au cas par cas sur les exemples. Ici, elle est
promue au rang d'axiome, sous la forme R2, et c'est le seul axiome des quatre
qui porte sur ce que $r$ ne perd pas.

\subsection*{Proposition 3, et l'exemple 1 de 1967}

\textbf{Proposition 3.} Soit $P$ un objet \emph{galoisien de groupe $G$} dans
$R_{S'}$. Le foncteur $X \mapsto r(X \times P)/r(P)$ est une \emph{équivalence}
de la catégorie des objets de $R_{S'}$ subordonnés à $P$ avec la catégorie des
objets constants par morceaux sur $r(P)$ munis d'une action de $G$ compatible
avec celle qu'il a sur $r(P)$.

La démonstration est en deux temps : utilisant R2, le foncteur est pleinement
fidèle ; utilisant R2 et R3, il est essentiellement surjectif. Et comme $G$
opère sur $R$ il opère sur $r(P)$, l'équivalence étant compatible.

\textbf{Corollaire 1.} Soit $P$ galoisien de groupe $G$ dans $R_{S}$. Le même
foncteur induit une équivalence de la catégorie des objets de $R_{S}$
\emph{localement} subordonnés à $P$ avec la catégorie des revêtements étales de
$Z = r(P)$ munis d'une action de $G$ compatible, satisfaisant la condition
suivante --- automatiquement vérifiée si $Z$ est fini sur $S$ : les $S' \to S$
étales tels que $Z_{S'}$ y devienne constant par morceaux recouvrent $S$.

Grothendieck ajoute lui-même, après un trait oblique : « Localisant sur $S$, on
trouve une description du champ $R^{P}$ en termes des schémas $Z$ à groupe
d'opérateurs $G$ sur $S$. »

\begin{quote}
\emph{C'est l'exemple 1 du tapuscrit de 1967}, page 14 du carton : un schéma
$Z$ sur $S$ muni d'un groupe fini $G$ d'opérateurs, la catégorie des
revêtements étales de $Z$ à opérateurs, et $r(Z') = Z'/G$. Ce qui était le
premier et le plus simple des trois exemples devient ici la description du
sous-champ engendré par un objet galoisien. L'exemple est devenu un théorème.
\end{quote}

\pagerange{195}{196}
\section*{L'exemple 2, et pourquoi l'exemple 3 devait exister (pages 195 et 196)}

\subsection*{Proposition 4 : la réciproque}

\textbf{Proposition 4.} Pour tout schéma $Z$ sur $S$ muni d'un groupe fini $G$
d'opérateurs agissant de façon \emph{admissible}, le champ décrit ci-dessus est
un champ en $\mathcal{U}$-catégories multigaloisiennes, et le foncteur
« passage au quotient par $G$ » --- qui est défini, l'action étant admissible
--- en fait un \emph{domaine de ramification} sur $S$.

La vérification est brève, et il faut la lire pour voir combien les axiomes
sont peu coûteux. Que $R_{S'}$ soit multigaloisienne : on part de la catégorie
des revêtements à opérateurs $G$ de $Z_{S'}$, dont on prend une sous-catégorie
pleine stable par limites projectives finies et par facteurs directs. Que R1
soit vérifié : $r$ est un \emph{passage au quotient}, donc commute aux limites
inductives finies, calculées au sens des schémas. R2, R3 et R4 sont
immédiats.\footnote{Le mot que la transcription lit « donnée » dans la
conclusion de la proposition 4 pourrait être « domaine » --- c'est le terme que
la page 190 définit et que la chemise porte en titre --- la lettre ne tranchant
pas et l'article féminin allant contre. C'est le seul endroit du carton où le
choix entre les deux mots a une conséquence, puisqu'il s'agit précisément de
dire que la construction produit l'objet nouveau et non l'objet auxiliaire de
1967. Le contexte impose « domaine » et c'est ce qui est écrit ici.}

\subsection*{L'exemple 2}

Supposons $S$ affine et $R$ soumis à « une petite condition supplémentaire »
vérifiée en pratique. Alors tout objet de $R_{S}$ est subordonné à un objet
galoisien à opérateurs $(P, G)$ de $R_{S}$, de sorte que
\[
  R_{S} \;=\; \varinjlim_{P} R^{P}_{S}
\]
est \emph{limite inductive filtrante} de sous-catégories, et l'on obtient une
description de $R_{S}$ en termes d'un \emph{système projectif} d'objets à
opérateurs $(Z_{i}, G_{i})$.

\begin{quote}
\emph{C'est l'exemple 2 du tapuscrit de 1967}, page 16 du carton : un système
projectif $(Z_{i}, G_{i})$ de schémas à opérateurs, et
$\mathrm{Rev}^{\mathcal{R}}(S) = \varinjlim_{i}
\mathrm{Rev}^{\mathcal{R}_{i}}(S)$. Là encore, l'exemple est devenu une
description.
\end{quote}

Si l'objet final de $R_{S}$ est \emph{connexe}, on peut prendre les $P_{i}$
connexes ; les morphismes de transition $G_{j} \to G_{i}$ sont alors
\emph{surjectifs} et $Z_{i} \simeq Z_{j}/H_{ij}$, et les $Z_{i}/G_{i}$ sont
connexes. Ce sont exactement les conditions sous lesquelles, en 1967, les
foncteurs de transition de l'exemple 2 étaient pleinement fidèles.

C'est ici que la page annonce, relativement à un tel système projectif, « un
énoncé comme précédemment » sous le nom de \textbf{Proposition 5} --- et
n'énonce rien. Le numéro est suivi de points de suspension et la page passe à
autre chose.\footnote{La page 196 porte « \textbf{Prop.\ 5} \dots » et
c'est tout. Ce qu'elle aurait dit se devine de la phrase qui la précède ---
l'analogue, pour un système projectif $(Z_i, G_i)$, de la description que la
proposition 3 et son corollaire donnent pour un seul objet galoisien --- mais
se deviner n'est pas être écrit, et cette lecture ne l'écrit pas à sa place.
C'est le second endroit de la suite où un énoncé est annoncé et manque ; le
premier est à la page 206, où l'essentielle surjectivité de $\alpha$ est
annoncée et sa démonstration ne suit pas.}

\subsection*{Et pourquoi il faut recoller}

Vient alors l'observation qui justifie tout ce qui suit :

\begin{quote}
Partant de $R$ et construisant les $(Z_{i}, G_{i})$, on obtient un nouveau
$(R', r')$ et une inclusion $R' \hookrightarrow R$ ; il n'est \emph{pas} vrai
en général que ce soit une équivalence : \emph{le champ $R$ n'est pas
nécessairement engendré par ses sections globales}. Donc pour définir un $R$ en
général, il faut localiser et \emph{recoller}.
\end{quote}

\begin{quote}
\emph{C'est la raison d'être de l'exemple 3 de 1967}, et c'est ici seulement
qu'elle est dite. La page 35 du carton l'expliquait déjà en termes concrets ---
« sometimes the ramification models are not defined globally (for instance in
the Kummer case, when we give divisors which are not principal), only locally »
--- mais sans énoncé. L'exemple 3, avec ses quatre conditions, puis ses cinq
conditions a) à e) de la page 32, était l'appareil de recollement écrit à la
main. Ici, la nécessité du recollement est un fait sur le champ $R$, et le
recollement lui-même va être fait une fois pour toutes.
\end{quote}

\subsection*{Les homomorphismes}

Un \emph{homomorphisme} de domaines de ramification $(R, r) \to (R', r')$ est
un couple $(\varphi, \lambda)$ où $\varphi : R \to R'$ est un foncteur
\emph{continu et exact} et $\lambda$ un isomorphisme de foncteurs continus
$r'\varphi \simeq r$.

\pagerange{197}{198}
\section*{La 2-catégorie $\mathrm{Domram}(S)$, et le recollement (pages 197 et 198)}

\subsection*{Une 2-catégorie}

Les domaines de ramification sur $S$ forment une \emph{2-catégorie},
$\mathrm{Domram}(S)$ ; et pour $S$ variable, ces 2-catégories forment une
\emph{2-catégorie fibrée}.

C'est ici que se règle ce devant quoi la page 27 avait reculé : « On recule
devant la tâche d'expliciter une propriété de transitivité pour un composé
$U \to T \to S$, et jetant provisoirement un voile pudique sur ces abîmes ».
Les abîmes étaient la 2-catégorialité, et elle est ici assumée comme une
structure plutôt que subie comme une complication.

\subsection*{Proposition 6}

Si $(R, r)$ est défini par $(Z, G)$, on dispose dans $R_{S}$ d'un objet
\emph{galoisien canonique} $P$ de groupe $G$. Un homomorphisme
$(R, r) \to (R', r')$ définit alors un objet $\varphi(P)$ de $R'$, galoisien de
groupe $G$, et un isomorphisme compatible à $G$
\[
  r'(\varphi(P)) \;\overset{\alpha}{\simeq}\; Z .
\]

\textbf{Proposition 6.} Cela donne une \emph{équivalence} entre
$\underline{\mathrm{Hom}}\bigl((R,r),(R',r')\bigr)$ et la catégorie des couples
$(P', \alpha)$ où $P'$ est un objet galoisien de $R'_{S}$ de groupe $G$ et
$\alpha$ un $G$-isomorphisme $r'(P') \simeq Z$.

\textbf{Corollaire.} Dans tout homomorphisme $(\varphi, \lambda)$ de domaines
de ramification, $\varphi$ est \emph{nécessairement pleinement fidèle}.

C'est un énoncé fort, et il faut le mesurer : \emph{il n'y a pas de morphisme
non trivial qui écrase}. Un domaine de ramification ne s'envoie dans un autre
que comme sous-domaine, et le corollaire suivant dit lequel.

\textbf{Corollaire.} Si $(R', r')$ est défini par $(Z', G')$, alors
$\underline{\mathrm{Hom}}\bigl((R,r),(R',r')\bigr)$ est équivalent à la
catégorie des revêtements \emph{principaux} $P'$ de $Z'$ de groupe $G$, sur
lesquels $G'$ opère compatiblement avec ses opérations sur $Z'$, munis d'un
isomorphisme de $G$-schémas sur $S$
\[
  P'/G' \;\simeq\; Z .
\]

\subsection*{Le recollement}

Soient des $S'_{i}$ couvrant $S$, sur chacun un schéma à groupe d'opérateurs
admissible $(Z_{i}, G_{i})$ ; sur chaque $S_{i} \times_{S} S_{j}$ une
\emph{subordination} de $(Z_{ij}, G_{i})$ à $(Z_{ji}, G_{j})$, où
$Z_{ij} = Z_{i} \times_{S_{i}} (S_{i} \times_{S} S_{j})$ ; sur les triples, des
isomorphismes de transitivité ; et enfin une condition de compatibilité sur les
triples. Alors on peut \emph{recoller} les $(R_{i}, r_{i})$.

\subsection*{La construction générale, par les gerbes}

Soit $C$ une catégorie \emph{fibrée} de « modèles » sur $\mathrm{Et}(S)$ ---
pas nécessairement un champ --- et $s : C \to \mathrm{Fl}_{S}$ un foncteur
continu. On suppose :

\begin{itemize}
  \item $C$ fibrée en \emph{groupoïdes} ;
  \item deux objets d'un $C_{S'}$ localement isomorphes ;
  \item $C_{S'}$ localement non vide.
\end{itemize}

Le champ associé est alors une \emph{gerbe}. On suppose de plus que le
\emph{lien} de cette gerbe est \emph{localement constant fini} : pour tout
$x \in C_{S'}$, le faisceau associé au préfaisceau $S'' \mapsto
\operatorname{Aut}(x_{S''})$ est représentable par un revêtement étale de
$S'$.

On a alors une « représentation » de cette gerbe dans le préchamp
$\mathrm{Fl}_{S}$, et l'on associe au tout un domaine de ramification :

\begin{tikzcd}[column sep=small, row sep=small]
& \mathrm{Rev}\,\mathrm{Fl}_S \arrow[d] \\
\mathcal{G} \arrow[r] \arrow[d] & \mathrm{Fl}_S \\
\mathrm{Et}(S) &
\end{tikzcd}

\emph{les objets de $R_{S'}$ sont les homomorphismes continus
$\mathcal{G} \to \mathrm{Rev}\,\mathrm{Fl}_{S}$ qui relèvent
$\mathcal{G} \to \mathrm{Fl}_{S}$.}

\begin{quote}
C'est, littéralement, la définition de 1967 : la page 12 posait une catégorie fibrée $\mathcal{C}$ sur $\mathrm{Et}(S)$, un foncteur cartésien $s$ vers les
schémas relatifs, et définissait $\mathrm{Rev}^{\mathcal{R}}(S)$ comme la
catégorie des \emph{sections cartésiennes} de l'image inverse de la catégorie
fibrée des revêtements étales. Ce sont les mêmes objets, et le mot
« relèvement » est le mot juste pour « section cartésienne ». Ce qui a changé
tient en deux mots : $\mathcal{C}$ est maintenant une \emph{gerbe} à lien
localement constant, et le résultat est un \emph{champ}. La normalisation
demandée page 24 est faite.
\end{quote}

\pagerange{199}{200}
\section*{Les exemples kummériens, et la classe une quatrième fois (pages 199 et 200)}

\subsection*{Exemple 1 : des équations globales}

Donnons-nous une famille $\underline{a} = (a_{i})_{i \in I}$ de sections de
$\mathcal{O}_{S}$ et une famille $\underline{n} = (n_{i})_{i \in I}$ d'entiers
$\geqslant 1$. On prend pour $C_{S'}$ le groupoïde correspondant, avec
\[
  s(C_{S'}) \;=\; Z^{\underline{a}}_{\underline{n}} , \qquad
  \mu_{\underline{n}}(S') \;=\; \prod_{i \in I} \mu_{n_{i}}(S') ,
\]
et l'opération qui en dérive. Les axiomes sont vérifiés. Ici la gerbe est
\emph{triviale} --- il y a une section globale, à savoir $\underline{a}$
elle-même --- et le domaine se réduit aux revêtements étales de
$Z^{\underline{a}}_{\underline{n}}$ sur lesquels $\mu_{\underline{n}}$ opère
compatiblement.

\textbf{Exemple 1 bis.} La limite sur $\underline{n}$ : le \emph{domaine de
ramification kummérien} défini par $\underline{a}$.

\subsection*{Exemple 2 : des diviseurs, et la gerbe}

Donnons-nous une famille $(D_{i})_{i \in I}$ de \emph{diviseurs} et une famille
$(n_{i})$ d'entiers. On prend
\[
  \operatorname{Ob}(C_{S'}) \;=\; \bigl\{\, \underline{a} = (a_{i})_{i \in I}
  \;:\; a_{i} \text{ est une équation de } D_{i} \,\bigr\} ,
\]
\[
  \operatorname{Hom}(\underline{a}, \underline{b}) \;=\;
  \bigl\{\, \xi = (\xi_{i})_{i \in I} \;:\; b_{i} = a_{i}\,\xi_{i}^{\,n_{i}}
  \,\bigr\} ,
\]
la composition et les images inverses étant évidentes. Le groupe des
automorphismes d'un objet est $\{\xi : \xi_{i}^{n_{i}} = 1\} =
\mu_{\underline{n}}$, et le champ associé est donc une \emph{gerbe abélienne de
groupe $\mu_{\underline{n}}$}. Elle est localement non vide --- les équations
existent localement --- et sa classe obstrue l'existence d'une section globale.
Grothendieck la nomme :

\begin{quote}
L'existence d'une section dépend de la \emph{nullité} d'un élément de
$H^{2}(X, \mu_{\underline{n}})$, qui n'est autre, comme on devine, que la
\emph{classe de cohomologie des $D_{i}$} dans les $H^{1}(X, \mu_{n_{i}})$ !
\end{quote}

C'est le théorème des pages 87 à 97, énoncé pour la quatrième fois dans ce carton, et
sous sa forme la plus nue.\footnote{Les trois autres : le lemme 2 de la page 89
(la classe de l'extension des groupes fondamentaux est celle de
l'auto-intersection de la fibre spéciale) ; le programme de la page 100 (pour
une famille quelconque de diviseurs de Cartier, $\xi_{i} =
c(\mathcal{O}_{X}(D_{i})) \bmod n_{i}$) ; et le corollaire de la page 102. Les
deux exposants de $H$ sur cette page sont tracés de la même main, qui ne
distingue pas le 1 du chiffre romain ; le second est un 2 : la classe d'un
fibré en droites est dans $H^{1}(\mathbb{G}_{m})$, et sa classe de Chern dans
$H^{2}(\mu_{n})$ par la suite de Kummer.}

\textbf{Exemple 2 bis.} La limite sur $\underline{n}$ : le \emph{domaine de
ramification kummérien} défini par $D = (D_{i})$.

\subsection*{Ce que ces exemples sont, aujourd'hui}

La gerbe de l'exemple 2, avec les revêtements étales de son revêtement
kummérien équivariants sous $\mu_{\underline{n}}$, est ce qu'on appelle
aujourd'hui le \emph{champ des racines}
$\sqrt[\underline{n}]{\underline{D}/S}$ : le champ dont les points sont les
choix de racines $n_{i}$-ièmes des équations des $D_{i}$, et dont la gerbe
d'inertie le long de $D_{i}$ est $\mu_{n_{i}}$. Les revêtements de type $R$ ne
sont alors rien d'autre que les revêtements finis étales de ce champ, et le
groupe $\pi_{1}^{\mathcal{R}}(S)$ construit à la page 25 du carton est son
groupe fondamental.\footnote{Les champs des racines sont de Charles Cadman et
d'Abramovich, Graber et Vistoli, dans les années 2000 ; les gerbes sont de Jean
Giraud (1971) ; et l'énoncé que les revêtements modérément ramifiés d'un couple
$(S, D)$ à croisements normaux sont les revêtements étales d'un champ, ou
--- ce qui revient au même --- les revêtements Kummer-étales du schéma
logarithmique associé, appartient à la géométrie logarithmique de Kato, Illusie
et Nakayama, dans les années 1990. Aucun de ces mots n'est disponible ici, et
rien dans cette note n'affirme de filiation : elle constate que l'objet, sa
gerbe, sa classe d'obstruction et la description de ses revêtements sont
construits dans ces pages.}

\subsection*{Le boitement, avoué}

La page se termine sur une réserve : « Ces exemples boitent un tant soit peu ;
on voudrait au fait que les $n_{i}$ ne soient pas premiers avec les
[caractéristiques] résiduelles des $D_{i}$ seulement \dots ». C'est-à-dire :
la condition de modération devrait porter sur chaque $D_{i}$ séparément, et non
globalement sur $S$. C'est exactement l'assouplissement que la donnée kummérienne des pages 19 et 20 opérait en laissant les $u : G_{S'} \to
\mu_{\underline{n},S'}$ varier localement --- et les pages 201 à 206 vont reprendre la question sous le titre des \emph{familles génératrices}.

La page s'interrompt en cours de phrase, sur « Dans le cas de l'exemple 2 ». La suite est à la page 201, et elle ne revient pas sur les exemples.

\pagerange{201}{202}
\section*{Familles génératrices, et la donnée de 1967 retrouvée (pages 201 et 202)}

Grothendieck pagine cette suite 11 à 16, à la suite du 1--10 des pages 190 à 200.

\subsection*{Ce qu'est une famille génératrice}

Soit $(R, r)$ un domaine de ramification sur $S$. Une famille d'objets
$X_{i}/S_{i}$ de $R$, avec $S_{i} \in \mathrm{Et}(S)$, est dite
\emph{génératrice} si, pour tout objet $X$ de $R$ sur un $S'$, les $S''$ étales
sur $S'$ pour lesquels il existe un $i$, un morphisme $S'' \to S_{i}$ et une
subordination de $X_{S''}$ à $X_{i,S''}$, recouvrent $S'$.

En clair : \emph{tout revêtement est localement subordonné à un des
générateurs}. C'est la même condition que le crible couvrant de 1967, dite du
côté du champ.

\subsection*{Le refus de la présentation naïve}

On pourrait vouloir reconstituer $(R, r)$ à partir des seuls
$(Z_{i}, G_{i}) = (r(X_{i}), \operatorname{Aut}(X_{i}))$, en précisant les
compatibilités. Grothendieck écarte ce parti, et son jugement mérite d'être
retenu tel quel : « la formulation de telles données est \emph{pénible et
finalement peu utile} ».\footnote{La phrase de la page se défait entre « Si de
plus » et « les $Z_{i}$, $G_{i}$ doivent permettre » ; ce qui se lit est le
projet et son abandon. C'est le troisième refus de ce genre dans le carton,
après le « voile pudique » de la page 27 et le « pénible » implicite des cinq
conditions a) à e) de la page 32 --- et les trois portent sur le même objet :
l'explicitation en coordonnées de ce qui se dit sans coordonnées.}

\subsection*{Le couple $(C, s)$, et où le carton se referme}

Il lui préfère d'associer à $R$ une catégorie $C$ sur $\mathrm{Et}(S)$ et un
foncteur
\[
  \varphi : C \longrightarrow R ,
\]
dit \emph{générateur} si les $\varphi(M)$ forment une famille génératrice ---
c'est-à-dire si, pour tout $X \in R$ sur $S'$, les $S''$ étales sur $S'$ tels
que $X_{S''}$ soit subordonné à un objet de $C_{S''}$ recouvrent $S'$. Et il
pose
\[
  s \;=\; r\varphi \;:\; C \longrightarrow \mathrm{Fl}_{S} ,
\]
en annonçant qu'on peut alors, à équivalence près, reconstituer $(R, r)$ en
termes de $(C, s)$ --- et de même les $(R_{S'}, r_{S'})$, avec toutes les
fonctorialités.

\begin{quote}
\textbf{Le couple $(C, s)$ est la donnée de ramification de 1967.} Le tapuscrit
de la page 12 posait : « une donnée de ramification sur $S$ est un couple
$(\mathcal{C}, s)$, où $\mathcal{C}$ est une catégorie fibrée sur
$\mathrm{Et}(S)$ et $s : \mathcal{C} \to \mathrm{Fl}_{S}$ un foncteur
cartésien ». C'est le même couple, avec le même nom pour le second membre. Ce
qui a changé est son statut : en 1967 c'était la donnée dont on \emph{déduisait}
la catégorie ; ici c'est une \emph{présentation} de la catégorie, qui existe
d'abord, et dont on demande si elle la détermine.
\end{quote}

\subsection*{Les sections cartésiennes}

Soit $X$ un objet de $R_{S}$, et $C_{X}$ la sous-catégorie pleine de $C$ formée
des objets $M$ sur un $S'$ tels que $X_{S'}$ soit subordonné à $\varphi(M)$ et
que $\varphi(M) \to X_{S'}$ soit un épimorphisme. Pour un tel $M$,
\[
  r\bigl(X_{S'} \times \varphi(M)\bigr) \longrightarrow r\varphi(M) = s(M)
\]
est un \emph{revêtement étale}. Donc
\[
  M \longmapsto r\bigl(X \times_{S'} \varphi(M)\bigr) \big/ s(M)
\]
donne une \emph{section cartésienne} de la catégorie fibrée des revêtements
étales relatifs, tirée en arrière sur $C_{X}$ par
$C_{X} \xrightarrow{\ s\ } \mathrm{Fl}_{S}$.

\begin{quote}
\textbf{Ce sont les revêtements de type $\mathcal{R}$ de 1967.} La page 13 du
carton posait : « on appellera revêtement de type $\mathcal{R}$ de $S$ une
section cartésienne de $\mathcal{D}$ sur $\mathcal{C}$ », $\mathcal{D}$ étant
l'image inverse par $s$ de la catégorie fibrée des revêtements étales finis.
C'est mot pour mot la même construction, dans l'autre sens : là on partait de
$\mathcal{C}$ pour \emph{fabriquer} les sections cartésiennes, ici on part d'un
objet $X$ du champ pour lui \emph{associer} une section cartésienne.
\end{quote}

\subsection*{Les sous-catégories $R_{S}^{C_{0}}$}

Pour un crible \emph{admissible} $C_{0}$ de $C$, soit $R_{S}^{C_{0}}$ la
sous-catégorie \emph{pleine} de $R_{S}$ formée des objets $X$ tels que, pour
tout $S' \in \mathrm{Et}(S)$, $X_{S'}$ soit localement subordonné à \emph{tout}
objet de $C_{0}$ sur $S'$. C'est une sous-catégorie multigaloisienne de
$R_{S}$, stable par limites projectives et par sommes directes, et l'on
dispose d'un foncteur
\[
  (\ast) \qquad
  \alpha : R_{S}^{C_{0}} \longrightarrow
  \varprojlim_{C_{0}} \bigl(\mathcal{E}\,|\,C_{0}\bigr) ,
\]
$\mathcal{E}$ désignant la catégorie fibrée des revêtements étales
relatifs.\footnote{Le symbole que la page emploie pour cette catégorie fibrée
est un E de fantaisie qu'elle ne nomme nulle part ; la transcription le rend
$\mathcal{E}$ partout où il paraît, et cette lecture fait de même. C'est le
$\mathcal{D}$ de la page 12, pris au-dessus de $C_{0}$.}

\pagerange{203}{206}
\section*{$\alpha$ est pleinement fidèle ; et la suite manque (pages 203 à 206)}

\subsection*{La formule visée}

Faisons parcourir à $C_{0}$ les cribles admissibles de $C$, et supposons la
catégorie fibrée $C$ \emph{localement filtrante} --- ce qui entraîne que
l'intersection de deux cribles couvrants est couvrante. Alors les
$R_{S}^{C_{0}}$ forment une famille \emph{filtrante} de sous-catégories
multigaloisiennes pleines de $R_{S}$, de réunion $R_{S}$. Donc, si les
$\alpha_{0}$ sont des \emph{équivalences}, on aura
\[
  R_{S} \;\simeq\; \varinjlim_{C_{0}} \; \varprojlim_{C_{0}}
  \bigl(\mathcal{E}\,|\,C_{0}\bigr) ,
\]
formule qui explicite $R_{S}$ en termes de $(C, s)$.

\begin{quote}
\textbf{C'est la formule de 1967.} La page 13 du carton posait
$\mathrm{Rev}^{\mathcal{R}}(S) = \varinjlim_{\mathcal{C}_{0}}
\mathrm{Rev}^{\mathcal{R}}_{\mathcal{C}_{0}}(S)$, où
$\mathrm{Rev}^{\mathcal{R}}_{\mathcal{C}_{0}}(S)$ était la catégorie des
sections cartésiennes au-dessus de $\mathcal{C}_{0}$, c'est-à-dire précisément
$\varprojlim_{\mathcal{C}_{0}}$. Une limite inductive sur les cribles, de
limites projectives sur chaque crible : c'est la même formule, aux deux
extrémités du carton. Et les deux hypothèses (H1) et (H2) que la lecture des pages 12 à 14 avait dû isoler --- la pleine fidélité des foncteurs de transition,
l'existence de $r$ --- sont maintenant \emph{ce qu'il s'agit de démontrer}.
\end{quote}

Et l'on annonce que pour cela il suffit que $\varphi$ soit \emph{pleinement
fidèle}.

\subsection*{Ce qui est facile}

$\alpha$ est un foncteur exact de catégories multigaloisiennes ; il suffit donc,
pour la fidélité, de voir que
\[
  \alpha(X) = \varnothing \;\Longrightarrow\; X = \varnothing .
\]
C'est immédiat : la question est locale sur le site étale, il existe un
$M \in C_{0,S}$, et
$r(\varphi(M) \times X) = \varnothing$ entraîne
$\varphi(M) \times X = \varnothing$ par la proposition 2 de la page 192.

\subsection*{Ce qui demande du travail : la pleine fidélité}

Soit $X \in R_{S}^{C_{0}}$ et $Z = \alpha(X) = Z_{1} \sqcup Z_{2}$ ; il faut
montrer que cette décomposition provient d'une décomposition de $X$.

La question est locale sur $S$ ; on peut donc supposer qu'il existe
$M \in C_{S}$ avec $\varphi(M) \to e$ épimorphique. Considérons
\[
  r\bigl(X \times \varphi(M)\bigr) \;=\; \alpha(X)(M) \;=\; Z(M)
  \longrightarrow r(\varphi(M)) = s(M) .
\]
Il se coupe en $Z_{1}(M) \sqcup Z_{2}(M)$, donc \emph{par l'axiome R2}
$X \times \varphi(M)$ se coupe en $X_{1} \sqcup X_{2}$.

Reste à montrer que cette décomposition provient d'une décomposition de $X$ :
c'est-à-dire que pour tout objet $T$ de $R_{S}^{C_{0}}$ et deux morphismes
$T \rightrightarrows \varphi(M)$, les deux images inverses de $X_{1}$
coïncident.

\begin{tikzcd}[column sep=small, row sep=small]
T \arrow[d, bend left=12] \arrow[d, bend right=12]
  & X \times T \arrow[d, bend left=12] \arrow[d, bend right=12] & \\
\varphi(M) \arrow[d]
  & X \times \varphi(M) \arrow[l, no head] \arrow[r] \arrow[d] & X_1 \\
e & X \arrow[l, no head] &
\end{tikzcd}

L'argument est le suivant, et c'est le morceau de bravoure de cette suite.

Le morphisme $T \times \varphi(M) \to \varphi(M)$ est un épimorphisme. Or
$T \times \varphi(M)$, étant subordonné à $\varphi(M)$, est \emph{constant par
morceaux} sur $\varphi(M)$ ; et un objet constant par morceaux et épimorphique
est \emph{recouvert}, pour la topologie canonique, par les images des
\emph{sections} $w$ de $T \times \varphi(M)$ sur $\varphi(M)$. On peut donc
remplacer le couple $T \rightrightarrows \varphi(M)$ d'abord par les composés
\[
  T \times \varphi(M) \xrightarrow{\ \mathrm{pr}_{1}\ } T
  \rightrightarrows \varphi(M)
\]
--- ce qui est licite, $\varphi(M) \to e$ étant un épimorphisme --- puis par
les composés $\varphi(M) \xrightarrow{\ w\ } T \times \varphi(M) \to \cdots$.
\emph{On est ainsi ramené au cas $T = \varphi(M)$.} Mais alors, $\varphi$ étant
pleinement fidèle, les deux flèches $M \rightrightarrows M$ sont identiques.

Le mécanisme mérite d'être nommé : \emph{quand tout objet est localement
constant sur un modèle, on peut remplacer un objet quelconque par le modèle
lui-même, parce que les sections le recouvrent}. C'est exactement pour cela que
la subordination avait été définie.

On obtient donc une décomposition
\[
  X \;=\; X_{1} \sqcup X_{2}
\]
uniquement déterminée par la condition $Z_{1} = \alpha(X_{1})$. Il reste à voir
qu'elle ne dépend pas du choix de $M$ : c'est ce que fait la fin de la page 205
et la page 206, en localisant sur $S$ et en utilisant que $C$ est localement
filtrante.\footnote{Cette dernière partie est, sur la page 205, accolée à
gauche et barrée de deux longues diagonales ; elle se lit par endroits et par
endroits non. Le principe --- comparer les décompositions définies par $M$ et
par $M'$ en passant à un $M''$ qui les majore, ce que la filtration locale
fournit --- est ce que les phrases portantes autorisent ; le détail ne l'est
pas, et il n'est pas reconstruit ici.}

\subsection*{Et la suite manque}

La page 206 se termine sur une seule ligne :

\begin{quote}
Prouvons que $\alpha$ est \emph{essentiellement surjectif}.
\end{quote}

\emph{La démonstration ne suit pas.} La page s'arrête là et la suite du carton
est d'un autre sujet.

Il faut donc dire ce qui est acquis et ce qui ne l'est pas. Est acquise la
pleine fidélité de $\alpha$ --- $R_{S}^{C_{0}}$ s'envoie fidèlement et
pleinement dans la limite projective. N'est pas acquise l'essentielle
surjectivité, donc \emph{l'équivalence annoncée à la page 203 reste
conditionnelle}, et avec elle la reconstitution de $(R, r)$ à partir de
$(C, s)$. La boucle que le carton ferme --- de la donnée de 1967 au champ, et
du champ à la donnée --- est fermée dans un sens et pas dans l'autre.

\pagerange{208}{209}
\section*{Catégories multigaloisiennes : les conditions A et B (pages 208 et 209)}

À partir d'ici et jusqu'à la fin, le texte est écrit au net, à l'encre bleue,
et se lit presque de bout en bout --- à la différence de tout ce qui précède
dans ce carton. Grothendieck pagine 1 à 4, insère deux feuillets « 1 bis » et
« 1 ter », puis 5 et 6.

\subsection*{A) Ce que c'est}

Soit $\mathcal{U}$ un univers et $C$ une catégorie. Les conditions suivantes
sont équivalentes.

\begin{enumerate}
  \item[(i)] Il existe un $\mathcal{U}$-topos $\mathcal{T}$ et une équivalence
    $C \simeq \mathcal{T}_{\mathrm{cf}}$, où $\mathcal{T}_{\mathrm{cf}}$ est la
    sous-catégorie pleine des objets \emph{localement constants
    finis}.\footnote{L'indice est « cf », et c'est Grothendieck lui-même qui le
    glose à la ligne au-dessus : « constants finis ». La transcription le
    signale ; c'est cette glose, et non le tracé, qui fixe la lecture.}
  \item[(ii)] $C$ est une $\mathcal{U}$-catégorie dont l'ensemble des classes
    d'isomorphie d'objets est $\mathcal{U}$-petit, et il existe sur $C$ une
    topologie --- moins fine que la topologie canonique --- telle que
    $C \to \tilde{C}$ induise une équivalence
    $C \xrightarrow{\ \sim\ } \tilde{C}_{\mathrm{cf}}$.
  \item[(ii bis)] La même chose, en prenant la topologie \emph{la plus fine de
    la descente effective universelle}.
  \item[(iii)] Pour tout objet $X$ de $C$, les objets $S$ tels que $X_{S}$ soit
    \emph{constant fini} sur $S$ forment une famille de descente effective
    universelle.
  \item[(iv)] Quatre conditions, portées sur un feuillet inséré (page 212).
\end{enumerate}

La condition (i) est la définition conceptuelle : \emph{une catégorie
multigaloisienne est la catégorie des systèmes locaux finis d'un topos}. La
condition (iii) est la définition opératoire : \emph{tout objet devient
constant après un changement de base convenable}. Que les deux coïncident est
tout le contenu de A.

\subsection*{B) Le cas multigaloisien proprement dit}

Les conditions suivantes sont équivalentes.

\begin{enumerate}
  \item[(i)] $C$ satisfait aux conditions de A), et son objet final est somme
    \emph{effective universelle} d'objets connexes.
  \item[(ii)] $C$ satisfait aux axiomes \textbf{G0} à \textbf{G5} ci-dessous.
  \item[(iii)] $C$ est équivalente à une catégorie produit
    $\prod_{i \in I} C_{i}$, les $C_{i}$ étant \emph{galoisiennes}.
\end{enumerate}

Avec la note qui explique le mot : \emph{les $C_{i}$ correspondent aux
composantes connexes $e_{i}$ de l'objet final} ; et la donnée d'une équivalence
$C_{i} \simeq C(\pi_{i})$ avec la catégorie des ensembles finis à opérateurs
d'un groupe profini $\pi_{i}$ correspond à la donnée d'un \emph{foncteur fibre}
de $C_{/e_{i}}$, c'est-à-dire d'un « revêtement universel » de $e_{i}$.

\begin{quote}
Voilà ce que « multigaloisien » veut dire, et c'est la réponse à une question
que le carton posait à sa page 30, en 1967, en employant le mot sans le
définir. Une catégorie galoisienne est celle des ensembles finis sur lesquels
opère \emph{un} groupe profini ; elle exige que l'objet final soit connexe. Une
catégorie multigaloisienne est un \emph{produit} de catégories galoisiennes,
une par composante connexe de l'objet final : elle correspond donc non à un
groupe mais à un \emph{groupoïde} profini, et c'est ce qu'il faut quand on ne
veut pas supposer la base connexe. C'est exactement ce qui manquait au
paragraphe 2 du tapuscrit de 1967, dont le groupe
$\pi_{1}^{\mathcal{R}}(S,s)$ n'existait que sous une hypothèse de connexité du
site.
\end{quote}

\subsection*{Les axiomes G0 à G5}

\begin{itemize}
\item[\textbf{G0}] $C$ est une $\mathcal{U}$-catégorie, et l'ensemble des
  classes d'isomorphie d'objets est $\mathcal{U}$-petit.
\item[\textbf{G1}] Existence des limites projectives \emph{finies} et des
  sommes \emph{finies}.
\item[\textbf{G2}] Les relations d'équivalence sont \emph{effectives}, et les
  épimorphismes sont universels.
\item[\textbf{G3}] Tout morphisme $X \to Y$ se factorise en un
  \emph{épimorphisme} suivi d'un \emph{isomorphisme sur un facteur direct}.
\item[\textbf{G4}] L'objet final est somme \emph{effective universelle}
  d'objets connexes.
\item[\textbf{G5}] Soit $f : X \to Y$ avec $X$ et $Y$ connexes. Alors il existe
  un entier $N$ tel que, pour tout $Y' \to Y$ non vide, \dots\footnote{L'énoncé
  de G5 est le seul de la liste que la page laisse incomplet : la conclusion
  est biffée et réécrite, et la transcription ne la lit pas. Ce qu'il vise se
  déduit de l'usage : c'est une borne uniforme sur le nombre de composantes
  connexes de $X \times_{Y} Y'$, c'est-à-dire une condition de \emph{degré
  borné} qui interdit à un morphisme entre objets connexes de se scinder
  indéfiniment. La page 211 revient dessus, et pour en dire qu'elle le soupçonne
  redondant.}
\end{itemize}

L'axiome \textbf{G3} est le seul de la liste qui ait traversé tout le carton :
c'est la propriété que la page 24, en 1967, obtenait comme un \emph{théorème}
sur $\mathrm{Rev}^{\mathcal{R}}(S)$ --- « tout morphisme de la catégorie se
factorise en le produit d'un épimorphisme et d'un monomorphisme \dots\ tout
monomorphisme est un isomorphisme avec un sommande direct » --- et c'est
l'axiome G3 de SGA 1, exposé V.

\pagerange{209}{211}
\section*{Les démonstrations, et un axiome soupçonné (pages 209 à 211)}

\subsection*{Démonstration de A : le revêtement des repères, une seconde fois}

Les implications (iii) $\Rightarrow$ (ii bis) $\Rightarrow$ (ii) $\Rightarrow$
(i) sont triviales. Pour (i) $\Rightarrow$ (iii) :

Soit $X \in \mathcal{T}_{\mathrm{cf}}$. L'objet final de $\mathcal{T}$ est somme
d'objets $e_{i}$ ($i \in \mathbb{N}$) au-dessus desquels $X|e_{i}$ est de rang
$i$. Posons alors
\[
  P_{i} \;=\; \mathbf{Isom}_{e_{i}}\bigl([1,i]_{e_{i}},\ X\bigr) ,
\]
le faisceau des identifications de $X$ avec l'ensemble constant à $i$ éléments.
C'est un objet de $\mathcal{T}_{\mathrm{cf}}$, l'application
$P_{i} \to e_{i}$ est un épimorphisme --- autrement dit $P_{i}$ est un
\emph{torseur sous $(\mathfrak{S}_{i})_{e_{i}}$} --- et
\[
  X_{P_{i}} \;\simeq\; [1, i]_{P_{i}} .
\]
Les $P_{i} \to e$ forment donc une famille de descente effective universelle
pour $\mathcal{T}$, donc aussi pour $\mathcal{T}_{\mathrm{cf}}$, qui est stable
par limites projectives et inductives finies.

\begin{quote}
\textbf{C'est la démonstration de Serre, page 152 du carton.} Là, $X \to Y$
était un morphisme fini étale de rang $n$, et $P$ le complémentaire de la
diagonale dans $X^{n}$ --- c'est-à-dire le schéma des $n$-uplets de points deux
à deux distincts, ou encore le schéma des identifications de la fibre avec
$I_{n}$ ---, torseur sous $\mathfrak{S}_{n}$. Ici, $\mathbf{Isom}([1,i], X)$
est le même objet, écrit dans un topos quelconque. Le carton contient donc la
construction dans son cas géométrique et dans sa forme abstraite, à
soixante pages de distance, et l'attribution à Serre est portée sur la
première.
\end{quote}

\subsection*{Démonstration de B, et l'axiome G5}

Les conditions A(ii) impliquent aisément G0, G1, G2, G3. Pour G4, il faut
prouver que pour tout $X \in C$ les $S_{i} \to e$ tels que $X_{S_{i}}$ soit
constant fini recouvrent $e$ pour la topologie de la descente effective ;
cela s'obtient en utilisant G2 et G4, en sommant sur les composantes connexes
$e_{i}$ de l'objet final.

Puis l'énoncé de récurrence : si $X \to Y$ est tel que $Y$ soit non vide et $X$
connexe, alors il existe $Y' \to Y$ de descente effective universelle avec
$X_{Y'} \simeq [1,n]_{Y'}$ ; on procède par récurrence sur le rang, et c'est
là que G5 sert.

\textbf{Remarque.} Grothendieck note qu'il soupçonne \textbf{G5} d'être
\emph{redondant}, c'est-à-dire conséquence des autres axiomes, et indique ce
qu'il faudrait pour le montrer : pour toute composante connexe $e_{i}$ de
l'objet final, l'existence d'un foncteur exact
\[
  F : C \longrightarrow (\text{Ens.\ finis}) \qquad \text{tel que}
  \qquad F(e_{i}) \neq \varnothing .
\]
Il ne le démontre pas.\footnote{C'est le second axiome que le carton
soupçonne de trop, après celui de la page 157 --- « $F(X) = \varnothing
\Rightarrow X = 0$ » --- qui, lui, était biffé parce que \emph{démontré}
redondant, conséquence de Gal 5. Ici le soupçon reste un soupçon. Une bonne
partie de la page 211 est barrée de longues diagonales et ne se lit que par
fragments ; ce qui précède est ce que les phrases portantes autorisent.}

\pagerange{212}{214}
\section*{Compter les feuillets sans compter de points : les morphismes de degré $n$ (pages 212 et 214)}

Deux feuillets insérés, que Grothendieck numérote « 1 bis » et « 1 ter » dans
un cercle, portent les quatre conditions que la page 208 annonçait sous A(iv).
Elles sont numérotées à rebours --- (3), (2), (1), (0) --- chaque chiffre
cerclé.

Ce qu'elles font mérite d'être dit avant d'être lu : elles axiomatisent
directement la notion de \emph{morphisme de degré $n$}. Dire qu'un revêtement a
$n$ feuillets suppose d'ordinaire qu'on compte les points d'une fibre ; ici, il
n'y a pas de points, et il faut donc dire ce qu'un degré est par ses
propriétés seules.

\begin{itemize}
\item[\textbf{(3)}] Pour tout objet non vide, on se donne, pour chaque $n$, une
  classe de « morphismes de degré $n$ », telle que :
  \begin{enumerate}
    \item[a)] elle soit stable par changement de base ;
    \item[b)] elle soit \emph{additive} pour les sommes finies ;
    \item[c)] un isomorphisme soit de degré 1 ;
    \item[d')] pour tout $f : X \to Y$, $Y$ soit somme directe disjointe
      effective universelle $Y = \coprod_{n \in \mathbb{N}} Y_{n}$ avec
      $X_{n} = X \times_{Y} Y_{n} \to Y_{n}$ de degré $n$ ;
    \item[d)] si $f$ est de degré $n \geqslant 1$, $f$ soit un morphisme de
      descente \emph{effective} --- donc universelle, grâce à a) ;
    \item[e)] si $f$ est de degré $0$, alors $X = \varnothing$, et $f$ soit un
      isomorphisme sur un \emph{sommande direct universel}.
  \end{enumerate}
\item[\textbf{(2)}] Tous les monomorphismes sont d'un type prescrit.
\item[\textbf{(1)}] Existence des limites projectives \emph{finies}.
\item[\textbf{(0)}] $C$ est une $\mathcal{U}$-catégorie et ses classes
  d'isomorphie forment un ensemble $\mathcal{U}$-petit.
\end{itemize}

La condition d') est celle qui fait tout le travail : elle dit que le degré
n'est pas un nombre attaché au morphisme, mais une \emph{décomposition de la
base}. C'est la traduction exacte, sans points, du fait qu'un morphisme fini
étale a un rang localement constant.

\subsection*{La récurrence}

De ces conditions on déduit A(iii) : pour $X$ donné, les $S \to e$ tels que
$X_{S}$ soit constant fini forment une famille de descente effective
universelle. On se ramène, par a) et d'), au cas d'un $X \to Y$ de degré $n$.
Si $n = 0$ on conclut par 3 e), $X$ étant vide et $Y \to Y$ convenant. Si
$n > 0$, on procède par récurrence, l'énoncé étant supposé acquis pour les
entiers $< n$ :

Utilisant 3 d), posons $X' = X \times_{Y} X$ au-dessus de $Y' = X$ :

\begin{tikzcd}[column sep=small, row sep=small]
X \arrow[d] & X \times_Y X = X' \arrow[l] \arrow[d] \\
Y & X = Y' \arrow[l, no head]
\end{tikzcd}

La \emph{diagonale} $\Delta_{X,Y}$ existe par (1), et par (2) c'est un
isomorphisme sur un sommande direct universel : donc
$X' \simeq Y' \sqcup X'_{1}$. Par 3 a), $X' \to Y'$ est de degré $n$ ; par 3 c),
$Y' \to Y'$ est de degré 1 ; donc par additivité, $X'_{1} \to Y'$ est de degré
$n - 1$. L'hypothèse de récurrence fournit alors une famille
$Y''_{s} \to Y'$ de descente effective universelle qui trivialise $X'_{1}$ ;
et comme $X' \simeq Y' \sqcup X'_{1}$ \emph{universellement},
\[
  X'' \;=\; X \times_{Y} Y'' \;\simeq\; Y'' \sqcup
  \underbrace{\bigl(X'_{1} \times_{Y'} Y''\bigr)}_{\textstyle
  Y'' \times [1, n-1]} \;=\; Y'' \times [1, n] .
\]
Et $Y'' \to Y$ est de descente effective universelle comme composé de deux
tels.

\subsection*{L'unicité}

\textbf{Conclusion.} Dans les conditions de A(iv), la notion de « morphisme de
degré $n$ » est \emph{uniquement déterminée}, et elle équivaut à : il existe
$Y' \to Y$ épimorphisme effectif universel, et de descente effective, avec
\[
  X \times_{Y} Y' \;\simeq\; [1, n]_{Y'} .
\]

C'est-à-dire : \emph{un morphisme est de degré $n$ s'il devient l'ensemble
constant à $n$ éléments après un changement de base convenable}, et cette
définition-là est forcée. On ne pouvait pas la prendre pour définition dès le
départ --- il aurait fallu savoir que « $n$ éléments » a un sens --- et c'est
pour cela que les six propriétés a) à e) sont posées d'abord.

Une \emph{explicitation} accompagne : pour vérifier qu'un $f$ est un morphisme
de descente effective, il suffit de savoir que c'est un épimorphisme effectif
universel, et que les relations d'équivalence sont effectives universelles.

\begin{tikzcd}[column sep=small, row sep=small]
Z'' \arrow[d]
  & Z \arrow[l, dashed] \arrow[d]
  & D \arrow[l, bend left=12] \arrow[l, bend right=12] \arrow[d] \\
Y & X \arrow[l]
  & X \times_Y X \arrow[l, bend left=12] \arrow[l, bend right=12]
\end{tikzcd}

\pagerange{213}{213}
\section*{Un feuillet d'un autre sujet : puissances fibrées et partitions (page 213)}

Une feuille en travers, sans numéro de sa main, barrée de deux longues
diagonales, et sans continuité avec ce qui l'entoure. Elle est transcrite à son
rang et lue ici de même.

Soit $f : X \to Y$ et $E$ un ensemble fini ; soit $X^{E}$ la puissance fibrée
de $X$ au-dessus de $Y$ indexée par $E$. À toute partition $P$ de $E$ est
associée une \emph{immersion fermée}

\begin{tikzcd}[column sep=large]
X^P \arrow[r, hook, "\alpha_P"] & X^E
\end{tikzcd}

--- celle qui identifie les coordonnées à l'intérieur de chaque bloc --- et
l'on a la décomposition disjointe
\[
  X^{E} \;=\; \bigcup_{P} \alpha_{P}\bigl((X^{P})^{\ast}\bigr) ,
\]
où $(X^{P})^{\ast}$ est la partie ouverte des points à coordonnées deux à deux
distinctes. De même, pour la petite diagonale de $Y$,
\[
  (f^{E})^{-1}\bigl(\Delta_{y}^{(E)}\bigr) \;=\;
  \bigcup_{P} \alpha_{P}(F_{P}) \qquad \text{(réunion disjointe)} ,
\]
les morceaux étant de dimension
\[
  n - (\operatorname{card} P - 1)\, d \;\geqslant\;
  n - (\operatorname{card} E - 1)\, d ,
\]
avec égalité pour la seule partition maximale
$P_{0}$.\footnote{La page écrit la minoration sous la forme
$\geqslant n - \operatorname{card}(E)\,d$, ce qui est vrai mais non optimal et
incompatible avec le cas d'égalité qu'elle énonce dans la même phrase : pour
$P = P_{0}$, la partition en $\operatorname{card} E$ blocs, la formule donne
$n - (\operatorname{card} E - 1)d$. C'est cette borne qui est écrite ici. Les
lettres $n$ et $d$ ne sont pas définies sur la page ; le cas travaillé
ci-dessous les fixe : $n = \dim X$ et $d$ la dimension relative de $f$.}

\subsection*{Le cas $\operatorname{card} E = 3$}

Les partitions de $E$ à trois éléments sont : celle en un bloc (la petite
diagonale), les trois en deux blocs, et celle en trois blocs. D'où
\[
  (f^{3})^{-1}\bigl(\Delta_{y}^{(3)}\bigr) \;=\;
  \underset{n - 2d}{F_{3}} \;\cup\;
  \underset{n - d}{\bigcup_{1 \leqslant i < j \leqslant 3} \alpha_{ij}(F_{2})}
  \;\cup\; \underset{n}{\Delta_{X}^{(3)}} ,
\]
les dimensions notées étant les dimensions « génériques ». Le cas numérique
travaillé sur la page est $\dim X = 3$, $\dim Y = 4$, avec
$\dim F_{3} = 1$ et $\dim \alpha_{ij}(F_{2}) = 2$ --- ce qui donne $n = 3$ et
$d = 1$.\footnote{La page numérote deux de ses trois lignes de partitions
« en 3 », là où l'une des deux doit porter « en 1 » : c'est la partition en un
seul bloc, qui donne la petite diagonale. Le chiffre est écrit deux fois de la
même main ; la transcription le conserve avec un \emph{sic}, et cette lecture
donne les trois lignes dans leur ordre correct.}

\textbf{N.B.} Si $f$ est convenable, $\overline{F}_{3} \cap
\Delta_{X}^{(3)} = \varnothing$ et
$\overline{\alpha_{ij}(F_{2})} \cap \Delta_{X}^{(3)} = \varnothing$. Mais la
question que la page pose et laisse ouverte est celle des intersections
croisées :
\[
  \overline{F}_{3} \cap \overline{\alpha_{ij}(F_{2})} ,
  \qquad
  \overline{F}_{3} \cap \alpha_{ij}(\overline{F}_{2}) \;?
\]

\pagerange{215}{216}
\section*{Le cas galoisien, et la dernière proposition (pages 215 et 216)}

\subsection*{C) Le cas galoisien}

Conditions équivalentes, définissant les catégories \emph{galoisiennes} :

\begin{enumerate}
  \item[(i)] Comme A(i), avec l'objet final de $\mathcal{T}$ \emph{connexe}.
  \item[] \textbf{(ii), (ii bis), (iii), (iv)} Comme dans A, en supposant de
    plus l'objet final de $C$ connexe.
  \item[(v)] Les axiomes \textbf{G0}, \textbf{G1}, \textbf{G2}, \textbf{G3},
    et de plus : il existe un foncteur
    \[
      F : C \longrightarrow (\text{Ens.\ fini})
    \]
    \emph{conservatif} et \emph{exact} --- respectant les limites projectives
    finies et commutant aux sommes finies.
  \item[(vi)] Il existe un \emph{groupe profini} $\pi$ et une équivalence
    $C \simeq C(\pi)$.
\end{enumerate}

C'est la liste de SGA 1, exposé V, dans sa forme définitive : les axiomes de la
catégorie, l'existence d'un foncteur fibre exact et conservatif, et la
conclusion --- la catégorie est celle des ensembles finis à opérateurs d'un
groupe profini. La différence avec B tient en un mot : ici l'objet final est
connexe, et le groupoïde se réduit à un groupe.

Il faut mesurer ce que le carton vient de faire. La liste Gal 1 à Gal 6 de la
page 157 exigeait la décomposition en composantes connexes comme un axiome ;
la liste G0 à G5 de la page 209 la déduit ; et la condition (v) ci-dessus
remplace toute la machinerie par une seule hypothèse d'existence, celle du
foncteur fibre. Trois états de la même axiomatique, dans un seul carton, et
c'est le dernier qui est passé dans SGA 1.

\subsection*{La proposition finale}

\textbf{Proposition.} Soit $u : C \to C'$ un foncteur \emph{exact}, $C$ et
$C'$ multigaloisiennes. Alors :

\begin{enumerate}
  \item[(i)] $u$ est \emph{fidèle} $\Longleftrightarrow$ $u$ est
    \emph{conservatif} $\Longleftrightarrow$ pour tout objet $X$ de $C$,
    $u(X) = \varnothing$ entraîne $X = \varnothing$.
  \item[(ii)] $u$ est \emph{pleinement fidèle} $\Longleftrightarrow$ pour tout
    objet $X$ de $C$, $u$ induit une \emph{bijection} de l'ensemble des
    sous-objets de $X$ sur l'ensemble des sous-objets de $u(X)$.
  \item[(iii)] $u$ est une \emph{équivalence} $\Longleftrightarrow$ $u$ est
    pleinement fidèle et tout objet de $C'$ est \emph{subordonné} à un objet de
    l'image.
\end{enumerate}

C'est sur cette proposition, dont la démonstration ne suit pas, que la cote
s'arrête ; la moitié inférieure de la page est nue.

\begin{quote}
\textbf{Ce que dit cette dernière page, et pourquoi elle est à sa place.} Ce
sont les critères sous lesquels les axiomes du domaine de ramification ont été
choisis, page 190, cent trente pages plus tôt dans le même carton. L'axiome
\textbf{R2} --- « l'application $X' \mapsto r(X')$ est une bijection de
l'ensemble des sous-objets directs de $X$ sur l'ensemble des parties ouvertes
et fermées de $r(X)$ » --- est le critère (ii). La proposition 2 de la page
192 --- « $X = \varnothing$ si et seulement si $r(X) = \varnothing$ » --- est
le critère (i). Et le critère (iii), la subordination, est le mot même sur
lequel toute la théorie des familles génératrices est bâtie aux pages 201 à
206.

La réalisation géométrique $r$ n'est pas un foncteur entre catégories
multigaloisiennes --- son but est la catégorie des schémas --- de sorte que la
proposition ne s'y applique pas littéralement. Mais c'est bien elle qui donne
aux axiomes R1 et R2 leur raison d'être : ils demandent de $r$ exactement les
propriétés qui, entre catégories multigaloisiennes, caractérisent la fidélité
et la pleine fidélité. Le carton s'achève donc sur les raisons de son propre
point de départ.
\end{quote}

\end{document}
